% Leçon 21 — Expression orale : vie culturelle au Cameroun et en Chine — Chinois 1ère A — Cours signé OnBuch+
% Programme officiel MINESEC. Compiler avec : tectonic ZHA21-oral-vie-culturelle-au-cameroun-et-en-chine.tex
\newcommand{\DOCMATIERE}{Chinois}
\newcommand{\DOCNIVEAU}{1ère A}
\newcommand{\DOCMODULE}{Module 3 : Citoyenneté et ouverture au monde}
\newcommand{\DOCLECON}{ZHA21}
\newcommand{\DOCTITRE}{Leçon 21 — Expression orale : vie culturelle au Cameroun et en Chine}
% OnBuch+ — préambule « cours signés » ENRICHI (compilé avec tectonic / XeLaTeX)
% Système visuel « Pop » d'OnBuch+ : fond crème (jamais blanc), bordures encre
% épaisses, ombres « dures » décalées, titres Archivo Black en capitales.
% Version MATHÉMATIQUES : graphes (pgfplots/tikz), boîtes pédagogiques
% (définition, propriété, méthode, attention, le savais-tu, exercice résolu,
% exercice, TP, bilan), page de garde et signature OnBuch+ ancrée en bas.
%
% Macros attendues dans main.tex AVANT \input{preamble} :
%   \DOCMATIERE  \DOCNIVEAU  \DOCMODULE  \DOCLECON (numéro)  \DOCTITRE
\documentclass[11pt]{article}

\usepackage{fontspec}
\usepackage[french]{babel} % français = langue principale ; le chinois via \zh{..} / textebox (xeCJK)
\usepackage{xeCJK}
\usepackage{pifont} % \ding{51} \ding{55} utilisés par les modèles
\usepackage[a4paper,margin=2cm,top=3.4cm,bottom=2.8cm,headheight=1.4cm,footskip=1.3cm]{geometry}
\usepackage{amsmath,amssymb,amsfonts}
\usepackage{mathtools} % \xmapsto, \coloneqq... souvent utilisés par les modèles
\usepackage{cancel} % simplifications barrées (bilans de forces)
% \abs{z} (module, valeur absolue) et \norm{u} (norme), utilisés par les modèles
\providecommand{\abs}{}\let\abs\relax\DeclarePairedDelimiter{\abs}{\lvert}{\rvert}
\providecommand{\norm}{}\let\norm\relax\DeclarePairedDelimiter{\norm}{\lVert}{\rVert}
\usepackage[table]{xcolor} % [table] : fournit aussi \rowcolors (alternance de lignes)
\usepackage{pagecolor}
\usepackage{tikz}
\usetikzlibrary{arrows.meta,positioning,calc,shapes.geometric,shapes.misc,shapes.arrows,decorations.pathmorphing,decorations.pathreplacing,decorations.markings,patterns,shadows,mindmap,trees,fit,backgrounds,matrix,intersections,angles,quotes}
\usepackage{pgfplots}
\usepackage[version=4]{mhchem} % équations nucléaires \ce{^{235}_{92}U}
\usepackage[european,siunitx]{circuitikz} % schémas électriques
\pgfplotsset{compat=1.18}
\usepgfplotslibrary{fillbetween,groupplots} % aires entre courbes (calcul intégral)
\usepackage{enumitem}
\usepackage{fancyhdr}
% [most] charge toutes les bibliothèques tcolorbox (listings, minted, raster,
% documentation...) et gonfle inutilement la mémoire TeX sur un document long
% (~300 boîtes) : on ne charge que ce qui est réellement utilisé.
\usepackage[skins,breakable]{tcolorbox}
\usepackage{graphicx}
\usepackage{colortbl}
\usepackage{booktabs}
\usepackage{tabularx}
\usepackage{diagbox} % case d'en-tête diagonale des tableaux à double entrée
% Colonnes extensibles centrée / gauche / droite (C, L, R), souvent utilisées par les modèles
\newcolumntype{C}{>{\centering\arraybackslash}X}
\newcolumntype{L}{>{\raggedright\arraybackslash}X}
\newcolumntype{R}{>{\raggedleft\arraybackslash}X}
\usepackage{array}
\usepackage{multicol}
\usepackage{multirow}
\usepackage{siunitx}
% parse-numbers=false : accepte aussi \SI{2,5 \times 10^{-3}}{...}
\sisetup{locale=FR,output-decimal-marker={,},group-minimum-digits=4,parse-numbers=false}
\DeclareSIUnit\ohm{\ensuremath{\symup{\Omega}}} % U+2126 absent de la police
\DeclareSIUnit\division{div} % réglages d'oscilloscope (ms/div, V/div)
\DeclareSIUnit\tr{tr} % tours (vitesse de rotation, ex. \SI{1500}{\tr\per\minute})
\DeclareSIUnit\molar{M} % concentration molaire (ex. \SI{3}{\molar}, \SI{1}{\micro\molar})
\DeclareSIUnit\an{an} % année (le modèle confond parfois avec \year, un compteur TeX) — ex. \SI{5}{\kilo\meter\per\an}
\DeclareSIUnit\FCFA{FCFA} % franc CFA (ex. \SI{500}{\FCFA\per\kilo\gram})
\DeclareSIUnit\degreeBrix{\textdegree Brix} % teneur en sucre d'un sirop/jus (réfractométrie)
\DeclareSIUnit\month{mois} % le modèle utilise parfois \month, un compteur TeX, comme unité de durée
\usepackage{qrcode}
% \degree : commande fréquemment utilisée par les modèles pour un angle ou une
% température en dehors de \SI{}{} (non fournie par les paquets chargés).
\providecommand{\degree}{^{\circ}}
% \qed (fin de démonstration), utilisé par les modèles sans amsthm
\providecommand{\qed}{\hfill$\square$}
% \barre : double barre des valeurs interdites dans un tableau de variations (tkz-tab)
\providecommand{\barre}{\ensuremath{\|}}
% environnement proof (amsthm non chargé) utilisé par les modèles
\ifdefined\proof\else\newenvironment{proof}[1][Démonstration]{\par\noindent\textit{#1.}\ }{\qed\par}\fi
\usepackage{lastpage}
\usepackage{hyperref}
\hypersetup{hidelinks}

% ============================================================
% Palette « Pop » OnBuch+ — mêmes hex que la classe P (plus_ui.dart)
% ============================================================
\definecolor{poporange}{HTML}{F59321}
\definecolor{popdark}{HTML}{E07A0C}
\definecolor{popcream}{HTML}{FFF6E8}
\definecolor{popink}{HTML}{1C1714}
\definecolor{poppurple}{HTML}{7A5AE0}
\definecolor{popgreen}{HTML}{2E9E6A}
\definecolor{poppink}{HTML}{E0457B}
\definecolor{popblue}{HTML}{2F7DE1}
\definecolor{popgold}{HTML}{C98A12}
\definecolor{popmuted}{HTML}{8A7F76}
\definecolor{popline}{HTML}{EADFD2}
\definecolor{popgray}{HTML}{8A7F76}
\colorlet{popgrey}{popgray}
% Alias tolérants (le modèle nomme parfois les couleurs différemment)
\colorlet{popwhite}{white}
\colorlet{popblack}{popink}
\colorlet{popred}{poppink}
\colorlet{popyellow}{popgold}
% Teintes claires (fonds de tableaux/encadrés) — le modèle nomme parfois
% les couleurs avec un suffixe L (ex. popblueL) sans les avoir définies.
\colorlet{poporangeL}{poporange!20}
\colorlet{popdarkL}{popdark!20}
\colorlet{poppurpleL}{poppurple!20}
\colorlet{popgreenL}{popgreen!20}
\colorlet{poppinkL}{poppink!20}
\colorlet{popblueL}{popblue!20}
\colorlet{popgoldL}{popgold!20}
\colorlet{popredL}{popred!20}
\colorlet{popgrayL}{popgray!20}
% Teintes claires pour fonds de tableaux / zones
\colorlet{popblueL}{popblue!10!white}
\colorlet{poporangeL}{poporange!14!white}
\colorlet{popgreenL}{popgreen!12!white}
\colorlet{poppurpleL}{poppurple!12!white}
\colorlet{poppinkL}{poppink!12!white}
\colorlet{popgoldL}{popgold!14!white}

\pagecolor{popcream}

% ============================================================
% Typographie
% ============================================================
\defaultfontfeatures{Ligatures=TeX}
\setmainfont{PlusJakartaSans-Medium.ttf}[Path=../../../fonts/,BoldFont=PlusJakartaSans-Bold.ttf]
% Symboles, API et emojis absents de la police principale : repli sur DejaVu Sans (caractères actifs)
\newfontface\symfont{DejaVuSans.ttf}[Path=../../../fonts/]
\makeatletter
\newcommand\symactive[1]{\catcode#1=13 \begingroup\lccode`\~=#1 \lowercase{\endgroup\def~}{{\symfont\char#1}}}
\newcommand\symblank[1]{\catcode#1=13 \begingroup\lccode`\~=#1 \lowercase{\endgroup\def~}{}}
\newcommand\symhyph[1]{\catcode#1=13 \begingroup\lccode`\~=#1 \lowercase{\endgroup\def~}{\mbox{-}}}
\newcount\symc
\newcommand\symrange[2]{\symc=#1 \@whilenum\symc<#2 \do{\expandafter\symactive\expandafter{\the\symc}\advance\symc by 1 }}
\newcommand\symblankrange[2]{\symc=#1 \@whilenum\symc<#2 \do{\expandafter\symblank\expandafter{\the\symc}\advance\symc by 1 }}
\makeatother
\symrange{"0250}{"02FF}\symactive{"2713}\symactive{"2714}\symactive{"2717}\symactive{"2718}\symactive{"2610}\symactive{"2611}\symactive{"2612}
\symrange{"2600}{"27BF}
\catcode"2705=13 \begingroup\lccode`\~="2705 \lowercase{\endgroup\def~}{{\symfont\char"2713}}\symblank{"26BD}\symblank{"26A1}
\symrange{"2460}{"257F}\symactive{"223C}\symactive{"2248}\symactive{"2260}
\symactive{"01F8}\symactive{"01F9}\symactive{"1E3E}\symactive{"1E3F}
\symhyph{"2011}\symhyph{"2E1A}\symblankrange{"1F300}{"1FAFF}\symblank{"FE0F}
% Caractères chinois : WenQuanYi Zen Hei (sous-ensemble GB2312 + pinyin), gras/italique simulés
\setCJKmainfont{WenQuanYiZenHei-subset.ttf}[Path=../../../fonts/,AutoFakeBold=3,AutoFakeSlant=0.2]
\xeCJKsetup{CJKspace=false}
% Pinyin : voyelles à caron (ǎ ǐ ǒ ǔ ǖ ǘ ǚ ǜ) absentes de la police principale -> caractères actifs
\newfontface\pyfont{WenQuanYiZenHei-subset.ttf}[Path=../../../fonts/]
\newcommand\pyactive[1]{\catcode#1=13 \begingroup\lccode`\~=#1 \lowercase{\endgroup\def~}{{\pyfont\char#1}}}
\pyactive{"01CD}\pyactive{"01CE}\pyactive{"01CF}\pyactive{"01D0}\pyactive{"01D1}\pyactive{"01D2}\pyactive{"01D3}\pyactive{"01D4}\pyactive{"01D5}\pyactive{"01D6}\pyactive{"01D7}\pyactive{"01D8}\pyactive{"01D9}\pyactive{"01DA}\pyactive{"01DB}\pyactive{"01DC}
% Maths en sans-serif (Fira Math) avec chiffres et lettres droites de Plus
% Jakarta, pour que les formules se fondent dans le texte.
\usepackage[math-style=ISO,bold-style=ISO]{unicode-math}
\setmathfont{FiraMath-Regular.otf}[Path=../../../fonts/]
\setmathfont{PlusJakartaSans-Medium.ttf}[Path=../../../fonts/,range={up/num,up/{Latin,latin},\mathup}]
\usepackage{icomma} % virgule décimale sans espace en mode math
% Caractères mathématiques Unicode absents des polices de texte (Plus Jakarta,
% Archivo Black) : rendus par leur équivalent mathématique, en texte comme en
% formule — sinon ils s'affichent en carré vide (ex. « Arithmétique dans ℤ »).
\usepackage{newunicodechar}
\newunicodechar{ℝ}{\ensuremath{\mathbb{R}}}
\newunicodechar{ℤ}{\ensuremath{\mathbb{Z}}}
\newunicodechar{ℂ}{\ensuremath{\mathbb{C}}}
\newunicodechar{ℕ}{\ensuremath{\mathbb{N}}}
\newunicodechar{ℚ}{\ensuremath{\mathbb{Q}}}
\newunicodechar{𝔻}{\ensuremath{\mathbb{D}}}
\newunicodechar{α}{\ensuremath{\alpha}}
\newunicodechar{β}{\ensuremath{\beta}}
\newunicodechar{γ}{\ensuremath{\gamma}}
\newunicodechar{δ}{\ensuremath{\delta}}
\newunicodechar{ε}{\ensuremath{\varepsilon}}
\newunicodechar{θ}{\ensuremath{\theta}}
\newunicodechar{λ}{\ensuremath{\lambda}}
\newunicodechar{μ}{\ensuremath{\mu}}
\newunicodechar{σ}{\ensuremath{\sigma}}
\newunicodechar{τ}{\ensuremath{\tau}}
\newunicodechar{φ}{\ensuremath{\varphi}}
\newunicodechar{ψ}{\ensuremath{\psi}}
\newunicodechar{ω}{\ensuremath{\omega}}
\newunicodechar{Σ}{\ensuremath{\Sigma}}
% majuscules produites par \MakeUppercase dans les titres (α-aminés -> Α)
\newunicodechar{Α}{\ensuremath{\alpha}}
\newunicodechar{Β}{\ensuremath{\beta}}
\newunicodechar{Γ}{\ensuremath{\Gamma}}
\newunicodechar{Δ}{\ensuremath{\Delta}}
\newunicodechar{Ε}{\ensuremath{\varepsilon}}
\newunicodechar{Θ}{\ensuremath{\Theta}}
\newunicodechar{Λ}{\ensuremath{\Lambda}}
\newunicodechar{Μ}{\ensuremath{\mu}}
\newunicodechar{Τ}{\ensuremath{\tau}}
\newunicodechar{Φ}{\ensuremath{\Phi}}
\newunicodechar{Ψ}{\ensuremath{\Psi}}
\newunicodechar{Ω}{\ensuremath{\Omega}}
\newunicodechar{↦}{\ensuremath{\mapsto}}
\newunicodechar{∈}{\ensuremath{\in}}
\newunicodechar{∉}{\ensuremath{\notin}}
\newunicodechar{∧}{\ensuremath{\wedge}}
\newunicodechar{⇔}{\ensuremath{\Leftrightarrow}}
\newunicodechar{⟺}{\ensuremath{\Longleftrightarrow}}
\newunicodechar{⇒}{\ensuremath{\Rightarrow}}
\newunicodechar{⟹}{\ensuremath{\Longrightarrow}}
\newunicodechar{∘}{\ensuremath{\circ}}
\newunicodechar{∪}{\ensuremath{\cup}}
\newunicodechar{∩}{\ensuremath{\cap}}
\newunicodechar{≡}{\ensuremath{\equiv}}
\newunicodechar{⊂}{\ensuremath{\subset}}
\newunicodechar{⊥}{\ensuremath{\perp}}
\newunicodechar{∃}{\ensuremath{\exists}}
\newunicodechar{∀}{\ensuremath{\forall}}
\newunicodechar{∛}{\ensuremath{\sqrt[3]{\,}}}
\newunicodechar{∜}{\ensuremath{\sqrt[4]{\,}}}
\newunicodechar{⃗}{\ensuremath{{}^{\to}}}
\newunicodechar{ⁿ}{\ifmmode^{n}\else\textsuperscript{\ensuremath{n}}\fi}
\newunicodechar{ˣ}{\ifmmode^{x}\else\textsuperscript{\ensuremath{x}}\fi}
\newunicodechar{ᵇ}{\ifmmode^{b}\else\textsuperscript{\ensuremath{b}}\fi}
\newunicodechar{ʳ}{\ifmmode^{r}\else\textsuperscript{\ensuremath{r}}\fi}
\newunicodechar{ᵘ}{\ifmmode^{u}\else\textsuperscript{\ensuremath{u}}\fi}
\newunicodechar{ʸ}{\ifmmode^{y}\else\textsuperscript{\ensuremath{y}}\fi}
\newunicodechar{ᵃ}{\ifmmode^{a}\else\textsuperscript{\ensuremath{a}}\fi}
\newunicodechar{ᵉ}{\ifmmode^{e}\else\textsuperscript{\ensuremath{e}}\fi}
\newunicodechar{ᵏ}{\ifmmode^{k}\else\textsuperscript{\ensuremath{k}}\fi}
\newunicodechar{ᵖ}{\ifmmode^{p}\else\textsuperscript{\ensuremath{p}}\fi}
\newunicodechar{ᴺ}{\ifmmode^{N}\else\textsuperscript{\ensuremath{N}}\fi}
\newunicodechar{ᶜ}{\ifmmode^{c}\else\textsuperscript{\ensuremath{c}}\fi}
\newunicodechar{ʰ}{\ifmmode^{h}\else\textsuperscript{\ensuremath{h}}\fi}
\newunicodechar{ᵠ}{\ifmmode^{q}\else\textsuperscript{\ensuremath{q}}\fi}
\newunicodechar{ₙ}{\ifmmode_{n}\else\textsubscript{\ensuremath{n}}\fi}
\newunicodechar{ₐ}{\ifmmode_{a}\else\textsubscript{\ensuremath{a}}\fi}
\newunicodechar{ᵢ}{\ifmmode_{i}\else\textsubscript{\ensuremath{i}}\fi}
\newunicodechar{ᵣ}{\ifmmode_{r}\else\textsubscript{\ensuremath{r}}\fi}
\newunicodechar{ₖ}{\ifmmode_{k}\else\textsubscript{\ensuremath{k}}\fi}
\newunicodechar{ᵦ}{\ifmmode_{\beta}\else\textsubscript{\ensuremath{\beta}}\fi}
\newunicodechar{ₓ}{\ifmmode_{x}\else\textsubscript{\ensuremath{x}}\fi}
\newfontfamily\popdisplay{ArchivoBlack-Regular.ttf}[Path=../../../fonts/]
\newfontfamily\poplogo{SpaceGrotesk-Bold.ttf}[Path=../../../fonts/]
\newfontfamily\popsemi{PlusJakartaSans-SemiBold.ttf}[Path=../../../fonts/]
\newfontfamily\popxbold{PlusJakartaSans-ExtraBold.ttf}[Path=../../../fonts/]
\newfontfamily\popxboldls{PlusJakartaSans-ExtraBold.ttf}[Path=../../../fonts/,LetterSpace=3.0]

\setlist[enumerate]{leftmargin=1.7em,itemsep=5pt,topsep=4pt}
\setlist[itemize]{leftmargin=1.7em,itemsep=4pt,topsep=4pt,label={\color{poporange}\textbullet}}
\setlength{\parindent}{0pt}
\setlength{\parskip}{5pt}
\renewcommand{\arraystretch}{1.25}

% pgfplots : style Pop par défaut
\pgfplotsset{
  every axis/.append style={axis line style={popink,line width=1pt},tick style={popink},
    grid style={popline,line width=0.5pt},label style={font=\small},tick label style={font=\footnotesize}},
  popcurve/.style={poporange,line width=1.8pt,no markers,smooth},
}
% Même style disponible dans un \draw TikZ simple (hors pgfplots)
\tikzset{popcurve/.style={poporange,line width=1.8pt}}
\tikzset{popgrid/.style={popline,very thin}}
\colorlet{popcurve}{poporange} % le nom du style est parfois utilisé comme couleur

% ============================================================
% Pill / squiggle
% ============================================================
\newcommand{\poppill}[3]{%
  \tikz[baseline=-0.55ex]\node[draw=popink,line width=1.2pt,rounded corners=2.6pt,%
    fill=#2,inner xsep=6.5pt,inner ysep=3pt]{%
    \popxboldls\fontsize{7}{7}\selectfont\color{#3}\MakeUppercase{#1}};%
}
\newcommand{\popsquiggle}[1]{%
  \tikz[baseline=-0.4ex]{
    \draw[#1,line width=2.6pt,line cap=round]
      plot [smooth] coordinates {(0,0.09) (0.34,0.01) (0.68,0.21) (1.02,0.01) (1.36,0.10)};
  }%
}
% Mot-clé surligné (marqueur orange sous le texte)
\newcommand{\cle}[1]{{\popxbold\color{popdark}#1}}

% ============================================================
% En-tête / pied
% ============================================================
\pagestyle{fancy}
\fancyhf{}
\renewcommand{\headrulewidth}{0pt}
\renewcommand{\footrulewidth}{0pt}
\lhead{\raisebox{-0.15ex}{\poplogo\fontsize{13}{13}\selectfont\color{popink}OnBuch\color{poporange}+}}
\rhead{\poppill{\DOCMATIERE\ \textbullet\ \DOCNIVEAU\ \textbullet\ Leçon \DOCLECON}{white}{popink}}
\lfoot{\popxbold\fontsize{7.6}{7.6}\selectfont\color{popmuted}\MakeUppercase{Cours signé OnBuch+}\ \textbullet\ {\popsemi\DOCTITRE}}
\rfoot{\tikz[baseline=-0.6ex]\node[rounded corners=8.5pt,fill=popink,minimum height=17pt,minimum width=17pt,inner xsep=5pt,inner sep=0]{\popxbold\fontsize{8}{8}\selectfont\color{popcream}\thepage\,/\,\pageref*{LastPage}};}

% ============================================================
% Titres
% ============================================================
\newcommand{\chaptitre}[2]{%
  \begin{tcolorbox}[enhanced,colback=poporange,colframe=popink,boxrule=2.6pt,arc=11pt,
      drop shadow=popink,left=15pt,right=15pt,top=12pt,bottom=11pt,before skip=3pt,after skip=18pt]
    \poppill{#2}{popink}{popcream}\par\vspace{9pt}
    {\raggedright\hyphenpenalty=10000\exhyphenpenalty=10000\popdisplay\fontsize{22}{24}\selectfont\color{popcream}\MakeUppercase{#1}\par}
  \end{tcolorbox}
}
% Section numérotée : pastille orange avec le numéro + titre Archivo
\newcounter{popsec}
\newcommand{\coursec}[1]{%
  \stepcounter{popsec}\setcounter{popsub}{0}%
  \par\vspace{16pt}\noindent
  \begin{minipage}{\linewidth}
  \tikz[baseline=-0.7ex]\node[draw=popink,line width=1.6pt,fill=poporange,rounded corners=4pt,minimum size=22pt,inner sep=0]{\popdisplay\fontsize{12}{12}\selectfont\color{popcream}\thepopsec};%
  \hspace{8pt}{\popdisplay\fontsize{15}{17}\selectfont\color{popink}\MakeUppercase{#1}}\par
  \vspace{4pt}\noindent{\color{poporange}\rule{36pt}{3.6pt}}
  \end{minipage}\par\nopagebreak\vspace{8pt}
}
\newcounter{popsub}
\newcommand{\courssub}[1]{%
  \stepcounter{popsub}%
  \par\vspace{9pt}\noindent{\popxbold\fontsize{11.5}{13}\selectfont\color{popink}{\color{poporange}\thepopsec.\thepopsub}\hspace{6pt}#1}\par\nopagebreak\vspace{4pt}
}

% ============================================================
% Boîtes pédagogiques — même recette : fond blanc, bordure épaisse colorée,
% ombre dure encre, badge plein. Titre optionnel [#1].
% ============================================================
\tcbset{popbase/.style={breakable,enhanced,colback=white,boxrule=2.3pt,arc=9pt,
  shadow={3pt}{-3pt}{0pt}{popink},left=14pt,right=14pt,top=11pt,bottom=11pt,before skip=11pt,after skip=15pt,
  fontupper=\popsemi\fontsize{10.6}{14.5}\selectfont\color{popink},
  fontlower=\popsemi\fontsize{10.6}{14.5}\selectfont\color{popink}}}
% Coche dessinée (le glyphe ✓ n'existe pas dans les polices du document)
\newcommand{\popcheck}[1]{\tikz[baseline=-0.5ex]\draw[#1,line width=1.6pt,line cap=round,line join=round](-0.13,0.0)--(-0.04,-0.09)--(0.13,0.1);}
\providecommand{\coche}{\popcheck{popgreen}}
\providecommand{\resultat}[1]{\par\smallskip\noindent\fcolorbox{popgreen}{popgreenL}{\parbox{\dimexpr\linewidth-2\fboxsep-2\fboxrule\relax}{\textbf{Résultat.} #1}}\par\smallskip}
\newcommand{\popboxhead}[3]{\poppill{#1}{#2}{white}\ifx\relax#3\relax\else\hspace{6pt}{\popxbold\fontsize{10.6}{12}\selectfont\color{popink}#3}\fi\par\vspace{7pt}}

\newtcolorbox{aretenir}[1][]{popbase,colframe=popgreen,before upper={\popboxhead{À retenir}{popgreen}{#1}}}
% Texte en chinois (document de lecture, dialogue, extrait) : cadre
% d'étude. Un titre optionnel (source, auteur) s'affiche dans l'en-tête.
\newtcolorbox{textebox}[1][]{popbase,colframe=popblue,before upper={\popboxhead{文本 · Texte}{popblue}{#1}},after upper={}}
% Marques ✓ / ✗ (forme correcte / fautive) souvent utilisées par les modèles
\providecommand{\cross}{\textcolor{poppink}{\ensuremath{\times}}}
\providecommand{\xmark}{\cross}\providecommand{\crossmark}{\cross}\providecommand{\cmark}{\ensuremath{\checkmark}}
% Ligne de réponse / justification à compléter (inventée par les modèles)
\providecommand{\justification}[1]{\par\noindent\textit{Justification~:} \dotfill\par}
% \textipa (package tipa non chargé) : la police ne couvre pas l'API -> simple passage
\providecommand{\textipa}[1]{#1}
% Soulignement (ulem non chargé) : \uline, \ul
\providecommand{\uline}[1]{\underline{#1}}\providecommand{\ul}[1]{\underline{#1}}
% \glossaire{\item ...} inventée par les modèles : liste de glossaire
\providecommand{\glossaire}[1]{\begin{itemize}#1\end{itemize}}
% \alert (beamer) inventée par les modèles : texte mis en évidence
\providecommand{\alert}[1]{\textbf{\textcolor{poppink}{#1}}}
% variantes ASCII d'ordinaux écrites en macro
\providecommand{\iere}{\textsuperscript{re}}\providecommand{\ieme}{\textsuperscript{e}}\providecommand{\ere}{\textsuperscript{re}}\providecommand{\eme}{\textsuperscript{e}}\providecommand{\er}{\textsuperscript{er}}\providecommand{\ie}{\textsuperscript{e}}
% Ordinaux français écrits en macro par les modèles : 1\ière, 3\ième
\newcommand{\ière}{\textsuperscript{re}}\newcommand{\ième}{\textsuperscript{e}}
% \exemple (inventée par les modèles) : étiquette « Exemple : »
\providecommand{\exemple}{\textbf{Exemple~:}\ }
% \square utilisable aussi hors mode math (cases à cocher)
\AtBeginDocument{\let\squareMath\square\renewcommand{\square}{\ensuremath{\squareMath}}}
% Mot ou phrase en chinois à l'intérieur d'un paragraphe français
\newcommand{\zh}[1]{\ifmmode\text{#1}\else{#1}\fi}
\let\eng\zh \let\esp\zh \let\all\zh % alias tolérés
% flèches d'intonation inventées par les modèles
\providecommand{\textsearrow}{}\renewcommand{\textsearrow}{\ensuremath{\searrow}}\providecommand{\textnearrow}{}\renewcommand{\textnearrow}{\ensuremath{\nearrow}}
% \fr{..} : mot français (inventé par les modèles)
\providecommand{\fr}[1]{\textit{#1}}
% zhbox : encadré de texte chinois inventé par les modèles
\newenvironment{zhbox}{\begin{quote}}{\end{quote}}
% \courssubsub : titre de niveau 3 inventé par les modèles
\providecommand{\courssubsub}[1]{\par\medskip\noindent\textbf{#1}\par\smallskip}
% \vs : « versus » inventé par les modèles
\providecommand{\vs}{\textit{vs}\ }
% \blank : case à compléter inventée par les modèles
\providecommand{\blank}[1][]{\underline{\hspace{1.6em}}}
\newtcolorbox{exemplebox}[1][]{popbase,colframe=poporange,before upper={\popboxhead{Exemple}{poporange}{#1}}}
\newtcolorbox{definition}[1][]{popbase,colframe=popblue,before upper={\popboxhead{Définition}{popblue}{#1}}}
\newtcolorbox{propriete}[1][]{popbase,colframe=poppurple,before upper={\popboxhead{Propriété}{poppurple}{#1}}}
\newtcolorbox{methode}[1][]{popbase,colframe=popgold,before upper={\popboxhead{Méthode}{popgold}{#1}}}
\newtcolorbox{attention}[1][]{popbase,colframe=poppink,before upper={\popboxhead{Attention}{poppink}{#1}}}
\newtcolorbox{experience}[1][]{popbase,colframe=popblue,colback=popblueL,before upper={\popboxhead{Expérience}{popblue}{#1}}}
% « Le savais-tu ? » : carte violette pleine (contexte, culture, Cameroun)
\newtcolorbox{savaistu}[1][]{popbase,colback=poppurple,colframe=popink,
  fontupper=\popsemi\fontsize{10.4}{14.2}\selectfont\color{white},
  before upper={\poppill{Le savais-tu\,?}{white}{poppurple}\ifx\relax#1\relax\else\hspace{6pt}{\popxbold\color{white}#1}\fi\par\vspace{7pt}}}
% Exercice résolu : énoncé en haut, \tcblower puis la solution sur fond crème
\newcounter{popexo}\newcounter{popexr}
\newtcolorbox{exoresolu}[1][]{popbase,colframe=poporange,colbacklower=popcream,
  segmentation style={popink,line width=1.2pt,dashed},
  before upper={\stepcounter{popexr}\popboxhead{Exercice résolu \thepopexr}{poporange}{#1}},
  before lower={\poppill{Solution}{popink}{popcream}\par\vspace{6pt}}}
% Exercice (non résolu en place) — la correction est dans la partie Corrigés
\newtcolorbox{exercice}[1][]{popbase,colframe=popink,
  before upper={\stepcounter{popexo}\popboxhead{Exercice \thepopexo}{popink}{#1}}}
% Corrigé
\newtcolorbox{corrige}[1][]{popbase,colframe=popgreen,colback=popgreenL,
  before upper={\popboxhead{Corrigé}{popgreen}{#1}}}
% Objectifs / prérequis / bilan
\newtcolorbox{objectifs}{popbase,colframe=popink,colback=poporangeL,
  before upper={\popboxhead{Objectifs de la leçon}{popink}{}}}
\newtcolorbox{prerequis}{popbase,colframe=popink,colback=popblueL,
  before upper={\popboxhead{Prérequis}{popblue}{}}}
\newtcolorbox{bilan}{popbase,colframe=popink,colback=white,boxrule=2.8pt,
  before upper={\popboxhead{Fiche bilan}{poporange}{L'essentiel à connaître pour l'examen}}}
% Figure encadrée avec légende
\newtcolorbox{popfigure}[1][]{enhanced,breakable=false,colback=white,colframe=popline,boxrule=1.4pt,arc=8pt,
  left=10pt,right=10pt,top=10pt,bottom=8pt,before skip=10pt,after skip=12pt,halign=center,
  lower separated=false,halign lower=center,
  fontlower=\popsemi\fontsize{9}{11.5}\selectfont\color{popmuted},
  #1}
\newcommand{\legende}[1]{\par\vspace{6pt}{\popsemi\fontsize{9}{11.5}\selectfont\color{popmuted}#1}}

% ============================================================
% Signature OnBuch+
% ============================================================
\newcommand{\poplogoinline}[1]{{\poplogo\fontsize{#1}{#1}\selectfont\color{popink}OnBuch\color{poporange}+}}
\newcommand{\signatureonbuch}{%
  \begin{tcolorbox}[enhanced,colback=popink,colframe=popink,boxrule=2.2pt,arc=10pt,
      drop shadow=poporange,left=14pt,right=14pt,top=10pt,bottom=10pt,before skip=0pt,after skip=0pt]
    \tikz[baseline=-0.7ex]\node[circle,fill=popgreen,draw=popcream,line width=1.3pt,minimum size=22pt,inner sep=0]{\popcheck{white}};%
    \hspace{9pt}\begin{minipage}[c]{0.62\linewidth}
      {\popxbold\fontsize{12}{13}\selectfont\color{popcream}Signé\ \textbullet\ {\poplogo OnBuch\color{poporange}+}}\par\vspace{2pt}
      {\raggedright\popsemi\fontsize{8}{10.5}\selectfont\color{popline}\MakeUppercase{Équipe pédagogique OnBuch+}\\\MakeUppercase{Conforme au programme officiel MINESEC}\par}
    \end{minipage}\hfill
    {\poplogo\fontsize{22}{22}\selectfont\color{popcream}OnBuch\color{poporange}+}
  \end{tcolorbox}%
}

% ============================================================
% Page de garde
% #1 titre · #2 matière · #3 niveau · #4 module · #5 n° leçon · #6 durée ·
% #7 description / objectifs (texte ou itemize)
% ============================================================
\newcommand{\pagegarde}[7]{%
  \thispagestyle{empty}
  \newgeometry{a4paper,margin=1.7cm,top=1.6cm,bottom=1.4cm}%
  \begin{tikzpicture}[remember picture,overlay]
    \fill[poporange!18] ([xshift=-3.2cm,yshift=-3.6cm]current page.north east) circle (4.6cm);
    \fill[poppurple!14] ([xshift=2.4cm,yshift=7.6cm]current page.south west) circle (3.8cm);
  \end{tikzpicture}%
  {\poplogo\fontsize{30}{30}\selectfont\color{popink}OnBuch\color{poporange}+}\hfill
  \raisebox{0.6ex}{\poppill{Cours signé \textbullet\ Premium}{popink}{popcream}}\par
  \vspace{1pt}\popsquiggle{poppurple}\par
  \vspace{8pt}
  \begin{tcolorbox}[enhanced,colback=poporange,colframe=popink,boxrule=2.8pt,arc=13pt,
      drop shadow=popink,left=18pt,right=18pt,top=12pt,bottom=13pt,after skip=10pt]
    \poppill{#2 \textbullet\ #3}{popink}{popcream}\hspace{5pt}\poppill{#4}{white}{popink}\par\vspace{8pt}
    {\popdisplay\fontsize{14}{14}\selectfont\color{popink}LEÇON #5}\par\vspace{4pt}
    {\raggedright\hyphenpenalty=10000\exhyphenpenalty=10000\popdisplay\fontsize{25}{27}\selectfont\color{popcream}\MakeUppercase{#1}\par}
    \vspace{8pt}
    {\popxbold\fontsize{9.5}{9.5}\selectfont\color{popink}\MakeUppercase{Durée conseillée : #6}}
  \end{tcolorbox}
  \begin{tcolorbox}[enhanced,colback=white,colframe=popink,boxrule=2.2pt,arc=10pt,
      drop shadow=popink,left=16pt,right=16pt,top=10pt,bottom=10pt,after skip=8pt]
    \poppill{Dans cette leçon}{poppurple}{white}\par\vspace{5pt}
    \popsemi\fontsize{9.3}{12.6}\selectfont\color{popink}#7
  \end{tcolorbox}
  \noindent\begin{minipage}[t]{0.487\linewidth}
    \begin{tcolorbox}[enhanced,colback=popgreen,colframe=popink,boxrule=2.2pt,arc=10pt,
        drop shadow=popink,left=11pt,right=11pt,top=10pt,bottom=10pt,height=3.1cm,valign=center]
      \tcbox[colback=white,colframe=popink,boxrule=1.3pt,arc=5pt,boxsep=0pt,nobeforeafter,
        left=3pt,right=3pt,top=3pt,bottom=3pt,valign=center]{\qrcode[height=1.9cm]{https://play.google.com/store/apps/details?id=com.luvvix.onbuch&hl=fr}}%
      \hspace{8pt}\begin{minipage}[c]{0.5\linewidth}
        {\popdisplay\fontsize{11}{12.5}\selectfont\color{white}\MakeUppercase{Télécharge OnBuch}}\par\vspace{4pt}
        {\raggedright\popsemi\fontsize{8.4}{11}\selectfont\color{white}Gratuit \textbullet\ Google Play\\Tuteur IA, annales, concours, résultats\par}
      \end{minipage}
    \end{tcolorbox}
  \end{minipage}\hfill
  \begin{minipage}[t]{0.487\linewidth}
    \begin{tcolorbox}[enhanced,colback=poppurple,colframe=popink,boxrule=2.2pt,arc=10pt,
        drop shadow=popink,left=11pt,right=11pt,top=10pt,bottom=10pt,height=3.1cm,valign=center]
      \tikz[baseline=-0.7ex]\node[circle,fill=white,draw=popink,line width=1.3pt,minimum size=1.6cm,inner sep=0]{\popdisplay\fontsize{24}{24}\selectfont\color{poppurple}+};%
      \hspace{8pt}\begin{minipage}[c]{0.6\linewidth}
        {\popdisplay\fontsize{11}{12.5}\selectfont\color{white}\MakeUppercase{Passe à OnBuch+}}\par\vspace{4pt}
        {\raggedright\popsemi\fontsize{8.4}{11}\selectfont\color{white}Tous les cours signés, Tuteur IA illimité, fiches PDF — onglet OnBuch+ de l'app\par}
      \end{minipage}
    \end{tcolorbox}
  \end{minipage}
  \vspace{8pt}
  \popcontenu
  \vfill
  \signatureonbuch
  \restoregeometry
  \newpage
}

% Rangée « Ce cours contient » (page de garde)
\newcommand{\poptile}[3]{% #1 couleur #2 gros texte #3 légende
  \begin{tcolorbox}[enhanced,colback=#1,colframe=popink,boxrule=2pt,arc=9pt,shadow={2.5pt}{-2.5pt}{0pt}{popink},
      left=6pt,right=6pt,top=9pt,bottom=9pt,halign=center,valign=center,height=1.75cm,width=\linewidth,nobeforeafter]
    {\popdisplay\fontsize{12.5}{13}\selectfont\color{white}\MakeUppercase{#2}}\par\vspace{3pt}
    {\centering\popsemi\fontsize{7.8}{9.5}\selectfont\color{white}#3\par}
  \end{tcolorbox}}
\newcommand{\popcontenu}{%
  \par\noindent{\popxboldls\fontsize{7.5}{7.5}\selectfont\color{popmuted}CE COURS SIGNÉ CONTIENT}\par\vspace{7pt}
  \noindent\begin{minipage}[t]{0.235\linewidth}\poptile{popblue}{Cours}{complet, illustré, pas à pas}\end{minipage}\hfill
  \begin{minipage}[t]{0.235\linewidth}\poptile{poporange}{Méthodes}{et exercices résolus}\end{minipage}\hfill
  \begin{minipage}[t]{0.235\linewidth}\poptile{poppink}{Exercices}{type examen + corrigés}\end{minipage}\hfill
  \begin{minipage}[t]{0.235\linewidth}\poptile{popgreen}{Fiche bilan}{l'essentiel à retenir}\end{minipage}\par}

% Sommaire
\newtcolorbox{sommaire}{enhanced,colback=white,colframe=popink,boxrule=2.2pt,arc=10pt,
  drop shadow=popink,left=18pt,right=18pt,top=13pt,bottom=13pt,before skip=6pt,after skip=20pt,
  before upper={\poppill{Sommaire}{poppurple}{white}\par\vspace{9pt}\popsemi\fontsize{10.6}{14.5}\selectfont\color{popink}}}

% Signature de fin de cours — poussée tout en bas de la dernière page
\newcommand{\findecours}{%
  \par\vfill\nopagebreak
  \begin{center}\popsquiggle{poporange}\end{center}\vspace{4pt}
  \signatureonbuch
}
\AtBeginDocument{\def\llbracket{\mathopen{[\mkern-3.5mu[}}\def\rrbracket{\mathclose{]\mkern-3.5mu]}}} % intervalles d'entiers (glyphe ⟦ absent de la police)
% Glyphes absents de Fira Math : redéfinis à partir de glyphes disponibles
\AtBeginDocument{%
  \def\setminus{\mathbin{\backslash}}\let\smallsetminus\setminus
  \def\vdots{\mathord{\vbox{\baselineskip3.2pt\lineskiplimit0pt\kern4pt\hbox{.}\hbox{.}\hbox{.}}}}%
  \def\ddots{\mathinner{\mkern1mu\raise6.5pt\hbox{.}\mkern2mu\raise3.6pt\hbox{.}\mkern2mu\raise0.7pt\hbox{.}\mkern1mu}}%
  \def\triangle{\mathord{\scalebox{0.9}{$\symup{\Delta}$}}}%
  \def\nabla{\mathord{\raisebox{\depth}{\scalebox{1}[-1]{$\symup{\Delta}$}}}}%
  \def\checkmark{\popcheck{popdark}}%
  \def\star{\mathbin{\ast}}\def\bigstar{\mathord{\ast}}%
  \def\diamond{\mathbin{\scalebox{0.6}{$\Diamond$}}}%
  \def\bigcup{\mathop{\mathchoice{\raisebox{-0.3em}{\scalebox{1.7}{$\cup$}}}{\scalebox{1.25}{$\cup$}}{\cup}{\cup}}\displaylimits}%
  \def\bigcap{\mathop{\mathchoice{\raisebox{-0.3em}{\scalebox{1.7}{$\cap$}}}{\scalebox{1.25}{$\cap$}}{\cap}{\cap}}\displaylimits}%
  \def\lesseqqgtr{\mathrel{\lessgtr}}\def\bigoplus{\mathop{\oplus}\displaylimits}%
}
\providecommand{\ssi}{si et seulement si} % abréviation fréquente
\AtBeginDocument{\providecommand{\wideparen}{\overparen}} % arc de cercle

%% FIGFIX-BEGIN v4 — correctifs globaux des figures (docs/onbuch-generation/tools/figfix)
\usepackage{adjustbox}\usepackage{etoolbox}
% toute figure TikZ/pgfplots plus large que la ligne est réduite à la largeur disponible
\BeforeBeginEnvironment{tikzpicture}{\begin{adjustbox}{max width=\linewidth}}
\AfterEndEnvironment{tikzpicture}{\end{adjustbox}}
% légendes pgfplots lisibles par-dessus les courbes
\pgfplotsset{every axis/.append style={legend style={fill=white,fill opacity=0.92,text opacity=1,draw=black!15,inner sep=3pt}}}
% étiquettes d'arêtes (syntaxe edge["..."]) avec fond blanc
\tikzset{every edge quotes/.append style={fill=white,inner sep=1.2pt}}
%% FIGFIX-END
\begin{document}

\pagegarde{Leçon 21 — Expression orale : vie culturelle au Cameroun et en Chine}{Chinois}{1ère A}{Module 3 : Citoyenneté et ouverture au monde}{ZHA21}{2 h}{Cette leçon d'expression orale apprend aux élèves de 1ère A à présenter un exposé simple et à participer à un débat guidé sur la vie culturelle au Cameroun et en Chine. Elle fournit le lexique culturel, les formules fonctionnelles pour présenter, opiner et relancer, ainsi qu'un entraînement ciblé sur les tons.
\vspace{4pt}
\begin{multicols}{2}\small
\begin{itemize}
\item À la fin de cette leçon, je sais utiliser le vocabulaire de base de la vie culturelle (fêtes, musique, danse, cinéma, cuisine) en chinois
\item À la fin de cette leçon, je sais présenter oralement un exposé simple de 8 à 10 phrases sur une fête camerounaise ou chinoise
\item À la fin de cette leçon, je sais donner mon avis avec 我觉得 / 我认为 et nuancer avec 但是 / 不过
\item À la fin de cette leçon, je sais exprimer l'accord et le désaccord poliment (我同意 / 我不太同意)
\item À la fin de cette leçon, je sais relancer un débat avec des questions ouvertes (你呢？你觉得怎么样？)
\item À la fin de cette leçon, je sais prononcer correctement les quatre tons et les distinguer dans des paires minimales
\end{itemize}\end{multicols}}

\begin{sommaire}
\begin{multicols}{2}
\begin{enumerate}
\item Lexique de la vie culturelle
\item Dialogue modèle 1 : présenter une fête
\item Formules fonctionnelles : présenter, opiner, relancer
\item Les tons : entraînement ciblé
\item Dialogue modèle 2 : donner son avis et débattre
\item Jeux de rôle et débat guidé
\item Activité d'intégration
\item Méthodes et exercices résolus
\item Exercices
\item Corrigés des exercices
\item Fiche bilan
\end{enumerate}
\end{multicols}
\end{sommaire}

\begin{objectifs}
\begin{itemize}
\item À la fin de cette leçon, je sais utiliser le vocabulaire de base de la vie culturelle (fêtes, musique, danse, cinéma, cuisine) en chinois
\item À la fin de cette leçon, je sais présenter oralement un exposé simple de 8 à 10 phrases sur une fête camerounaise ou chinoise
\item À la fin de cette leçon, je sais donner mon avis avec 我觉得 / 我认为 et nuancer avec 但是 / 不过
\item À la fin de cette leçon, je sais exprimer l'accord et le désaccord poliment (我同意 / 我不太同意)
\item À la fin de cette leçon, je sais relancer un débat avec des questions ouvertes (你呢？你觉得怎么样？)
\item À la fin de cette leçon, je sais prononcer correctement les quatre tons et les distinguer dans des paires minimales
\item À la fin de cette leçon, je sais comparer deux réalités culturelles avec 比 et 跟……一样
\item À la fin de cette leçon, je sais tenir un dialogue de 2 minutes avec un camarade sur un sujet culturel
\end{itemize}
\end{objectifs}

\begin{prerequis}
\begin{itemize}
\item Les quatre tons et le pinyin (Seconde)
\item Les pronoms personnels et les verbes courants 是、有、喜欢、觉得
\item La structure interrogative avec 吗、呢、什么、怎么
\item Le comparatif simple avec 比 (début d'année de 1ère)
\item Le vocabulaire des pays et nationalités (中国、喀麦隆)
\end{itemize}
\end{prerequis}

\newpage

\coursec{Lexique de la vie culturelle}

C'est la semaine culturelle au lycée de Yaoundé. Dans la cour, les stands se préparent : au stand n° 3, les élèves de Première A présentent les danses traditionnelles du Cameroun — le bikutsi du Centre et du Sud, le makossa du Littoral, les danses des Grassfields de l'Ouest. Un groupe d'étudiants chinois en échange à l'Université de Yaoundé I s'approche, curieux. L'un d'eux demande : \zh{你们有什么传统节日？} \emph{(Nǐmen yǒu shénme chuántǒng jiérì ?)} « Quelles fêtes traditionnelles avez-vous ? ». Pour répondre, il faut du vocabulaire : les mots des fêtes, de la musique, des spectacles, de la cuisine, et les adjectifs pour dire « c'est beau », « c'est animé », « c'est célèbre ». C'est précisément l'objectif de cette première section : construire le \cle{lexique de la vie culturelle}, le prononcer correctement, et savoir l'employer dans des phrases simples. Problématique : comment présenter la richesse culturelle du Cameroun en chinois ?

\courssub{Observer d'abord : un mini-texte de présentation}

Avant d'apprendre les mots un par un, lisons le texte que les élèves du stand n° 3 ont préparé pour accueillir leurs visiteurs chinois.

\begin{textebox}[Au stand de la semaine culturelle]
喀麦隆有很多传统节日，音乐很好听，菜也很好吃。\\
\emph{Kāmàilóng yǒu hěn duō chuántǒng jiérì, yīnyuè hěn hǎotīng, cài yě hěn hǎochī.}\\
« Le Cameroun a beaucoup de fêtes traditionnelles, la musique est belle et les plats sont bons. »\\
我们喜欢跳舞、唱歌。Makossa 和 bikutsi 很有意思！\\
\emph{Wǒmen xǐhuan tiàowǔ, chànggē. Makossa hé bikutsi hěn yǒuyìsi !}\\
« Nous aimons danser et chanter. Le makossa et le bikutsi sont très intéressants ! »
\end{textebox}

Observons ce texte. Quatre remarques :

\begin{enumerate}
\item On retrouve la structure déjà connue \cle{Sujet + 很 + adjectif} : \zh{音乐很好听} « la musique est belle », \zh{菜也很好吃} « les plats aussi sont bons ». En chinois, l'adjectif seul fait office de verbe : pas besoin de « être ».
\item L'adverbe \zh{也} (yě, « aussi ») se place \textbf{avant} l'adverbe de degré : \zh{菜也很好吃} « les plats aussi sont bons », jamais \zh{菜很也好吃}.
\item Le mot \zh{和} (hé, « et ») relie deux noms : \zh{makossa 和 bikutsi}. Attention : \zh{和} ne relie jamais deux phrases complètes, seulement des mots ou des groupes nominaux.
\item Les mots étrangers comme \emph{makossa} et \emph{bikutsi} s'intègrent directement dans la phrase chinoise, sans traduction : c'est la méthode que nous utiliserons pour parler des réalités camerounaises.
\end{enumerate}

\courssub{Le lexique : quatre familles de mots}

\textbf{Famille 1 : les fêtes et festivals}\par

\begin{center}
\begin{tabularx}{\linewidth}{|X|X|X|X|}
\hline
\rowcolor{popblueL} Caractères & Pinyin & Nature & Traduction \\
\hline
节日 & jiérì & nom & fête, jour de fête \\
\hline
春节 & Chūnjié & nom propre & Fête du Printemps (Nouvel An chinois) \\
\hline
中秋节 & Zhōngqiūjié & nom propre & Fête de la Lune (mi-automne) \\
\hline
文化节 & wénhuàjié & nom & semaine / festival culturel \\
\hline
传统 & chuántǒng & nom / adj. & tradition ; traditionnel \\
\hline
\end{tabularx}
\end{center}

Remarque de formation : le mot \zh{节} (jié, « fête ») sert de suffixe : \zh{春} (printemps) + \zh{节} = Fête du Printemps ; \zh{中秋} (mi-automne) + \zh{节} = Fête de la Lune ; \zh{文化} (culture) + \zh{节} = festival culturel. Comprendre ce mécanisme permet de deviner le sens de nombreux mots nouveaux.

\textbf{Famille 2 : les arts et les spectacles}\par

\begin{center}
\begin{tabularx}{\linewidth}{|X|X|X|X|}
\hline
\rowcolor{popblueL} Caractères & Pinyin & Nature & Traduction \\
\hline
音乐 & yīnyuè & nom & musique \\
\hline
跳舞 & tiàowǔ & verbe & danser \\
\hline
唱歌 & chànggē & verbe & chanter \\
\hline
京剧 & Jīngjù & nom propre & opéra de Pékin \\
\hline
电影 & diànyǐng & nom & film, cinéma \\
\hline
表演 & biǎoyǎn & nom / verbe & spectacle, représentation ; jouer, se produire \\
\hline
\end{tabularx}
\end{center}

Remarque de formation : \zh{跳} (sauter) + \zh{舞} (danse) = danser ; \zh{唱} (chanter) + \zh{歌} (chanson) = chanter. Ces verbes sont dits « verbe + objet figé » : on peut insérer un complément entre les deux caractères, par exemple \zh{跳一个舞} \emph{(tiào yí ge wǔ)} « danser une danse ».

\textbf{Famille 3 : la cuisine et la vie quotidienne}\par

\begin{center}
\begin{tabularx}{\linewidth}{|X|X|X|X|}
\hline
\rowcolor{popblueL} Caractères & Pinyin & Nature & Traduction \\
\hline
菜 & cài & nom & plat ; légumes \\
\hline
米饭 & mǐfàn & nom & riz (cuit) \\
\hline
饺子 & jiǎozi & nom & raviolis chinois \\
\hline
好吃 & hǎochī & adjectif & bon à manger, délicieux \\
\hline
\end{tabularx}
\end{center}

\textbf{Famille 4 : les adjectifs d'appréciation}\par

\begin{center}
\begin{tabularx}{\linewidth}{|X|X|X|X|}
\hline
\rowcolor{popblueL} Caractères & Pinyin & Nature & Traduction \\
\hline
有意思 & yǒuyìsi & adjectif & intéressant, amusant \\
\hline
漂亮 & piàoliang & adjectif & beau, joli \\
\hline
有名 & yǒumíng & adjectif & célèbre \\
\hline
热闹 & rènao & adjectif & animé, bruyant et joyeux \\
\hline
特别 & tèbié & adverbe / adj. & spécial ; particulièrement \\
\hline
\end{tabularx}
\end{center}

\begin{definition}[热闹 (rènao) : le mot des fêtes]
\zh{热闹} est un adjectif typique pour décrire une fête, un marché, un stade : il signifie « animé, plein de monde, bruyant et joyeux ». Il est formé de \zh{热} (chaud) et \zh{闹} (bruyant) : littéralement « chaud et bruyant ». On l'emploie avec la structure Sujet + 很 + 热闹 : \zh{春节很热闹。} \emph{(Chūnjié hěn rènao.)} « La Fête du Printemps est très animée. » À Yaoundé, on dira : \zh{市场很热闹。} \emph{(Shìchǎng hěn rènao.)} « Le marché est très animé. »
\end{definition}

\begin{definition}[La famille des 好 + verbe : apprécier avec les sens]
Le chinois forme des adjectifs d'appréciation avec \zh{好} (hǎo, « bon ») + un verbe de perception : \zh{好吃} (hǎochī, « bon à manger », de \zh{吃} manger), \zh{好听} (hǎotīng, « beau à écouter », de \zh{听} écouter), \zh{好看} (hǎokàn, « beau à voir », de \zh{看} regarder). Ainsi : la musique est \zh{好听}, un plat est \zh{好吃}, un film ou un spectacle est \zh{好看}. Choisir le bon adjectif, c'est choisir le bon sens : on ne dit pas \zh{音乐很好吃} !
\end{definition}

\courssub{Carte mentale du lexique}

\begin{popfigure}
\begin{center}
\begin{tikzpicture}[
  mindmap,
  grow cyclic,
  every node/.style={concept, align=center, font=\small},
  root concept/.append style={concept color=popblue, font=\bfseries},
  level 1/.append style={level distance=3.2cm, sibling angle=90},
  level 2/.append style={level distance=2.4cm, sibling angle=45, font=\scriptsize}
]
\node[root concept, align=center]{文化生活\\ \emph{wénhuà shēnghuó}\\ vie culturelle}
  child[concept color=poporange]{ node[align=center]{节日\\ fêtes}
    child { node[align=center]{春节\\ Fête du Printemps} }
    child { node[align=center]{中秋节\\ Fête de la Lune} }
    child { node[align=center]{文化节\\ festival culturel} }
  }
  child[concept color=popgreen]{ node[align=center]{音乐舞蹈\\ musique et danse}
    child { node {音乐 yīnyuè} }
    child { node {跳舞 tiàowǔ} }
    child { node {唱歌 chànggē} }
  }
  child[concept color=poppurple]{ node[align=center]{电影表演\\ cinéma et spectacles}
    child { node {京剧 Jīngjù} }
    child { node {电影 diànyǐng} }
    child { node {表演 biǎoyǎn} }
  }
  child[concept color=popgold]{ node[align=center]{美食\\ cuisine}
    child { node {米饭 mǐfàn} }
    child { node {饺子 jiǎozi} }
    child { node {好吃 hǎochī} }
  };
\end{tikzpicture}
\end{center}
\legende{Carte mentale du lexique de la vie culturelle : quatre familles de mots autour du thème central.}
\end{popfigure}

\courssub{La grammaire du lexique : apprécier avec 很 + adjectif}

\begin{propriete}[Structure d'appréciation : Sujet + 很 + adjectif]
Pour donner une appréciation simple, le chinois utilise : \textbf{Sujet + 很 + adjectif (+ 。)}. Ici, \zh{很} (hěn) ne signifie pas toujours « très » : il sert surtout de liaison obligatoire entre le sujet et l'adjectif. Sans \zh{很}, la phrase sonne comparative ou incorrecte.\\
Exemples :\\
\zh{京剧很有名。} \emph{(Jīngjù hěn yǒumíng.)} « L'opéra de Pékin est (très) célèbre. »\\
\zh{饺子很好吃。} \emph{(Jiǎozi hěn hǎochī.)} « Les raviolis sont bons. »\\
\zh{文化节特别热闹。} \emph{(Wénhuàjié tèbié rènao.)} « Le festival culturel est particulièrement animé. »
\end{propriete}

Comparaison avec le français : le français exige le verbe « être » (« la fête \emph{est} animée ») ; le chinois n'en a pas besoin, l'adjectif fonctionne comme un verbe. En revanche, le chinois exige presque toujours un adverbe de degré (\zh{很}, \zh{特别}, \zh{真} zhēn « vraiment ») devant l'adjectif. Pour la négation, on remplace \zh{很} par \zh{不} : \zh{这个电影不好看。} \emph{(Zhège diànyǐng bù hǎokàn.)} « Ce film n'est pas beau (pas intéressant à regarder). »

\begin{exemplebox}[Exemples traduits]
\zh{春节很热闹。} \emph{(Chūnjié hěn rènao.)} « La Fête du Printemps est très animée. »\\
\zh{京剧很有名。} \emph{(Jīngjù hěn yǒumíng.)} « L'opéra de Pékin est très célèbre. »\\
\zh{这个电影很有意思。} \emph{(Zhège diànyǐng hěn yǒuyìsi.)} « Ce film est très intéressant. »\\
\zh{中国的菜特别好吃。} \emph{(Zhōngguó de cài tèbié hǎochī.)} « La cuisine chinoise est particulièrement bonne. »\\
\zh{她的表演很好看。} \emph{(Tā de biǎoyǎn hěn hǎokàn.)} « Son spectacle est très beau (à regarder). »\\
\zh{这个表演真有意思！} \emph{(Zhège biǎoyǎn zhēn yǒuyìsi !)} « Ce spectacle est vraiment intéressant ! »
\end{exemplebox}

\courssub{Application au Cameroun : intégrer makossa et bikutsi}

Pour parler des réalités camerounaises qui n'ont pas de mot chinois courant, on insère le mot tel quel dans la phrase chinoise, comme un nom propre. C'est exactement ce que font les Chinois eux-mêmes.

\begin{exemplebox}[Parler du Cameroun en chinois]
\zh{我喜欢喀麦隆的 makossa。} \emph{(Wǒ xǐhuan Kāmàilóng de makossa.)} « J'aime le makossa du Cameroun. »\\
\zh{Bikutsi 很有意思。} « Le bikutsi est très intéressant. »\\
\zh{喀麦隆的音乐很好听。} \emph{(Kāmàilóng de yīnyuè hěn hǎotīng.)} « La musique camerounaise est belle. »\\
\zh{Ndolé 特别好吃。} « Le ndolé est particulièrement bon. »
\end{exemplebox}

\begin{center}
\begin{tabularx}{\linewidth}{|X|X|X|}
\hline
\rowcolor{popblueL} Domaine & Cameroun & Chine \\
\hline
Musique et danse & makossa, bikutsi, danse des Grassfields & 京剧 Jīngjù (opéra de Pékin), danses du dragon \\
\hline
Grande fête & ngondo (peuple sawa), fêtes traditionnelles des chefferies & 春节 Chūnjié (Fête du Printemps) \\
\hline
Fête de récolte & fêtes de l'igname dans l'Ouest & 中秋节 Zhōngqiūjié (Fête de la Lune) \\
\hline
Plat emblématique & ndolé, poulet DG, eru & 饺子 jiǎozi (raviolis), 米饭 mǐfàn (riz) \\
\hline
\end{tabularx}
\end{center}

\courssub{Prononciation guidée : travailler les tons}

Pour chaque mot, procédons en trois temps : écoute du modèle du professeur, répétition chorale de toute la classe, puis répétition individuelle de trois ou quatre élèves. Points de vigilance :

\begin{itemize}
\item \zh{音乐} yīnyuè : 1er ton (haut et plat) puis 4e ton (descendant). Le premier caractère ne descend pas.
\item \zh{跳舞} tiàowǔ : 4e ton puis 3e ton (bas, qui descend puis remonte).
\item \zh{饺子} jiǎozi : 3e ton puis ton neutre (zi, court et léger).
\item \zh{热闹} rènao : 4e ton puis ton neutre.
\item \zh{有意思} yǒuyìsi : 3e ton, 4e ton, ton neutre.
\end{itemize}

\begin{attention}[Erreur fréquente des francophones : 音乐]
Beaucoup d'élèves francophones prononcent \zh{音乐} avec deux tons descendants, comme si c'était « yìnyuè ». Or le premier caractère \zh{音} est au \textbf{1er ton} : haut et plat, comme une note tenue. Dites \emph{yīn} en tenant la voix, puis laissez tomber sur \emph{yuè}. Autre piège : \zh{好吃} se prononce hǎochī (3e + 1er ton), jamais hàochī ; \zh{好} au 4e ton (hào) signifie « aimer », comme dans \zh{爱好} (passe-temps).
\end{attention}

\begin{savaistu}[L'opéra de Pékin, trésor de l'humanité]
L'opéra de Pékin (\zh{京剧} Jīngjù) est inscrit au patrimoine culturel immatériel de l'humanité de l'UNESCO depuis 2010. Ses acteurs portent des maquillages codés par couleurs : le rouge pour la loyauté, le noir pour la droiture, le blanc pour la ruse. Comme lui, plusieurs danses et savoir-faire traditionnels africains sont reconnus par l'UNESCO : la culture orale et les danses de nombreux peuples d'Afrique centrale et de l'Ouest figurent aussi sur ces listes. Parler de \zh{京剧} à un correspondant chinois, c'est donc lui parler de son « makossa à lui » : un art populaire devenu symbole national.
\end{savaistu}

\begin{exoresolu}[Compléter avec le bon mot du lexique]
Énoncé : complétez chaque phrase avec le mot qui convient : \zh{热闹}, \zh{有名}, \zh{好吃}, \zh{有意思}, \zh{音乐}.\\
1. 春节很 \_\_\_\_ 。 (La Fête du Printemps est très animée.)\\
2. 京剧很 \_\_\_\_ 。 (L'opéra de Pékin est très célèbre.)\\
3. 饺子很 \_\_\_\_ 。 (Les raviolis sont très bons.)\\
4. 这个电影很 \_\_\_\_ 。 (Ce film est très intéressant.)\\
5. 喀麦隆的 \_\_\_\_ 很好听。 (La musique du Cameroun est belle.)
\tcblower
Solution détaillée :\\
1. \zh{春节很热闹。} — on décrit une fête pleine de monde : \zh{热闹} est l'adjectif typique des fêtes.\\
2. \zh{京剧很有名。} — « célèbre » se dit \zh{有名} (littéralement « avoir un nom »).\\
3. \zh{饺子很好吃。} — pour la nourriture, « bon à manger » = \zh{好吃}.\\
4. \zh{这个电影很有意思。} — « intéressant » = \zh{有意思} (littéralement « avoir du sens »).\\
5. \zh{喀麦隆的音乐很好听。} — « musique » = \zh{音乐} ; on l'apprécie avec \zh{好听} « beau à écouter ».\\
Méthode : repérer d'abord le domaine (fête, art, nourriture, film, musique), puis choisir l'adjectif ou le nom de la bonne famille de la carte mentale.
\end{exoresolu}

\begin{exercice}[À vous de parler]
En binôme, préparez quatre phrases pour le stand de la semaine culturelle : une sur une fête camerounaise, une sur le makossa ou le bikutsi, une sur un plat camerounais, une sur l'opéra de Pékin. Utilisez la structure Sujet + 很 / 特别 + adjectif. Prononcez chaque phrase à voix haute en veillant aux tons ; votre binôme corrige la prononciation.
\end{exercice}

\begin{corrige}[Exercice « À vous de parler »]
Modèles possibles (toute phrase correcte avec la structure est acceptée) :\\
\zh{Ngondo 很热闹。} « Le ngondo est très animé. »\\
\zh{我喜欢喀麦隆的 bikutsi。} « J'aime le bikutsi du Cameroun. »\\
\zh{Ndolé 特别好吃。} « Le ndolé est particulièrement bon. »\\
\zh{京剧很有名，也很好看。} « L'opéra de Pékin est célèbre et beau à voir. »\\
Points de correction : présence de \zh{很} ou \zh{特别} devant l'adjectif ; tons de \zh{音乐} (1er ton plat) et de \zh{好吃} (hǎo, 3e ton) ; ordre Sujet + adverbe + adjectif respecté ; \zh{也} placé avant \zh{很}.
\end{corrige}

\begin{aretenir}[L'essentiel de la section]
\begin{itemize}
\item Quatre familles de mots : fêtes (\zh{节日}, \zh{春节}, \zh{中秋节}, \zh{文化节}), arts (\zh{音乐}, \zh{跳舞}, \zh{唱歌}, \zh{京剧}, \zh{电影}, \zh{表演}), cuisine (\zh{菜}, \zh{米饭}, \zh{饺子}, \zh{好吃}), appréciation (\zh{有意思}, \zh{漂亮}, \zh{有名}, \zh{热闹}, \zh{特别}).
\item Structure d'appréciation : Sujet + 很 / 特别 / 真 + adjectif, sans verbe « être » ; négation avec 不.
\item La famille 好 + verbe de perception : \zh{好吃} (manger), \zh{好听} (écouter), \zh{好看} (regarder).
\item Les mots camerounais (makossa, bikutsi, ndolé) s'insèrent directement dans la phrase chinoise.
\item Prononciation : \zh{音乐} = 1er ton plat + 4e ton ; \zh{好吃} = hǎochī, jamais hàochī.
\end{itemize}
\end{aretenir}

\coursec{Lexique de la vie culturelle}

Avant d'aborder le dialogue, il est indispensable de maîtriser le vocabulaire de base pour parler des fêtes, des lieux, des actions et des impressions. Ce lexique couvre les notions du programme (HSK 2-3) et le contexte camerounais.

\begin{center}
\begin{tabularx}{\linewidth}{|X|X|X|X|}
\hline
\rowcolor{popblueL} Caractères & Pinyin & Nature & Traduction \\
\hline
节日 & jiérì & nom & fête, festival \\
\hline
传统节日 & chuántǒng jiérì & nom & fête traditionnelle \\
\hline
文化节 & wénhuà jié & nom & semaine culturelle, festival culturel \\
\hline
春节 & Chūnjié & nom propre & Fête du Printemps (Nouvel An lunaire) \\
\hline
中秋节 & Zhōngqiūjié & nom propre & Fête de la Mi-Automne \\
\hline
圣诞节 & Shèngdànjié & nom propre & Noël \\
\hline
Ngondo 节 & Ngondo jié & nom propre & Fête du Ngondo (peuple Sawa) \\
\hline
赛龙舟 & sài lóngzhōu & verbe + obj & faire la course de bateaux-dragons \\
\hline
海边 & hǎibiān & nom & bord de mer, plage \\
\hline
教堂 & jiàotáng & nom & église \\
\hline
家里 & jiā lǐ & nom & à la maison, chez soi \\
\hline
举行 & jǔxíng & verbe & organiser, tenir (cérémonie, fête, réunion) \\
\hline
庆祝 & qìngzhù & verbe & célébrer, fêter \\
\hline
吃饺子 & chī jiǎozi & verbe + obj & manger des raviolis \\
\hline
吃月饼 & chī yuèbǐng & verbe + obj & manger des gâteaux de lune \\
\hline
看表演 & kàn biǎoyǎn & verbe + obj & regarder un spectacle \\
\hline
看月亮 & kàn yuèliang & verbe + obj & admirer la lune \\
\hline
放鞭炮 & fàng biānpào & verbe + obj & tirer des pétards / faire exploser des feux d'artifice \\
\hline
热闹 & rènao & adj & animé, bruyant, vivant \\
\hline
有意思 & yǒu yìsi & expr & intéressant \\
\hline
重要 & zhòngyào & adj & important \\
\hline
欢迎 & huānyíng & verbe & accueillir, souhaiter la bienvenue \\
\hline
比如 & bǐrú & conjonction & par exemple \\
\hline
时候 & shíhou & nom & moment, quand (après une phrase) \\
\hline
\end{tabularx}
\end{center}

\begin{aretenir}[Classificateurs utiles pour les événements]
\begin{itemize}
\item \zh{个} (gè) : universel. \zh{一个节日} (un festival), \zh{一个表演} (un spectacle).
\item \zh{场} (chǎng) : pour les événements publics, spectacles, matchs. \zh{一场比赛} (un match), \zh{一场表演} (une représentation).
\item \zh{次} (cì) : pour le nombre de fois. \zh{去过两次} (y être allé deux fois).
\end{itemize}
\end{aretenir}

\coursec{Dialogue modèle 1 : présenter une fête}

Nous mettons maintenant ce lexique en action. La situation : une élève camerounaise accueille un étudiant chinois pendant la semaine culturelle de son lycée à Douala et lui présente la fête du \cle{Ngondo}. Lis d'abord le dialogue en silence, puis à voix haute, avant de l'étudier phrase par phrase.

\courssub{Le dialogue à lire et à écouter}

\begin{textebox}[À la semaine culturelle du lycée — 在文化节上]
\textbf{A}（élève camerounaise）：你好！欢迎来到我们的文化节。\\
\emph{Nǐ hǎo! Huānyíng lái dào wǒmen de wénhuà jié.}\\
« Bonjour ! Bienvenue à notre semaine culturelle. »\\[0.4em]
\textbf{B}（étudiant chinois）：谢谢！喀麦隆有什么传统节日？\\
\emph{Xièxie! Kāmàilóng yǒu shénme chuántǒng jiérì?}\\
« Merci ! Quelles fêtes traditionnelles y a-t-il au Cameroun ? »\\[0.4em]
\textbf{A}：有很多，比如 Ngondo 节，是 Sawa 人的节日，在海边举行，很热闹。\\
\emph{Yǒu hěn duō, bǐrú Ngondo jié, shì Sawa rén de jiérì, zài hǎibiān jǔxíng, hěn rènao.}\\
« Il y en a beaucoup, par exemple la fête Ngondo, fête du peuple Sawa, célébrée au bord de la mer, très animée. »\\[0.4em]
\textbf{B}：有意思！中国最重要的节日是春节。\\
\emph{Yǒu yìsi! Zhōngguó zuì zhòngyào de jiérì shì Chūnjié.}\\
« Intéressant ! La fête la plus importante en Chine est la Fête du Printemps. »\\[0.4em]
\textbf{A}：春节的时候你们做什么？\\
\emph{Chūnjié de shíhou nǐmen zuò shénme?}\\
« Que faites-vous pendant la Fête du Printemps ? »\\[0.4em]
\textbf{B}：我们吃饺子、看表演、放鞭炮。\\
\emph{Wǒmen chī jiǎozi, kàn biǎoyǎn, fàng biānpào.}\\
« Nous mangeons des raviolis, regardons des spectacles et tirons des pétards. »
\end{textebox}

\begin{methode}[Comment travailler ce dialogue à voix haute]
\begin{enumerate}
\item \textbf{Première lecture} : lis le dialogue à voix haute en suivant le \cle{pinyin}, lentement, en marquant bien chaque ton.
\item \textbf{Deuxième lecture} : lis en regardant \emph{uniquement les caractères} ; ne redescends vers le pinyin qu'en cas d'hésitation.
\item \textbf{Troisième lecture} : ferme les yeux sur la traduction et vérifie que tu comprends chaque phrase globalement, sans traduire mot à mot.
\end{enumerate}
\end{methode}

\courssub{Observation : les cinq étapes de l'échange}

Avant toute règle, observons la \emph{logique} du dialogue. Il ne s'agit pas de phrases jetées au hasard : chaque réplique a une fonction précise.

\begin{center}
\begin{tabularx}{\linewidth}{|X|X|X|}
\hline
\rowcolor{popblueL} Étape & Réplique (caractères + pinyin) & Fonction \\
\hline
1. Saluer & 你好！欢迎来到我们的文化节。\newline \emph{Nǐ hǎo! Huānyíng lái dào wǒmen de wénhuà jié.} & Accueillir l'invité \\
\hline
2. Poser une question & 喀麦隆有什么传统节日？\newline \emph{Kāmàilóng yǒu shénme chuántǒng jiérì?} & S'intéresser à l'autre \\
\hline
3. Présenter & 有很多，比如 Ngondo 节……\newline \emph{Yǒu hěn duō, bǐrú Ngondo jié...} & Donner l'information \\
\hline
4. Réagir & 有意思！\newline \emph{Yǒu yìsi!} & Montrer son intérêt \\
\hline
5. Relancer & 春节的时候你们做什么？\newline \emph{Chūnjié de shíhou nǐmen zuò shénme?} & Faire parler l'autre à son tour \\
\hline
\end{tabularx}
\end{center}

\begin{aretenir}[Un bon échange = un escalier à cinq marches]
Un échange oral réussi monte toujours les mêmes marches : \cle{saluer} → \cle{poser une question} → \cle{présenter} → \cle{réagir} → \cle{relancer}. Celui qui parle ne monopolise jamais la parole : il rend la main par une question.
\end{aretenir}

\begin{popfigure}
\begin{tikzpicture}[
  stepS/.style={rectangle, draw=popblue, fill=popblueL, rounded corners=2pt,
               minimum width=3.1cm, minimum height=0.85cm, align=center, font=\small},
  arr/.style={-{Stealth[length=2.5mm]}, thick, poporange}
]
\node[stepS, align=center] (s1) at (0,0){\textbf{1. Saluer}\\ 你好！欢迎！};
\node[stepS, align=center] (s2) at (3.4,0.9){\textbf{2. Question}\\ 有什么节日？};
\node[stepS, align=center] (s3) at (6.8,1.8){\textbf{3. Présenter}\\ 比如 Ngondo 节……};
\node[stepS, align=center] (s4) at (3.4,2.7){\textbf{4. Réagir}\\ 有意思！};
\node[stepS, align=center] (s5) at (0,3.6){\textbf{5. Relancer}\\ 你们做什么？};
\draw[arr] (s1.east) -- (s2.west);
\draw[arr] (s2.east) -- (s3.west);
\draw[arr] (s3.west) -- (s4.east);
\draw[arr] (s4.west) -- (s5.east);
\end{tikzpicture}
\legende{Les cinq étapes d'un échange réussi : la conversation monte comme un escalier.}
\end{popfigure}

\courssub{Étude des structures du dialogue}

\textbf{1. Accueillir : 欢迎来到……}

\zh{欢迎来到我们的文化节。} \emph{Huānyíng lái dào wǒmen de wénhuà jié.} « Bienvenue à notre semaine culturelle. »

\begin{propriete}[La formule 欢迎 + (来到) + lieu]
Structure : \cle{欢迎} (+ \cle{来到}) + \cle{lieu / événement}.
\begin{itemize}
\item \zh{欢迎} seul signifie « bienvenue ! / soyez le bienvenu ! » : \zh{欢迎！} \emph{Huānyíng!}
\item Pour préciser la destination, on ajoute \zh{来到} (arriver à) ou directement le lieu : \zh{欢迎来喀麦隆！} \emph{Huānyíng lái Kāmàilóng!} « Bienvenue au Cameroun ! »
\item Le possesseur se construit avec \zh{的} : \zh{我们的文化节} « notre semaine culturelle ».
\end{itemize}
Exemples :
\begin{itemize}
\item \zh{欢迎来到我们学校！} \emph{Huānyíng lái dào wǒmen xuéxiào!} « Bienvenue dans notre école ! »
\item \zh{欢迎你来杜阿拉！} \emph{Huānyíng nǐ lái Dù'ālā!} « Bienvenue à Douala ! » (on peut insérer la personne accueillie : 欢迎你).
\end{itemize}
\end{propriete}

\begin{attention}[La prononciation de 欢迎]
\zh{欢迎} se prononce \cle{huānyíng} : \textbf{1\textsuperscript{er} ton + 2\textsuperscript{e} ton}. Le premier caractère \zh{欢} (huān, joyeux) est \emph{plat et haut}, pas descendant. Erreur typique des francophones : dire « huànyíng » (4\textsuperscript{e} ton), ce qui change le sens. Répète : huān (à plat, comme un chant tenu) — yíng (qui monte, comme une question).
\end{attention}

\textbf{2. Demander et donner un exemple : 有什么……？ et 比如}

\zh{喀麦隆有什么传统节日？} \emph{Kāmàilóng yǒu shénme chuántǒng jiérì?} « Quelles fêtes traditionnelles y a-t-il au Cameroun ? »

La question repose sur \zh{有} (yǒu, il y a / avoir) + \zh{什么} (shénme, quoi, quel). Contrairement au français, le mot interrogatif \zh{什么} reste \emph{à la place de la réponse}, en fin de phrase : littéralement « Le Cameroun a quelles fêtes traditionnelles ? ». On ne déplace jamais \zh{什么} en tête de phrase.

Pour répondre en citant un exemple, le chinois utilise \cle{比如} (bǐrú, par exemple), placé devant l'exemple.

\begin{exemplebox}[比如 = par exemple]
\begin{itemize}
\item \zh{喀麦隆有很多节日，比如 Ngondo。} \emph{Kāmàilóng yǒu hěn duō jiérì, bǐrú Ngondo.} « Le Cameroun a beaucoup de fêtes, par exemple le Ngondo. »
\item \zh{中国有很多传统节日，比如春节和中秋节。} \emph{Zhōngguó yǒu hěn duō chuántǒng jiérì, bǐrú Chūnjié hé Zhōngqiūjié.} « La Chine a beaucoup de fêtes traditionnelles, par exemple la Fête du Printemps et la fête de la Mi-Automne. »
\item \zh{我喜欢喀麦隆菜，比如 ndolé。} \emph{Wǒ xǐhuan Kāmàilóng cài, bǐrú ndolé.} « J'aime la cuisine camerounaise, par exemple le ndolé. »
\end{itemize}
\end{exemplebox}

\begin{attention}[比如 vs 例如]
\zh{比如} (bǐrú) s'emploie à l'oral et à l'écrit courant. \zh{例如} (lìrú) est plus formel, écrit. En expression orale (HSK 3-4), privilégie \zh{比如}.
\end{attention}

\textbf{3. Localiser un événement : 在 + lieu + 举行}

\zh{在海边举行} \emph{zài hǎibiān jǔxíng} « se tient au bord de la mer ». Observons : en français, « la fête se tient au bord de la mer » place le lieu \emph{après} le verbe ; en chinois, le complément de lieu introduit par \zh{在} se place \emph{avant} le verbe.

\begin{propriete}[Structure 在 + lieu + 举行 (se tenir, avoir lieu)]
Structure : \cle{Sujet} + \cle{在} + \cle{lieu} + \cle{举行}.\\
\zh{举行} (jǔxíng) signifie « organiser, tenir (une cérémonie, une fête, une réunion) » ; à la voix passive française « se tenir », le chinois dit simplement : sujet + 在 + lieu + 举行.
\begin{itemize}
\item \zh{节日在海边举行。} \emph{Jiérì zài hǎibiān jǔxíng.} « La fête se tient au bord de la mer. »
\item \zh{春节在中国举行。} \emph{Chūnjié zài Zhōngguó jǔxíng.} « La Fête du Printemps se déroule en Chine. »
\item \zh{比赛在雅温得举行。} \emph{Bǐsài zài Yǎwēndé jǔxíng.} « Le match a lieu à Yaoundé. »
\end{itemize}
Règle générale à retenir : en chinois, \textbf{le complément de lieu avec 在 se met toujours avant le verbe}, jamais après.
\end{propriete}

\begin{attention}[Piège de traduction : la place du lieu]
Ne dis jamais \zh{举行在海边} sur le modèle français « se tenir au bord de la mer ». L'ordre chinois est : \zh{在海边举行}. De même : \zh{我在杜阿拉学习} \emph{Wǒ zài Dù'ālā xuéxí.} « J'étudie à Douala » — et non \zh{我学习在杜阿拉}.
\end{attention}

\textbf{4. Décrire l'ambiance : 很热闹}

\zh{很热闹} \emph{hěn rènao} « très animé ». \zh{热闹} (rènao) est un adjectif d'état très courant pour décrire une ambiance festive, un marché, une rue pleine de monde. \zh{很} (hěn) fait office de lien sujet-adjectif (pas de verbe « être »).

\textbf{5. Réagir et comparer : 有意思！ et 最重要的}

\zh{有意思！} \emph{Yǒu yìsi!} « Intéressant ! » (litt. « il y a du sens »). C'est la réaction la plus naturelle pour montrer son intérêt. On peut aussi dire \zh{很好！} « Très bien ! » ou \zh{真的吗？} \emph{Zhēn de ma?} « Vraiment ? ».

\zh{中国最重要的节日是春节。} \emph{Zhōngguó zuì zhòngyào de jiérì shì Chūnjié.} « La fête la plus importante en Chine est la Fête du Printemps. » Le superlatif se forme avec \cle{最} (zuì, le plus) + adjectif + \zh{的} + nom : \zh{最重要的节日} « la fête la plus importante ».

\begin{propriete}[Superlatif : 最 + adjectif + 的 + nom]
Structure : \cle{最} + \cle{adjectif} + \cle{的} + \cle{nom}.
\begin{itemize}
\item \zh{最热闹的节日} \emph{zuì rènao de jiérì} « la fête la plus animée »
\item \zh{最好吃的菜} \emph{zuì hǎochī de cài} « le plat le plus délicieux »
\end{itemize}
\end{propriete}

\textbf{6. Situer dans le temps : ……的时候}

\zh{春节的时候} \emph{Chūnjié de shíhou} « pendant la Fête du Printemps / au moment de la Fête du Printemps ». La structure \cle{Nom/Verbe + 的时候} signifie « quand / au moment où ». Elle se place généralement en début de phrase (sujet de temps) ou après le sujet.

\begin{exemplebox}[……的时候]
\begin{itemize}
\item \zh{我小时候喜欢看电视。} \emph{Wǒ xiǎo de shíhou xǐhuan kàn diànshì.} « Quand j'étais petit, j'aimais regarder la télé. »
\item \zh{吃饭的时候不说话。} \emph{Chīfàn de shíhou bù shuōhuà.} « On ne parle pas pendant le repas. »
\end{itemize}
\end{exemplebox}

\textbf{7. Énumérer des actions : 、 entre les verbes}

\zh{我们吃饺子、看表演、放鞭炮。} \emph{Wǒmen chī jiǎozi, kàn biǎoyǎn, fàng biānpào.} « Nous mangeons des raviolis, regardons des spectacles et tirons des pétards. » La virgule chinoise \zh{、} (dùnhào) sépare les éléments d'une énumération ; on ne met pas de « et » (和) entre des verbes en série — \zh{和} ne relie que des noms.

\begin{attention}[和 vs 、]
\begin{itemize}
\item Noms : \zh{饺子和月饼} (raviolis \textbf{et} gâteaux de lune) — \zh{和} obligatoire.
\item Verbes / propositions : \zh{吃饺子、看表演} (manger des raviolis, regarder un spectacle) — \zh{、} (ou simple juxtaposition), \textbf{pas de 和}.
\end{itemize}
\end{attention}

\coursec{Formules fonctionnelles : présenter, opiner, relancer}

Au-delà du dialogue unique, tu dois pouvoir varier tes formulations. Voici les boîtes à outils pour les trois fonctions clés de l'échange oral (niveau HSK 3 / Première A).

\courssub{Présenter un événement / une fête}

\begin{center}
\begin{tabularx}{\linewidth}{|X|X|X|}
\hline
\rowcolor{popblueL} Structure / Formule & Exemple (caractères + pinyin) & Traduction \\
\hline
……是……的节日。 & Ngondo 节是 Sawa 人的节日。\newline \emph{Ngondo jié shì Sawa rén de jiérì.} & Le Ngondo est la fête du peuple Sawa. \\
\hline
……在……举行。 & 节日在海边举行。\newline \emph{Jiérì zài hǎibiān jǔxíng.} & La fête se tient au bord de la mer. \\
\hline
……的时候，我们……。 & 春节的时候，我们吃饺子。\newline \emph{Chūnjié de shíhou, wǒmen chī jiǎozi.} & Pendant la Fête du Printemps, nous mangeons des raviolis. \\
\hline
有很多……，比如……。 & 有很多传统活动，比如赛龙舟。\newline \emph{Yǒu hěn duō chuántǒng huódòng, bǐrú sài lóngzhōu.} & Il y a beaucoup d'activités traditionnelles, par exemple la course de bateaux-dragons. \\
\hline
\end{tabularx}
\end{center}

\courssub{Donner son avis / Réagir (Opiner)}

\begin{center}
\begin{tabularx}{\linewidth}{|X|X|X|}
\hline
\rowcolor{popgreenL} Structure / Formule & Exemple (caractères + pinyin) & Traduction \\
\hline
有意思！/ 很有意思！ & 有意思！\newline \emph{Yǒu yìsi!} & Intéressant ! / Très intéressant ! \\
\hline
我觉得……很……。 & 我觉得春节很热闹。\newline \emph{Wǒ juéde Chūnjié hěn rènao.} & Je trouve que la Fête du Printemps est très animée. \\
\hline
听起来…… & 听起来很热闹。\newline \emph{Tīng qǐlái hěn rènao.} & Ça a l'air animé / Ça sonne bien. \\
\hline
真……啊！ & 真热闹啊！\newline \emph{Zhēn rènao a!} & C'est vraiment animé ! \\
\hline
\end{tabularx}
\end{center}

\courssub{Relancer / Poursuivre la conversation}

\begin{center}
\begin{tabularx}{\linewidth}{|X|X|X|}
\hline
\rowcolor{poporangeL} Structure / Formule & Exemple (caractères + pinyin) & Traduction \\
\hline
……的时候你们做什么？ & 春节的时候你们做什么？\newline \emph{Chūnjié de shíhou nǐmen zuò shénme?} & Que faites-vous pendant la Fête du Printemps ? \\
\hline
你们怎么过……？ & 你们怎么过中秋节？\newline \emph{Nǐmen zěnme guò Zhōngqiūjié?} & Comment fêtez-vous la Mi-Automne ? \\
\hline
还有别的活动吗？ & 还有别的活动吗？\newline \emph{Hái yǒu bié de huódòng ma?} & Y a-t-il d'autres activités ? \\
\hline
你喜欢……吗？ & 你喜欢吃月饼吗？\newline \emph{Nǐ xǐhuan chī yuèbǐng ma?} & Tu aimes manger des gâteaux de lune ? \\
\hline
\end{tabularx}
\end{center}

\begin{aretenir}[Mots-outils de liaison à l'oral]
Pour fluidifier ton exposé ou ton intervention dans un débat, utilise : \zh{然后} (ránhòu, ensuite), \zh{另外} (lìngwài, de plus), \zh{但是} (dànshì, mais), \zh{因为……所以……} (yīnwèi... suǒyǐ..., parce que... donc...).
\end{aretenir}

\coursec{Les tons : entraînement ciblé}

La prononciation des tons est la condition \emph{sine qua non} de l'intelligibilité. En Première A, on exige une maîtrise des quatre tons + ton neutre sur le vocabulaire connu. Travaille les séries ci-dessous en t'enregistrant si possible.

\courssub{Rappel des quatre tons + ton neutre}

\begin{center}
\begin{tabularx}{\linewidth}{|c|X|X|}
\hline
\rowcolor{poppurpleL} Ton & Description & Exemples du lexique \\
\hline
1\textsuperscript{er} (ˉ) & Haut, plat, long & \zh{欢} (huān), \zh{春} (chūn), \zh{中} (zhōng), \zh{放} (fàng), \zh{饺} (jiǎo - citation) \\
\hline
2\textsuperscript{e} (ˊ) & Montant, interrogatif & \zh{迎} (yíng), \zh{什么} (shénme \rightarrow shénme), \zh{如} (rú), \zh{举} (jǔ - citation), \zh{春} (chūn) \\
\hline
3\textsuperscript{e} (ˇ) & Descendant-montant (bas) & \zh{比} (bǐ), \zh{海} (hǎi), \zh{很} (hěn), \zh{饺} (jiǎo), \zh{看} (kàn), \zh{表} (biǎo) \\
\hline
4\textsuperscript{e} (ˋ) & Court, sec, descendant & \zh{节} (jié), \zh{日} (rì), \zh{在} (zài), \zh{行} (xíng), \zh{热} (rè), \zh{闹} (nào), \zh{意} (yì), \zh{思} (sì), \zh{最} (zuì), \zh{重} (zhòng), \zh{要} (yào), \zh{放} (fàng), \zh{炮} (pào) \\
\hline
Neutre (·) & Court, léger, sans ton propre & \zh{的} (de), \zh{么} (me), \zh{里} (li), \zh{子} (zi), \zh{上} (shang - dans shíhou), \zh{们} (men) \\
\hline
\end{tabularx}
\end{center}

\courssub{Exercice de discrimination tonale (Écoute / Répétition)}

Le professeur lit une série de mots ; l'élève identifie le ton du caractère souligné ou répète la série.

\begin{enumerate}
\item \zh{欢迎} (huānyíng) — \zh{换影} (huànyǐng, ombre changeante) — \zh{患疫} (huànyì, épidémie). \textbf{Focus : 1+2 vs 4+3 vs 4+4.}
\item \zh{节日} (jiérì) — \zh{接人} (jiē rén, accueillir qqn) — \zh{解热} (jiěrè, faire baisser la fièvre). \textbf{Focus : 2+4 vs 1+2 vs 3+4.}
\item \zh{热闹} (rènao) — \zh{惹脑} (rěnǎo, embêter le cerveau) — \zh{热恼} (rènǎo, chagrin ardent). \textbf{Focus : 4+4 vs 4+3 vs 4+3.}
\item \zh{举行} (jǔxíng) — \zh{局限} (júxiàn, limité) — \zh{巨星} (jùxīng, superstar). \textbf{Focus : 3+2 vs 2+4 vs 4+1.}
\item \zh{饺子} (jiǎozi) — \zh{教子} (jiàozǐ, éduquer son fils) — \zh{脚子} (jiǎozi, pied / argot). \textbf{Focus : 3+neutre vs 4+3 vs 3+3.}
\end{enumerate}

\courssub{Entraînement phraséologique (phrases du dialogue)}

Répète chaque phrase en exagérant légèrement la courbe des tons, puis à vitesse normale.

\begin{center}
\begin{tabularx}{\linewidth}{|X|X|}
\hline
\rowcolor{poppinkL} Phrase (caractères) & Pinyin avec tons \\
\hline
欢迎来到我们的文化节。 & Huānyíng lái dào wǒmen de wénhuà jié. \\
\hline
喀麦隆有什么传统节日？ & Kāmàilóng yǒu shénme chuántǒng jiérì? \\
\hline
比如 Ngondo 节，在海边举行。 & Bǐrú Ngondo jié, zài hǎibiān jǔxíng. \\
\hline
很热闹，有意思！ & Hěn rènao, yǒu yìsi! \\
\hline
中国最重要的节日是春节。 & Zhōngguó zuì zhòngyào de jiérì shì Chūnjié. \\
\hline
春节的时候我们吃饺子。 & Chūnjié de shíhou wǒmen chī jiǎozi. \\
\hline
我们看表演、放鞭炮。 & Wǒmen kàn biǎoyǎn, fàng biānpào. \\
\hline
\end{tabularx}
\end{center}

\begin{methode}[Auto-correction des tons en 3 étapes]
\begin{enumerate}
\item \textbf{Enregistrement} : enregistre-toi en lisant le tableau ci-dessus.
\item \textbf{Comparaison} : réécoute l'enregistrement du professeur (ou la synthèse vocale standard) et le tien. Note les écarts (tons plats au lieu de montants, 3\textsuperscript{e} ton pas assez bas, 4\textsuperscript{e} ton pas assez sec).
\item \textbf{Ciblage} : isole les 2-3 syllabes problématiques et répète-les 10 fois de suite (ex: \zh{huān-yíng}, \zh{jiǔ-xíng}, \zh{rè-nao}).
\end{enumerate}
\end{methode}

\courssub{Exercice résolu et entraînement}

\begin{exoresolu}[Complète et traduis]
Complète le mini-dialogue avec : 比如 / 在 / 欢迎 / 什么 / 的时候.\\
1. \zh{……来到喀麦隆！}\\
2. \zh{你们有……传统节日？}\\
3. \zh{我们有很多节日，…… Ngondo 节。}\\
4. \zh{这个节日……海边举行。}\\
5. \zh{春节……，我们吃饺子。}
\tcblower
\textbf{Solution :}\\
1. \zh{欢迎来到喀麦隆！} \emph{Huānyíng lái dào Kāmàilóng!} « Bienvenue au Cameroun ! »\\
2. \zh{你们有什么传统节日？} \emph{Nǐmen yǒu shénme chuántǒng jiérì?} « Quelles fêtes traditionnelles avez-vous ? »\\
3. \zh{我们有很多节日，比如 Ngondo 节。} \emph{Wǒmen yǒu hěn duō jiérì, bǐrú Ngondo jié.} « Nous avons beaucoup de fêtes, par exemple la fête Ngondo. »\\
4. \zh{这个节日在海边举行。} \emph{Zhège jiérì zài hǎibiān jǔxíng.} « Cette fête se tient au bord de la mer. »\\
5. \zh{春节的时候，我们吃饺子。} \emph{Chūnjié de shíhou, wǒmen chī jiǎozi.} « Pendant la Fête du Printemps, nous mangeons des raviolis. »
\end{exoresolu}

\begin{methode}[Mémoriser un dialogue en cinq étapes]
\begin{enumerate}
\item \textbf{Écouter} : écoute le dialogue lu par le professeur ou un camarade, livre fermé, pour saisir le sens global.
\item \textbf{Répéter phrase par phrase} : répète chaque réplique après le modèle, en soignant les tons (surtout huānyíng, 1\textsuperscript{er} + 2\textsuperscript{e} ton).
\item \textbf{Jouer le rôle A} : lis le rôle de l'élève camerounaise pendant qu'un camarade joue B.
\item \textbf{Jouer le rôle B} : inversez les rôles pour mémoriser les deux moitiés du dialogue.
\item \textbf{Jouer sans le texte} : reproduis le dialogue de mémoire, en regardant seulement le schéma des cinq étapes (saluer → question → présenter → réagir → relancer).
\end{enumerate}
\end{methode}

\begin{exercice}[À toi de jouer]
Par groupes de deux : récitez le dialogue modèle par cœur, puis inversez les rôles. Ensuite, remplacez le Ngondo par une autre fête de ta région (ex: \zh{薰衣草节} à Buea, \zh{Ngondo} à Douala, \zh{Fête du Maïs} à l'Ouest) et la Fête du Printemps par une autre fête chinoise (\zh{中秋节}, \zh{端午节} Duānwǔjié - Fête des Bateaux-Dragons), en gardant les mêmes structures : 欢迎来到…… / 有什么……？ / 比如…… / 在……举行 / ……的时候你们做什么？
\end{exercice}

\begin{corrige}[Exercice — modèle de production attendue]
A : \zh{你好！欢迎来到我们的文化节。} B : \zh{谢谢！喀麦隆有什么传统节日？} A : \zh{有很多，比如圣诞节，在教堂和家里庆祝，很热闹。} \emph{Yǒu hěn duō, bǐrú Shèngdànjié, zài jiàotáng hé jiā lǐ qìngzhù, hěn rènao.} « Il y en a beaucoup, par exemple Noël, célébré à l'église et à la maison, très animé. » B : \zh{有意思！中国最重要的节日是中秋节。} A : \zh{中秋节的时候你们做什么？} B : \zh{我们吃月饼、看月亮。} \emph{Wǒmen chī yuèbǐng, kàn yuèliang.} « Nous mangeons des gâteaux de lune et admirons la lune. »
\end{corrige}

\begin{savaistu}[Le Ngondo et la Fête du Printemps]
Le \cle{Ngondo} est la grande fête traditionnelle du peuple \cle{Sawa} (littoral camerounais) : célébrée au bord de l'eau, notamment à Douala, elle comprend des courses de pirogues, des danses et des rituels de communion avec les ancêtres. La \cle{Fête du Printemps} (春节, Chūnjié), ou Nouvel An lunaire, est en Chine le moment des retrouvailles familiales : on partage les raviolis (饺子, jiǎozi), symbole de richesse, on regarde les spectacles télévisés et autrefois on faisait éclater des pétards (鞭炮, biānpào) pour chasser les mauvais esprits. Des deux côtés, la fête réunit la famille, le peuple et l'eau ou le ciel : un beau sujet de conversation entre amis camerounais et chinois !
\end{savaistu}

\coursec{Formules fonctionnelles : présenter, opiner, relancer}

Dans la section précédente, nous avons appris à présenter une fête grâce au dialogue modèle 1 : \zh{我来介绍一下春节。} (\emph{Wǒ lái jièshào yíxià Chūnjié.} « Je vais vous présenter la fête du Printemps. »). Vous savez maintenant \emph{quoi} dire pour ouvrir un exposé. Mais un bon orateur ne se contente pas de présenter : il donne son avis, il réagit aux autres, il pose des questions pour faire vivre la discussion. C'est exactement ce que nous allons construire ici : votre \cle{boîte à outils du débatteur}, en trois familles de formules — \zh{介绍} (présenter), \zh{说看法} (donner son avis), \zh{提问} (relancer).

\courssub{Observer d'abord : un mini-débat en classe}

Lisons d'abord ce court échange entre deux élèves d'un lycée de Yaoundé. Soulignez mentalement les formules qui reviennent.

\begin{textebox}[Mini-débat : la musique au Cameroun et en Chine]
A：我来介绍一下喀麦隆的音乐。我觉得摩科萨很有意思。\\
\emph{Wǒ lái jièshào yíxià Kāmàilóng de yīnyuè. Wǒ juéde Mókēsà hěn yǒuyìsi.}\\
« Je vais vous présenter la musique du Cameroun. Je trouve que le makossa est très intéressant. »\\[4pt]
B：你说得对，不过我觉得京剧也很好听。\\
\emph{Nǐ shuō de duì, búguò wǒ juéde Jīngjù yě hěn hǎotīng.}\\
« Tu as raison, cependant je trouve que l'opéra de Pékin est aussi très beau à écouter. »\\[4pt]
A：我不太同意。对我来说，摩科萨比京剧有意思。你呢？\\
\emph{Wǒ bú tài tóngyì. Duì wǒ lái shuō, Mókēsà bǐ Jīngjù yǒu yìsi. Nǐ ne?}\\
« Je ne suis pas tout à fait d'accord. Pour moi, le makossa est plus intéressant que l'opéra de Pékin. Et toi ? »\\[4pt]
B：我有不同的看法：京剧比摩科萨有名。你能举个例子吗？\\
\emph{Wǒ yǒu bùtóng de kànfǎ: Jīngjù bǐ Mókēsà yǒumíng. Nǐ néng jǔ ge lìzi ma?}\\
« J'ai un avis différent : l'opéra de Pékin est plus célèbre que le makossa. Peux-tu donner un exemple ? »
\end{textebox}

Observez : chaque prise de parole contient \textbf{une formule d'opinion} (\zh{我觉得}, \zh{对我来说}) et souvent \textbf{une question finale} (\zh{你呢？}, \zh{你能举个例子吗？}). C'est le moteur du débat. Décortiquons maintenant chaque famille de formules.

\courssub{Présenter : ouvrir son exposé}

\begin{propriete}[Trois formules pour présenter]
\begin{itemize}
\item \zh{我来介绍一下……} (\emph{Wǒ lái jièshào yíxià...}) « Je vais (vous) présenter... » — formule d'ouverture la plus courante. \zh{一下} adoucit le verbe, comme un « un peu » de politesse.
\item \zh{这是……} (\emph{Zhè shì...}) « Voici... / Ceci est... » — pour montrer un objet, une photo, une affiche.
\item \zh{我要说的是……} (\emph{Wǒ yào shuō de shì...}) « Ce dont je vais parler, c'est... » — pour annoncer le sujet précis.
\end{itemize}
Structure : \textbf{formule + nom du sujet}. Exemple : \zh{我要说的是中国的春节。} (\emph{Wǒ yào shuō de shì Zhōngguó de Chūnjié.}) « Ce dont je vais parler, c'est la fête du Printemps chinoise. »
\end{propriete}

\begin{exemplebox}[Présenter en contexte camerounais]
\zh{我来介绍一下杜阿拉的文化节。} \emph{Wǒ lái jièshào yíxià Dù'ālā de wénhuà jié.} « Je vais vous présenter le festival culturel de Douala. »\\
\zh{这是一种喀麦隆的传统舞蹈。} \emph{Zhè shì yì zhǒng Kāmàilóng de chuántǒng wǔdǎo.} « Voici une sorte de danse traditionnelle camerounaise. »
\end{exemplebox}

\courssub{Donner son avis : 我觉得, 我认为, 对我来说}

\begin{propriete}[La règle d'or : pas de « que » en chinois]
En français on dit « je trouve \textbf{que} la cuisine chinoise est bonne ». En chinois, on \cle{juxtapose} simplement les deux phrases, sans mot de liaison :\\
\textbf{Sujet + 觉得 / 认为 + phrase complète (Sujet + Verbe)}\\
\zh{我觉得中国的菜很好吃。} \emph{Wǒ juéde Zhōngguó de cài hěn hǎochī.} « Je trouve que la cuisine chinoise est très bonne. »\\
Littéralement : « Moi + trouver + les plats chinois très bons ». Le « que » français n'a \textbf{aucun équivalent} : ne cherchez pas à le traduire.
\end{propriete}

Trois outils d'opinion, du plus courant au plus formel :

\begin{itemize}
\item \zh{我觉得……} (\emph{wǒ juéde}) : « je trouve que... », avis personnel, oral, le plus fréquent.
\item \zh{我认为……} (\emph{wǒ rènwéi}) : « je pense que... / j'estime que... », plus soutenu, idéal pour un exposé.
\item \zh{对我来说，……} (\emph{duì wǒ lái shuō}) : « pour moi,... / en ce qui me concerne,... » — la virgule est obligatoire après \zh{说}.
\end{itemize}

\begin{exemplebox}[Nuancer : 但是, 不过, 虽然……但是……]
\zh{京剧很有名，但是我不太喜欢。} \emph{Jīngjù hěn yǒumíng, dànshì wǒ bú tài xǐhuan.} « L'opéra de Pékin est très célèbre, mais je ne l'aime pas beaucoup. »\\
\zh{春节很热闹，不过我更喜欢圣诞节。} \emph{Chūnjié hěn rènao, búguò wǒ gèng xǐhuan Shèngdànjié.} « La fête du Printemps est très animée, cependant je préfère Noël. »\\
\zh{虽然摩科萨很有意思，但是京剧也很有名。} \emph{Suīrán Mókēsà hěn yǒuyìsi, dànshì Jīngjù yě hěn yǒumíng.} « Bien que le makossa soit intéressant, l'opéra de Pékin est aussi célèbre. »\\
\textbf{Attention, différence majeure avec le français :} en chinois, \zh{虽然……但是……} s'emploient \textbf{ensemble} dans la même phrase (« bien que... mais... »), alors qu'en français « bien que » et « mais » sont incompatibles. C'est normal et correct en chinois.
\end{exemplebox}

\courssub{Accord et désaccord : réagir à l'avis d'autrui}

\begin{propriete}[Quatre réactions à connaître par cœur]
\begin{itemize}
\item \zh{我同意。} (\emph{Wǒ tóngyì.}) « Je suis d'accord. »
\item \zh{你说得对。} (\emph{Nǐ shuō de duì.}) « Tu as raison / ce que tu dis est juste. »
\item \zh{我不太同意。} (\emph{Wǒ bú tài tóngyì.}) « Je ne suis pas tout à fait d'accord. » — désaccord poli.
\item \zh{我有不同的看法。} (\emph{Wǒ yǒu bùtóng de kànfǎ.}) « J'ai un avis différent. » — désaccord formel.
\end{itemize}
\end{propriete}

\begin{attention}[Ne dites jamais 我是同意 !]
Erreur typique des francophones : traduire « je suis d'accord » mot à mot par \zh{我是同意}. En chinois, \zh{同意} (\emph{tóngyì}) est un \textbf{verbe} à part entière (« être d'accord »), pas un adjectif. On dit donc simplement \zh{我同意}, sans \zh{是}. De même : \zh{我不同意} « je ne suis pas d'accord ». Retenez : \textbf{是 ne se met jamais devant un verbe}.
\end{attention}

\begin{exemplebox}[Exemple complet traduit]
\zh{我不太同意，我觉得摩科萨比京剧有意思。}\\
\emph{Wǒ bú tài tóngyì, wǒ juéde Mókēsà bǐ Jīngjù yǒu yìsi.}\\
« Je ne suis pas tout à fait d'accord, je trouve le makossa plus intéressant que l'opéra de Pékin. »\\
Remarquez l'enchaînement : désaccord poli + avis personnel + comparaison. C'est le schéma type d'une bonne réplique de débat.
\end{exemplebox}

\courssub{Comparer : A 比 B + adjectif et A 跟 B 一样}

\begin{propriete}[La comparaison avec 比]
Structure : \textbf{A + 比 + B + adjectif} = « A est plus... que B ».\\
\zh{京剧比摩科萨有名。} \emph{Jīngjù bǐ Mókēsà yǒumíng.} « L'opéra de Pékin est plus célèbre que le makossa. »\\
\zh{杜阿拉比雅温得热。} \emph{Dù'ālā bǐ Yǎwēndé rè.} « Douala est plus chaude que Yaoundé. »\\
Pour l'égalité : \textbf{A + 跟 + B + 一样 (+ adjectif)} = « A est comme B ».\\
\zh{春节跟圣诞节一样热闹。} \emph{Chūnjié gēn Shèngdànjié yíyàng rènao.} « La fête du Printemps est aussi animée que Noël. »\\
\textbf{Piège :} avec \zh{比}, l'adjectif seul suffit ; ne mettez pas \zh{很} (\zh{京剧比摩科萨很有名} est fautif). Pour renforcer, utilisez \zh{更} ou \zh{……得多} : \zh{京剧比摩科萨有名得多} « bien plus célèbre ».
\end{propriete}

\begin{popfigure}
\begin{center}
\begin{tikzpicture}[node distance=0.9cm and 1.1cm,
box/.style={draw=popblue, fill=popblueL, rounded corners, thick, align=center, minimum height=0.9cm, inner sep=5pt},
adj/.style={draw=poporange, fill=poporangeL, rounded corners, thick, align=center, inner sep=5pt}]
\node[box, align=center] (a){A\\京剧\\ \emph{Jīngjù}};
\node[box, right=of a, fill=poppinkL, draw=poppink, align=center] (bi){比\\ \emph{bǐ}\\ « que »};
\node[box, right=of bi, align=center] (b){B\\摩科萨\\ \emph{Mókēsà}};
\node[adj, right=of b, align=center] (adj){adjectif\\有名\\ \emph{yǒumíng}};
\draw[-{Stealth}, very thick, popdark] (a) -- (bi);
\draw[-{Stealth}, very thick, popdark] (bi) -- (b);
\draw[-{Stealth}, very thick, popdark] (b) -- (adj);
\end{tikzpicture}
\end{center}
\legende{Schéma de la comparaison : 京剧 比 摩科萨 有名。« L'opéra de Pékin est plus célèbre que le makossa. »}
\end{popfigure}

\courssub{Relancer : redonner la parole}

Un débat n'est pas un monologue. Après avoir donné votre avis, vous devez \cle{relancer} votre partenaire par une question courte :

\begin{itemize}
\item \zh{你呢？} (\emph{Nǐ ne?}) « Et toi ? » — la relance la plus simple et la plus utile.
\item \zh{你觉得怎么样？} (\emph{Nǐ juéde zěnmeyàng?}) « Qu'en penses-tu ? »
\item \zh{为什么？} (\emph{Wèishénme?}) « Pourquoi ? » — pour demander une justification.
\item \zh{你能举个例子吗？} (\emph{Nǐ néng jǔ ge lìzi ma?}) « Peux-tu donner un exemple ? » — pour demander une preuve concrète.
\end{itemize}

\begin{methode}[Relancer un débat en trois étapes]
\begin{enumerate}
\item \textbf{Donnez votre avis} avec une formule d'opinion : \zh{我觉得……} ou \zh{我认为……}.
\item \textbf{Justifiez ou comparez} : ajoutez une raison (\zh{因为……} « parce que ») ou une comparaison avec \zh{比}.
\item \textbf{Terminez toujours par une question} : \zh{你呢？} ou \zh{你觉得怎么样？} — c'est la question qui redonne la parole et fait vivre le débat. Sans elle, l'échange meurt.
\end{enumerate}
Exemple complet : \zh{我觉得喀麦隆的足球比中国的有名，因为喀麦隆有很多好球员。你呢？} \emph{Wǒ juéde Kāmàilóng de zúqiú bǐ Zhōngguó de yǒumíng, yīnwèi Kāmàilóng yǒu hěn duō hǎo qiúyuán. Nǐ ne?} « Je trouve que le football camerounais est plus célèbre que le chinois, parce que le Cameroun a beaucoup de bons joueurs. Et toi ? »
\end{methode}

\courssub{Tableau récapitulatif de la boîte à outils}

\begin{center}
\begin{tabularx}{\linewidth}{|X|X|X|X|}
\hline
\rowcolor{popblueL} Caractères & Pinyin & Fonction & Traduction \\
\hline
我来介绍一下…… & wǒ lái jièshào yíxià & présenter & je vais vous présenter... \\
\hline
这是…… & zhè shì & présenter & voici... \\
\hline
我要说的是…… & wǒ yào shuō de shì & présenter & ce dont je vais parler, c'est... \\
\hline
我觉得…… & wǒ juéde & opiner & je trouve que... \\
\hline
我认为…… & wǒ rènwéi & opiner (soutenu) & je pense que... \\
\hline
对我来说，…… & duì wǒ lái shuō & opiner & pour moi,... \\
\hline
但是 / 不过 & dànshì / búguò & nuancer & mais / cependant \\
\hline
虽然……但是…… & suīrán... dànshì... & nuancer & bien que... (mais)... \\
\hline
我同意 & wǒ tóngyì & accord & je suis d'accord \\
\hline
你说得对 & nǐ shuō de duì & accord & tu as raison \\
\hline
我不太同意 & wǒ bú tài tóngyì & désaccord poli & je ne suis pas tout à fait d'accord \\
\hline
我有不同的看法 & wǒ yǒu bùtóng de kànfǎ & désaccord & j'ai un avis différent \\
\hline
你呢？ & nǐ ne? & relancer & et toi ? \\
\hline
你觉得怎么样？ & nǐ juéde zěnmeyàng? & relancer & qu'en penses-tu ? \\
\hline
你能举个例子吗？ & nǐ néng jǔ ge lìzi ma? & relancer & peux-tu donner un exemple ? \\
\hline
A 比 B + adj. & A bǐ B & comparer & A est plus... que B \\
\hline
A 跟 B 一样 & A gēn B yíyàng & comparer & A est comme B \\
\hline
\end{tabularx}
\end{center}

\begin{popfigure}
\begin{center}
\begin{tikzpicture}[mindmap, grow cyclic,
every node/.style={concept, align=center, font=\small},
concept color=popblue,
level 1/.append style={level distance=3.4cm, sibling angle=120},
level 2/.append style={level distance=2.4cm, sibling angle=55}]
\node[concept, fill=popblue, text=white, align=center]{辩论工具箱\\ Boîte à outils\\ du débatteur}
child[concept color=popgreenL] { node[concept, align=center]{介绍\\ présenter}
  child { node[concept, font=\scriptsize, align=center]{我来介绍\\ 一下……} }
  child { node[concept, font=\scriptsize]{这是……} }
}
child[concept color=poporangeL] { node[concept, align=center]{说看法\\ donner son avis}
  child { node[concept, font=\scriptsize]{我觉得……} }
  child { node[concept, font=\scriptsize, align=center]{对我来说\\ ……} }
}
child[concept color=poppinkL] { node[concept, align=center]{提问\\ relancer}
  child { node[concept, font=\scriptsize]{你呢？} }
  child { node[concept, font=\scriptsize, align=center]{你觉得\\ 怎么样？} }
};
\end{tikzpicture}
\end{center}
\legende{Carte mentale : la boîte à outils du débatteur — trois branches, deux formules chacune.}
\end{popfigure}

\courssub{Entraînement : du français au chinois, à voix haute}

\begin{exoresolu}[Traduire et dire à voix haute]
Énoncé. Traduisez en chinois, puis prononcez à voix haute en soignant les tons : « Je trouve que la fête du Printemps est plus animée que Noël. Mais je ne suis pas tout à fait d'accord avec toi : pour moi, Noël est aussi intéressant. Et toi, qu'en penses-tu ? »
\tcblower
Solution détaillée. On découpe la phrase française en blocs chinois :
\begin{itemize}
\item « Je trouve que la fête du Printemps est plus animée que Noël » → \zh{我觉得春节比圣诞节热闹。} \emph{Wǒ juéde Chūnjié bǐ Shèngdànjié rènao.} (opinion \zh{觉得} + comparaison \zh{比}, sans « que », sans \zh{很}).
\item « Mais je ne suis pas tout à fait d'accord avec toi » → \zh{但是我也不太同意。} \emph{Dànshì wǒ yě bú tài tóngyì.} (\zh{也} « aussi » ; pas de \zh{是} devant \zh{同意}).
\item « Pour moi, Noël est aussi intéressant » → \zh{对我来说，圣诞节也很有意思。} \emph{Duì wǒ lái shuō, Shèngdànjié yě hěn yǒuyìsi.}
\item « Et toi, qu'en penses-tu ? » → \zh{你呢？你觉得怎么样？} \emph{Nǐ ne? Nǐ juéde zěnmeyàng?}
\end{itemize}
Phrase complète à dire d'un souffle : \zh{我觉得春节比圣诞节热闹。但是我也不太同意。对我来说，圣诞节也很有意思。你呢？你觉得怎么样？}
\end{exoresolu}

\begin{exercice}[À vous de débattre]
Traduisez à voix haute, puis jouez la scène à deux : 1. « Je vais vous présenter la musique camerounaise. » 2. « Tu as raison, cependant je pense que la musique chinoise est aussi très belle. » 3. « J'ai un avis différent. Peux-tu donner un exemple ? »
\end{exercice}

\begin{corrige}[Exercice : À vous de débattre]
1. \zh{我来介绍一下喀麦隆的音乐。} \emph{Wǒ lái jièshào yíxià Kāmàilóng de yīnyuè.}\\
2. \zh{你说得对，不过我认为中国的音乐也很好听。} \emph{Nǐ shuō de duì, búguò wǒ rènwéi Zhōngguó de yīnyuè yě hěn hǎotīng.}\\
3. \zh{我有不同的看法。你能举个例子吗？} \emph{Wǒ yǒu bùtóng de kànfǎ. Nǐ néng jǔ ge lìzi ma?}
\end{corrige}

\begin{savaistu}[Dire « non » avec délicatesse]
En Chine, dire directement « non » ou « je ne suis pas d'accord » (\zh{我不同意}) peut sembler brusque, voire impoli. On préfère des formules adoucies comme \zh{我不太同意} « je ne suis pas tout à fait d'accord » ou \zh{可能吧} (\emph{kěnéng ba}) « peut-être ». Cette manière de ménager l'interlocuteur et de préserver l'harmonie du groupe n'est pas étrangère au Cameroun : dans beaucoup de cultures camerounaises, on évite aussi le refus frontal devant un aîné ou un invité, et l'on préfère les réponses nuancées. Le débat poli est donc un terrain où le Cameroun et la Chine se comprennent très bien.
\end{savaistu}

\begin{aretenir}[L'essentiel de la section]
\begin{itemize}
\item Présenter : \zh{我来介绍一下……} / \zh{这是……} / \zh{我要说的是……}.
\item Opiner : \zh{我觉得} et \zh{我认为} + phrase complète, \textbf{sans} traduire le « que » français.
\item Réagir : \zh{我同意} (jamais \zh{我是同意} !), \zh{你说得对}, \zh{我不太同意}, \zh{我有不同的看法}.
\item Comparer : A \zh{比} B + adjectif (sans \zh{很}) ; A \zh{跟} B \zh{一样}.
\item Relancer : terminez toujours par \zh{你呢？} ou \zh{你觉得怎么样？}.
\end{itemize}
\end{aretenir}

\coursec{Les tons : entraînement ciblé}

Dans la section précédente, nous avons appris à présenter une fête, à donner notre avis et à relancer la parole avec des formules comme \zh{我认为} (\emph{wǒ rènwéi}, « je pense que ») ou \zh{你觉得呢？} (\emph{nǐ juéde ne?}, « et toi, qu'en penses-tu ? »). Mais une belle phrase, même parfaitement construite, devient incompréhensible si les \cle{tons} sont faux. Avant de passer au débat, nous devons donc entraîner l'oreille et la bouche : c'est l'objectif de cette section, véritable « gymnastique phonétique » indispensable à tout exposé oral réussi.

\courssub{Rappel : les quatre tons et le ton neutre}

\begin{definition}[Qu'est-ce qu'un ton ?]
En chinois mandarin, chaque syllabe est prononcée avec une \cle{mélodie} (un contour de hauteur de voix) qui fait partie du mot, au même titre que ses consonnes et ses voyelles. Changer le ton, c'est changer de mot. Le mandarin possède \cle{quatre tons} et un \cle{ton neutre} (léger et bref, sans marque).
\end{definition}

\begin{center}
\begin{tabularx}{\linewidth}{|X|X|X|X|}
\hline
\rowcolor{popblueL} Ton & Contour mélodique & Marque & Exemple \\
\hline
1\textsuperscript{er} ton & haut et plat (—) & ā & \zh{妈} mā « maman » \\
\hline
2\textsuperscript{e} ton & montant (/) & á & \zh{麻} má « chanvre » \\
\hline
3\textsuperscript{e} ton & descendant-montant (V) & ǎ & \zh{马} mǎ « cheval » \\
\hline
4\textsuperscript{e} ton & descendant, sec (\textbackslash) & à & \zh{骂} mà « gronder » \\
\hline
ton neutre & bref et léger & (aucune) & \zh{吗} ma (particule de question) \\
\hline
\end{tabularx}
\end{center}

\begin{popfigure}
\begin{tikzpicture}[scale=1.0]
% Ton 1
\draw[popline, dashed] (0,0) -- (2.6,0);
\draw[popline, dashed] (0,1.6) -- (2.6,1.6);
\draw[line width=2pt, popblue] (0.2,1.4) -- (2.4,1.4);
\node[align=center] at (1.3,-0.7) {\zh{妈} mā\\ 1\textsuperscript{er} ton};
% Ton 2
\begin{scope}[xshift=3.4cm]
\draw[popline, dashed] (0,0) -- (2.6,0);
\draw[popline, dashed] (0,1.6) -- (2.6,1.6);
\draw[line width=2pt, poporange] (0.2,0.3) -- (2.4,1.4);
\node[align=center] at (1.3,-0.7) {\zh{麻} má\\ 2\textsuperscript{e} ton};
\end{scope}
% Ton 3
\begin{scope}[xshift=6.8cm]
\draw[popline, dashed] (0,0) -- (2.6,0);
\draw[popline, dashed] (0,1.6) -- (2.6,1.6);
\draw[line width=2pt, popgreen] (0.2,1.1) -- (1.3,0.2) -- (2.4,1.3);
\node[align=center] at (1.3,-0.7) {\zh{马} mǎ\\ 3\textsuperscript{e} ton};
\end{scope}
% Ton 4
\begin{scope}[xshift=10.2cm]
\draw[popline, dashed] (0,0) -- (2.6,0);
\draw[popline, dashed] (0,1.6) -- (2.6,1.6);
\draw[line width=2pt, poppink] (0.2,1.4) -- (2.4,0.2);
\node[align=center] at (1.3,-0.7) {\zh{骂} mà\\ 4\textsuperscript{e} ton};
\end{scope}
\end{tikzpicture}
\legende{Les contours mélodiques des quatre tons, avec un mot exemple sous chaque courbe.}
\end{popfigure}

\begin{attention}[Le français n'a pas de tons !]
En français, la mélodie de la voix exprime l'émotion ou la question (« Tu viens ? » monte à la fin), mais elle ne change jamais le sens du mot. En chinois, c'est l'inverse : un mauvais ton produit un \emph{autre mot}. Erreur classique et dangereuse : \zh{买} (\emph{mǎi}, 3\textsuperscript{e} ton, « acheter ») et \zh{卖} (\emph{mài}, 4\textsuperscript{e} ton, « vendre »). Au marché de Mokolo, dire \zh{我想卖芒果} (\emph{Wǒ xiǎng mài mángguǒ}) alors qu'on veut \emph{acheter} des mangues peut créer un vrai malentendu ! Autre paire célèbre : \zh{问} (\emph{wèn}, « demander ») et \zh{吻} (\emph{wěn}, « embrasser »).
\end{attention}

\courssub{Paires minimales : entraîner l'oreille}

Une \cle{paire minimale} est un couple de mots qui ne diffèrent que par un seul trait — ici, uniquement le ton. La série \emph{ma} est le grand classique :

\begin{center}
\begin{tabularx}{\linewidth}{|X|X|X|X|}
\hline
\rowcolor{popblueL} Caractère & Pinyin & Ton & Traduction \\
\hline
妈 & mā & 1\textsuperscript{er} & maman \\
\hline
麻 & má & 2\textsuperscript{e} & chanvre \\
\hline
马 & mǎ & 3\textsuperscript{e} & cheval \\
\hline
骂 & mà & 4\textsuperscript{e} & gronder, insulter \\
\hline
吗 & ma & neutre & particule interrogative \\
\hline
\end{tabularx}
\end{center}

La comptine traditionnelle chinoise résume tout : \zh{妈妈骑马，马慢，妈妈骂马。} (\emph{Māma qí mǎ, mǎ màn, māma mà mǎ.}) « Maman monte à cheval, le cheval est lent, maman gronde le cheval. » Répétez-la en exagérant chaque contour mélodique !

Appliquons maintenant au lexique culturel de notre leçon :

\begin{center}
\begin{tabularx}{\linewidth}{|X|X|X|X|}
\hline
\rowcolor{popblueL} Mot & Pinyin & Piège de ton & Traduction \\
\hline
音乐 & yīnyuè & yīn = 1\textsuperscript{er} ton plat, pas montant & musique \\
\hline
唱 & chàng & 4\textsuperscript{e} ton, sec et descendant & chanter \\
\hline
长 & cháng & 2\textsuperscript{e} ton, montant & long \\
\hline
画 & huà & 4\textsuperscript{e} ton & peindre ; un tableau \\
\hline
舞 & wǔ & 3\textsuperscript{e} ton & danser (dans 跳舞 tiàowǔ) \\
\hline
戏 & xì & 4\textsuperscript{e} ton & théâtre, opéra (dans 京剧 Jīngjù) \\
\hline
\end{tabularx}
\end{center}

Confondre \zh{唱} (\emph{chàng}, chanter) et \zh{长} (\emph{cháng}, long) dans la phrase \zh{我喜欢唱歌} (\emph{Wǒ xǐhuan chànggē}, « j'aime chanter ») rendrait votre exposé sur la musique incompréhensible. De même, \zh{音乐} commence par un 1\textsuperscript{er} ton bien haut et plat : ne le laissez pas « tomber » à la française.

\begin{exemplebox}[Phrases à répéter à voix haute]
\zh{我喜欢听音乐。} \emph{Wǒ xǐhuan tīng yīnyuè.} « J'aime écouter de la musique. »\\
\zh{她会唱喀麦隆的歌。} \emph{Tā huì chàng Kāmàilóng de gē.} « Elle sait chanter des chansons camerounaises. »\\
\zh{这个节目很长。} \emph{Zhège jiémù hěn cháng.} « Ce spectacle est très long. »
\end{exemplebox}

\courssub{Le sandhi du troisième ton}

Observez d'abord ces deux phrases que vous connaissez par cœur : \zh{你好} et \zh{我很好}. Vous les avez apprises avec les tons \emph{nǐ hǎo} et \emph{wǒ hěn hǎo}. Pourtant, à l'oral, personne ne prononce ainsi. Pourquoi ? Parce que deux 3\textsuperscript{e} tons consécutifs sont presque impossibles à enchaîner : la voix devrait descendre, remonter, redescendre, remonter — un vrai slalom !

\begin{propriete}[Règle du sandhi du 3\textsuperscript{e} ton]
Quand deux syllabes au 3\textsuperscript{e} ton se suivent, la \cle{première} se prononce comme un 2\textsuperscript{e} ton (montant). La seconde garde son 3\textsuperscript{e} ton.\\
\textbf{Schéma :} 3\textsuperscript{e} + 3\textsuperscript{e} → 2\textsuperscript{e} + 3\textsuperscript{e}\\
Ainsi \zh{你好} s'écrit \emph{nǐ hǎo} mais se prononce \emph{ní hǎo}. On écrit toujours le ton d'origine dans le pinyin ; seule la prononciation change.\\
\textbf{Cas de trois 3\textsuperscript{e} tons :} on applique la règle de proche en proche, en suivant le sens. Dans \zh{我很好} (\emph{wǒ hěn hǎo}), les mots se groupent \emph{wǒ} + \emph{hěnhǎo} : \emph{hěn} devient \emph{hén} devant \emph{hǎo}, puis \emph{wǒ} devient \emph{wó} devant cette syllabe devenue montante — d'où la prononciation \emph{wó hén hǎo}.
\end{propriete}

\begin{exemplebox}[Application de la règle]
\zh{我很好。} \emph{Wǒ hěn hǎo} → prononcé \emph{wó hén hǎo}. « Je vais très bien. » (deux sandhis successifs !)\\
\zh{你好！} \emph{Nǐ hǎo} → \emph{ní hǎo}. « Bonjour ! »\\
\zh{我也可以。} \emph{Wǒ yě kěyǐ} → \emph{wó yě kéyǐ}. « Moi aussi, je peux. » (\emph{kěyǐ} : \emph{kě} devient \emph{ké} devant \emph{yǐ}.)\\
\zh{请你给我。} \emph{Qǐng nǐ gěi wǒ} → \emph{qíng nǐ géi wǒ}. « Donne-le-moi, s'il te plaît. » (\emph{qǐng} change devant \emph{nǐ} ; \emph{gěi} change devant \emph{wǒ}.)
\end{exemplebox}

\begin{attention}[Écriture et prononciation]
Dans vos rédactions, écrivez toujours le pinyin avec les tons \emph{d'origine} : \emph{nǐ hǎo}, jamais \emph{ní hǎo}. Le sandhi est une règle de \cle{prononciation}, pas d'orthographe. À l'examen, écrire \emph{ní hǎo} est une faute.
\end{attention}

\courssub{Le sandhi de 不}

Le mot \zh{不} (\emph{bù}, « ne... pas »), au 4\textsuperscript{e} ton, subit lui aussi une transformation.

\begin{propriete}[Règle de 不]
\zh{不} se prononce \emph{bù} (4\textsuperscript{e} ton) partout, \textbf{sauf devant une syllabe au 4\textsuperscript{e} ton}, où il devient \emph{bú} (2\textsuperscript{e} ton montant). Deux 4\textsuperscript{e} tons secs qui se suivent sont trop « durs » ; la langue adoucit le premier en le faisant monter.\\
\textbf{Schéma :} bù + 4\textsuperscript{e} ton → bú + 4\textsuperscript{e} ton
\end{propriete}

\begin{center}
\begin{tabularx}{\linewidth}{|X|X|X|}
\hline
\rowcolor{popblueL} Mot & Prononciation réelle & Traduction \\
\hline
不是 & bú shì (et non bù shì) & ne pas être \\
\hline
不对 & bú duì & ne pas être correct \\
\hline
不要 & bú yào & ne pas vouloir ; il ne faut pas \\
\hline
不去 & bú qù & ne pas aller \\
\hline
不好 & bù hǎo (hǎo = 3\textsuperscript{e} ton, pas de changement) & pas bon \\
\hline
不来 & bù lái (lái = 2\textsuperscript{e} ton, pas de changement) & ne pas venir \\
\hline
\end{tabularx}
\end{center}

Phrase utile pour le débat : \zh{我不是很喜欢京剧。} \emph{Wǒ bú shì hěn xǐhuan Jīngjù.} « Je n'aime pas beaucoup l'opéra de Pékin. » — remarquez le \emph{bú shì} et le sandhi \emph{hěn xǐhuan} prononcé \emph{hén xǐhuan} !

\courssub{Méthode : la main qui dessine les tons}

\begin{methode}[Le geste mélodique]
Pour ancrer chaque ton dans le corps, dessinez son contour avec la main pendant que vous prononcez :
\begin{enumerate}
\item \textbf{1\textsuperscript{er} ton} : main plate qui avance horizontalement, voix haute et stable — \emph{mā}.
\item \textbf{2\textsuperscript{e} ton} : main qui monte en diagonale, comme quand on pose une question surprise — \emph{má}.
\item \textbf{3\textsuperscript{e} ton} : main qui descend puis remonte (un V), voix qui creuse — \emph{mǎ}.
\item \textbf{4\textsuperscript{e} ton} : main qui tombe d'un coup sec, comme un « non ! » catégorique — \emph{mà}.
\item Répétez chaque mot du vocabulaire culturel (\emph{yīnyuè, chàng, wǔ, huà, xì}) avec le geste, d'abord lentement, puis à vitesse normale.
\end{enumerate}
\end{methode}

\courssub{Exercices de discrimination auditive}

\begin{experience}[Le jeu des cartes de tons]
Chaque élève prépare cinq cartes : 1, 2, 3, 4 et « neutre ». L'enseignant lit un mot du lexique culturel ; les élèves lèvent la carte du ton entendu (pour les mots de deux syllabes, on se concentre sur la syllabe demandée). Série proposée : \emph{mā / cháng / mǎi / chàng / huà / má / wǔ / mài / xì / mà}. Variante en binômes : un élève lit, l'autre montre la carte, puis on inverse. Celui qui se trompe doit prononcer correctement trois fois le mot avec le geste de la main.
\end{experience}

\begin{exoresolu}[Quel ton entendez-vous ?]
Énoncé : l'enseignant prononce les quatre syllabes suivantes : (a) \emph{mài} ; (b) \emph{cháng} ; (c) \emph{wǔ} ; (d) \emph{mā}. Pour chacune, indiquez le ton et donnez un mot du lexique ou de la vie quotidienne qui correspond.
\tcblower
Solution détaillée : (a) \emph{mài} : 4\textsuperscript{e} ton, descendant et sec → \zh{卖}, « vendre » (attention à ne pas confondre avec \zh{买} \emph{mǎi}, « acheter », 3\textsuperscript{e} ton). (b) \emph{cháng} : 2\textsuperscript{e} ton, montant → \zh{长}, « long » (à distinguer de \zh{唱} \emph{chàng}, « chanter »). (c) \emph{wǔ} : 3\textsuperscript{e} ton, descendant-montant → \zh{舞}, « danser », comme dans \zh{跳舞} (\emph{tiàowǔ}). (d) \emph{mā} : 1\textsuperscript{er} ton, haut et plat → \zh{妈}, « maman ». Astuce : le 4\textsuperscript{e} ton est le plus facile à reconnaître (il « claque ») ; le plus difficile pour les francophones est de distinguer le 2\textsuperscript{e} ton (montant) du 3\textsuperscript{e} ton (qui descend d'abord).
\end{exoresolu}

\begin{exercice}[À vous de chanter les tons !]
1. Classez ces mots par ton : \zh{歌} (gē), \zh{节} (jié), \zh{舞} (wǔ), \zh{唱} (chàng), \zh{画} (huà), \zh{音} (yīn).\\
2. Indiquez la prononciation réelle (avec sandhi) de : \zh{很好}, \zh{不是}, \zh{我想}, \zh{不要}, \zh{你可以}.\\
3. Apprenez par cœur la comptine \zh{妈妈骑马} et récitez-la à votre voisin avec les gestes de la main.
\end{exercice}

\begin{corrige}[Exercice « À vous de chanter les tons ! »]
1. 1\textsuperscript{er} ton : \zh{歌} gē, \zh{音} yīn ; 2\textsuperscript{e} ton : \zh{节} jié ; 3\textsuperscript{e} ton : \zh{舞} wǔ ; 4\textsuperscript{e} ton : \zh{唱} chàng, \zh{画} huà.\\
2. \zh{很好} → \emph{hén hǎo} ; \zh{不是} → \emph{bú shì} ; \zh{我想} → \emph{wó xiǎng} ; \zh{不要} → \emph{bú yào} ; \zh{你可以} → \emph{ní kéyǐ} (deux sandhis : \emph{nǐ} devient \emph{ní} devant \emph{kě}, et \emph{kě} devient \emph{ké} devant \emph{yǐ}).\\
3. Comptine : \zh{妈妈骑马，马慢，妈妈骂马。} \emph{Māma qí mǎ, mǎ màn, māma mà mǎ.} — vérifier que le 1\textsuperscript{er} ton reste plat, que \emph{mǎ} creuse bien, et que \emph{mà} tombe sec.
\end{corrige}

\begin{savaistu}[Les tons multiplient les mots]
Le mandarin ne possède qu'environ 400 syllabes différentes si l'on ignore les tons — très peu comparé au français ! Mais en ajoutant les quatre tons, on obtient plus de 1 200 syllabes distinctes, et donc des milliers de mots possibles. Les tons ne sont pas un « décor » : ils sont une machine à fabriquer du vocabulaire. C'est aussi pour cela que les Chinois eux-mêmes, au téléphone, précisent parfois un mot en citant une expression connue qui le contient !
\end{savaistu}

\begin{aretenir}[L'essentiel de la section]
\begin{itemize}
\item Quatre tons : haut-plat (ā), montant (á), descendant-montant (ǎ), descendant (à), plus le ton neutre.
\item Un mauvais ton = un autre mot : \zh{买} mǎi « acheter » / \zh{卖} mài « vendre ».
\item 3\textsuperscript{e} + 3\textsuperscript{e} → 2\textsuperscript{e} + 3\textsuperscript{e} : \zh{你好} se prononce \emph{ní hǎo} ; avec trois 3\textsuperscript{e} tons, on applique la règle de proche en proche : \zh{我很好} → \emph{wó hén hǎo}.
\item \zh{不} devient \emph{bú} devant un 4\textsuperscript{e} ton : \zh{不是} → \emph{bú shì}.
\item On écrit toujours le ton d'origine ; le sandhi ne concerne que la prononciation.
\item La main qui dessine le contour mélodique est votre meilleur outil de mémorisation.
\end{itemize}
\end{aretenir}

Maintenant que notre oreille et notre bouche sont entraînées, nous pouvons passer au \textbf{dialogue modèle 2}, où vous donnerez votre avis sur les fêtes et la musique — avec des tons impeccables !

\coursec{Dialogue modèle 2 : donner son avis et débattre}

Après avoir travaillé les tons sur des mots isolés et des phrases courtes (voir \emph{Les tons : entraînement ciblé}), il est temps de les mettre en pratique dans une vraie situation de communication : le \cle{débat}. Débattre en chinois, c'est savoir donner son avis, dire qu'on n'est pas d'accord \emph{poliment}, défendre son idée avec un argument, puis conclure. C'est exactement ce que font les deux élèves du dialogue ci-dessous, sur un sujet qui vous parle : la musique camerounaise et la musique chinoise.

\courssub{Le dialogue à observer}

\begin{textebox}[La musique camerounaise ou la musique chinoise ?]
A : 我觉得喀麦隆的音乐比中国的音乐热闹。\\
\emph{Wǒ juéde Kāmàilóng de yīnyuè bǐ Zhōngguó de yīnyuè rènao.}\\
« Je trouve la musique camerounaise plus animée que la musique chinoise. »\\[0.4em]
B : 我不太同意。我觉得京剧很有意思，虽然有点难懂。\\
\emph{Wǒ bù tài tóngyì. Wǒ juéde Jīngjù hěn yǒuyìsi, suīrán yǒudiǎn nán dǒng.}\\
« Je ne suis pas tout à fait d'accord. Je trouve l'opéra de Pékin très intéressant, bien qu'un peu difficile à comprendre. »\\[0.4em]
A : 但是马库萨很有名，很多非洲人都喜欢。\\
\emph{Dànshì Mǎkùsà hěn yǒumíng, hěn duō Fēizhōurén dōu xǐhuan.}\\
« Mais le makossa est célèbre, beaucoup d'Africains l'aiment. »\\[0.4em]
B : 你说得对。不过中国的传统音乐也很美。你呢，你喜欢什么？\\
\emph{Nǐ shuō de duì. Bùguò Zhōngguó de chuántǒng yīnyuè yě hěn měi. Nǐ ne, nǐ xǐhuan shénme ?}\\
« Tu as raison. Mais la musique traditionnelle chinoise est aussi très belle. Et toi, qu'aimes-tu ? »\\[0.4em]
A : 我两个都喜欢！对我来说，音乐没有国界。\\
\emph{Wǒ liǎng gè dōu xǐhuan ! Duì wǒ lái shuō, yīnyuè méiyǒu guójiè.}\\
« J'aime les deux ! Pour moi, la musique n'a pas de frontières. »
\end{textebox}

\textbf{Glossaire du dialogue}\par

\begin{center}
\begin{tabularx}{\linewidth}{|X|X|X|X|}
\hline
\rowcolor{popblueL} Caractères & Pinyin & Nature & Traduction \\
\hline
觉得 & juéde & verbe & trouver, penser, estimer \\
\hline
热闹 & rènao & adjectif & animé, bruyant et joyeux \\
\hline
同意 & tóngyì & verbe & être d'accord \\
\hline
京剧 & Jīngjù & nom & opéra de Pékin \\
\hline
有意思 & yǒuyìsi & expression & intéressant, amusant \\
\hline
难懂 & nán dǒng & adjectif & difficile à comprendre \\
\hline
有名 & yǒumíng & adjectif & célèbre \\
\hline
马库萨 & Mǎkùsà & nom propre & makossa (genre musical) \\
\hline
不过 & bùguò & conjonction & cependant, mais (adouci) \\
\hline
传统 & chuántǒng & nom / adjectif & tradition ; traditionnel \\
\hline
国界 & guójiè & nom & frontière (entre pays) \\
\hline
\end{tabularx}
\end{center}

\courssub{Les tons : entraînement ciblé sur le dialogue}

Avant de jouer le dialogue, isolons les difficultés tonales de ce texte. Le chinois standard (putonghua) impose des règles de \cle{changement de ton} (sandhi) à l'oral que l'écriture pinyin ne marque pas toujours.

\begin{propriete}[Rappel des sandhi obligatoires pour ce dialogue]
\begin{enumerate}
\item \textbf{Troisième ton + Troisième ton} : le premier devient un deuxième ton (montant).\\
Ex. : \zh{很好} \rightarrow \emph{hén hǎo} (pas \emph{hěn hǎo}) ; \zh{马库萨} \emph{Mǎkùsà} (pas de sandhi ici, 3-4) ; \zh{你呢} \rightarrow \emph{ní ne} (3 + neutre).
\item \textbf{不 (bù) devant un 4\textsuperscript{e} ton} : \emph{bù} devient \emph{bú} (2\textsuperscript{e} ton).\\
Ex. : \zh{不太} \rightarrow \emph{bú tài} ; \zh{不懂} \rightarrow \emph{bú dǒng} ; \zh{不同意} \rightarrow \emph{bú tóngyì}.
\item \textbf{一 (yī) :} seul ou à la fin = 1\textsuperscript{er} ton ; devant 4\textsuperscript{e} ton = 2\textsuperscript{e} ton (\emph{yí}) ; devant 1/2/3\textsuperscript{e} ton = 4\textsuperscript{e} ton (\emph{yì}).\\
Ex. : \zh{一个} \rightarrow \emph{yí gè} ; \zh{一点} \rightarrow \emph{yì diǎn}.
\item \textbf{Ton neutre} : \zh{的 de}, \zh{呢 ne}, \zh{个 gè} (dans \zh{两个}) perdent leur ton propre et s'attachent à la syllabe précédente.
\end{enumerate}
\end{propriete}

\begin{experience}[Entraînement tonal en trois temps]
En binôme, travaillez ces phrases extraites du dialogue :
\begin{enumerate}
\item \textbf{Repérage} : soulignez les sandhi \emph{bù} (devant 4\textsuperscript{e} ton) et les 3\textsuperscript{e} tons consécutifs.
\item \textbf{Écoute-Modèle} : l'enseignant lit, les élèves marquent les contours de ton dans l'air (gestuelle main).
\item \textbf{Chorale puis individuel} : lecture chorale lente, puis passage individuel sur : \zh{我不太同意} (\emph{Wǒ \textbf{bú} tài tóngyì}), \zh{虽然有点难懂} (\emph{Suīrán yǒudiǎn nán dǒng}), \zh{音乐没有国界} (\emph{Yīnyuè méiyǒu guójiè}).
\end{enumerate}
\end{experience}

\courssub{Analyse du dialogue : la mécanique du débat}

Relisons le dialogue en suivant les \emph{coups} échangés, comme au ping-pong. Chaque réplique joue un rôle précis :

\begin{itemize}
\item \textbf{Avis de A} : \zh{我觉得……比……热闹。} — A donne son opinion avec \zh{我觉得} et la structure comparative \zh{A 比 B + adjectif}.
\item \textbf{Désaccord poli de B} : \zh{我不太同意。} — B n'est pas d'accord, mais le \zh{太} adoucit le refus : « pas tout à fait ».
\item \textbf{Contre-argument de B} : \zh{我觉得京剧很有意思，虽然有点难懂。} — B défend son camp avec la structure \zh{虽然……（但是）……}.
\item \textbf{Argument de A} : \zh{但是马库萨很有名，很多非洲人都喜欢。} — A répond avec un fait : la célébrité du makossa.
\item \textbf{Concession de B} : \zh{你说得对。不过……} — B reconnaît le point de A, puis nuance avec \zh{不过}.
\item \textbf{Relance} : \zh{你呢，你喜欢什么？} — B rend la parole à A.
\item \textbf{Conclusion commune} : \zh{我两个都喜欢！音乐没有国界。} — sortie élégante qui réconcilie les deux positions.
\end{itemize}

\begin{popfigure}
\begin{tikzpicture}[
  node distance=0.35cm and 0.9cm,
  coup/.style={rounded corners=4pt, draw=popline, thick, align=center, font=\small, inner sep=5pt, text width=3.1cm},
  fleche/.style={-{Stealth[length=2.5mm]}, thick, popdark}
]
\node[coup, fill=popblueL, align=center] (n1){Avis de A\\ 我觉得……比……};
\node[coup, fill=poppinkL, right=of n1, align=center] (n2){Désaccord poli\\ 我不太同意};
\node[coup, fill=popblueL, right=of n2, align=center] (n3){Argument + fait\\ 但是马库萨很有名};
\node[coup, fill=poppinkL, below=0.7cm of n3, align=center] (n4){Concession\\ 你说得对。不过……};
\node[coup, fill=poporangeL, left=of n4, align=center] (n5){Relance\\ 你呢？};
\node[coup, fill=popgreenL, left=of n5, align=center] (n6){Conclusion\\ 音乐没有国界};
\draw[fleche] (n1) -- (n2);
\draw[fleche] (n2) -- (n3);
\draw[fleche] (n3) -- (n4);
\draw[fleche] (n4) -- (n5);
\draw[fleche] (n5) -- (n6);
\end{tikzpicture}
\legende{Le « ping-pong du débat » : avis, désaccord poli, argument, concession, relance, conclusion commune.}
\end{popfigure}

\courssub{La structure 虽然……但是…… : « bien que... mais... »}

Dans la réplique de B, observez bien : \zh{我觉得京剧很有意思，虽然有点难懂。} Le chinois exprime les \emph{deux} moitiés de l'opposition, là où le français n'en garde souvent qu'une.

\begin{propriete}[虽然……但是…… : la concession à deux têtes]
\textbf{Structure} : 虽然 + proposition 1，但是／可是／不过 + proposition 2.\par
\textbf{Emploi} : la première partie (虽然) reconnaît un fait défavorable ou un contre-argument ; la seconde (但是) affirme l'idée principale. Contrairement au français, où « bien que » et « mais » ne se combinent pas dans la même phrase (*« Bien qu'il soit difficile, mais il est beau » est fautif), le chinois \textbf{utilise volontiers les deux ensemble} pour marquer clairement la concession et la rupture.\par
\textbf{Variantes} : 但是 peut être remplacé par 可是 (kěshì, plus oral) ou 不过 (bùguò, plus doux, « toutefois »). On peut aussi placer 虽 puis 但/却 dans un style plus littéraire (\zh{虽难，却美}), mais à l'oral standard, garder les deux mots complets est plus clair.
\end{propriete}

\begin{exemplebox}[虽然……但是…… en action]
\zh{虽然京剧有点难，但是很有意思。} \emph{Suīrán Jīngjù yǒudiǎn nán, dànshì hěn yǒuyìsi.} « Bien que l'opéra de Pékin soit un peu difficile, il est très intéressant. »\\[0.3em]
\zh{虽然汉语的声调很难，可是我很喜欢。} \emph{Suīrán Hànyǔ de shēngdiào hěn nán, kěshì wǒ hěn xǐhuan.} « Bien que les tons du chinois soient difficiles, je les aime beaucoup. »\\[0.3em]
\zh{虽然雅温得离杜阿拉很远，不过路很好走。} \emph{Suīrán Yǎwēndé lí Dù'ālā hěn yuǎn, bùguò lù hěn hǎo zǒu.} « Bien que Yaoundé soit loin de Douala, la route est bonne. »\\[0.3em]
\zh{虽然今天很热，但是我们要去踢足球。} \emph{Suīrán jīntiān hěn rè, dànshì wǒmen yào qù tī zúqiú.} « Bien qu'il fasse très chaud aujourd'hui, nous allons jouer au foot. »
\end{exemplebox}

\begin{attention}[« Je suis d'accord avec toi » : pas toujours 我同意你]
En français on dit « je suis d'accord \emph{avec toi} », mais en chinois on ne dit pas systématiquement \zh{我同意你} (sonne administratif ou incomplet). Le plus naturel est \zh{我同意。} tout court, ou \zh{我同意你的看法。} \emph{Wǒ tóngyì nǐ de kànfǎ.} « Je suis d'accord avec ton point de vue. » Pour reconnaître qu'un argument est juste dans une conversation, la formule du dialogue est idéale : \zh{你说得对。} \emph{Nǐ shuō de duì.} « Ce que tu dis est juste / Tu as raison. »
\end{attention}

\courssub{Tableau récapitulatif des formules du débat}

\begin{center}
\begin{tabularx}{\linewidth}{|X|X|X|}
\hline
\rowcolor{popblueL} Fonction & Formule chinoise + pinyin & Traduction \\
\hline
Donner son avis & 我觉得…… Wǒ juéde... & Je trouve que... / Je pense que... \\
\hline
Donner son avis (variante) & 对我来说…… Duì wǒ lái shuō... & Pour moi... / À mon avis... \\
\hline
Être d'accord & 我同意。/ 你说得对。 Wǒ tóngyì. / Nǐ shuō de duì. & Je suis d'accord. / Tu as raison. \\
\hline
Désaccord poli & 我不太同意。 Wǒ bù tài tóngyì. & Je ne suis pas tout à fait d'accord. \\
\hline
Désaccord franc (à éviter au début) & 我不同意。 Wǒ bù tóngyì. & Je ne suis pas d'accord. (trop direct) \\
\hline
Nuancer / Contredire & 不过…… / 可是…… Bùguò... / Kěshì... & Cependant... / Mais... \\
\hline
Concéder (structure) & 虽然……但是…… Suīrán... dànshì... & Bien que... (mais)... \\
\hline
Relancer l'échange & 你呢？ / 你喜欢什么？ Nǐ ne? / Nǐ xǐhuan shénme? & Et toi ? / Qu'aimes-tu ? \\
\hline
Conclure (universel) & 音乐没有国界。 Yīnyuè méiyǒu guójiè. & La musique n'a pas de frontières. \\
\hline
\end{tabularx}
\end{center}

La dernière formule est un petit trésor : \zh{音乐没有国界。} « La musique n'a pas de frontières. » C'est une \cle{phrase de conclusion} parfaite pour terminer un exposé ou un débat sur la culture, car elle dépasse l'opposition Cameroun / Chine et ouvre sur l'universel. Retenez aussi sa structure utile : \zh{X 没有 Y} « X n'a pas de Y » (ex. \zh{美食没有国界} « La gastronomie n'a pas de frontières »).

\courssub{Variantes de substitution : musique, cuisine, films, fêtes}

Le dialogue est un \emph{moule} : gardez les formules et remplacez le sujet. Voici les mots à injecter pour varier les thèmes culturels :

\begin{center}
\begin{tabularx}{\linewidth}{|X|X|X|X|}
\hline
\rowcolor{popblueL} Caractères & Pinyin & Nature & Traduction \\
\hline
音乐 & yīnyuè & nom & musique \\
\hline
菜 / 中国菜 & cài / Zhōngguó cài & nom & plat / cuisine chinoise \\
\hline
电影 & diànyǐng & nom & film, cinéma \\
\hline
节日 & jiérì & nom & fête, festival \\
\hline
好吃 & hǎochī & adjectif & bon (au goût) \\
\hline
好看 & hǎokàn & adjectif & beau, intéressant (à voir) \\
\hline
辣 & là & adjectif & piquant, épicé \\
\hline
精彩 & jīngcǎi & adjectif & splendide, palpitant (spectacle, match) \\
\hline
\end{tabularx}
\end{center}

Exemple de substitution avec la cuisine : \zh{我觉得喀麦隆的菜比中国的菜辣。} \emph{Wǒ juéde Kāmàilóng de cài bǐ Zhōngguó de cài là.} « Je trouve la cuisine camerounaise plus piquante que la cuisine chinoise. » — Réponse possible : \zh{我不太同意。虽然四川菜很辣，但是很好吃。} \emph{Wǒ bù tài tóngyì. Suīrán Sìchuān cài hěn là, dànshì hěn hǎochī.} « Je ne suis pas tout à fait d'accord. Bien que la cuisine du Sichuan soit très piquante, elle est délicieuse. »

Exemple avec les films : \zh{我觉得中国的电影比喀麦隆的电影精彩。} \emph{Wǒ juéde Zhōngguó de diànyǐng bǐ Kāmàilóng de diànyǐng jīngcǎi.} — \zh{我不太同意。虽然中国电影很精彩，但是喀麦隆的电影也很有意思。}

\begin{methode}[Structurer un mini-débat en quatre coups]
\begin{enumerate}
\item \textbf{Avis} : annoncez votre position avec \zh{我觉得} ou \zh{对我来说} (+ \zh{比} si vous comparez deux choses).
\item \textbf{Argument} : justifiez avec un fait objectif (\zh{很有名}, \zh{很多人都喜欢}, \zh{有历史}...).
\item \textbf{Exemple concret} : citez un cas précis connu de l'interlocuteur (le makossa, l'opéra de Pékin, le ndolé, le Nouvel An chinois, la Coupe d'Afrique...).
\item \textbf{Question à l'autre} : rendez la parole avec \zh{你呢？} ou \zh{你怎么看？} — un débat sans question est un monologue !
\end{enumerate}
\end{methode}

\begin{exoresolu}[Compléter un échange de débat]
Énoncé : Complétez la réplique de B avec une formule de désaccord poli et la structure 虽然……但是……\par
A : 我觉得中国的电影比喀麦隆的电影好看。\\
\emph{Wǒ juéde Zhōngguó de diànyǐng bǐ Kāmàilóng de diànyǐng hǎokàn.}\\
B : \_\_\_\_\_\_\_\_\_\_\_\_
\tcblower
Solution : B doit d'abord marquer un désaccord adouci, puis concéder un point avant de défendre son idée. Réponse modèle :\\
\zh{我不太同意。虽然中国的电影很好看，但是喀麦隆的电影也很有意思。}\\
\emph{Wǒ bù tài tóngyì. Suīrán Zhōngguó de diànyǐng hěn hǎokàn, dànshì Kāmàilóng de diànyǐng yě hěn yǒuyìsi.}\\
« Je ne suis pas tout à fait d'accord. Bien que les films chinois soient beaux, les films camerounais sont aussi très intéressants. »\\
Points de méthode : (1) \zh{不太同意} et non \zh{不同意} seul, trop sec ; (2) les deux moitiés \zh{虽然} / \zh{但是} sont présentes ; (3) l'adverbe \zh{也} « aussi » rééquilibre le débat ; (4) l'adjectif \zh{有意思} varie le lexique par rapport à \zh{好看}.
\end{exoresolu}

\begin{experience}[Jouer le dialogue en trois passes]
En binôme, jouez le dialogue \textbf{trois fois} :
\begin{enumerate}
\item \textbf{Lecture soignée} : lisez le texte en soignant les tons (attention aux sandhi : \zh{yīnyuè}, \zh{bú tài tóngyì}, \zh{suīrán}, \zh{liǎng gè}). Vérifiez la prononciation de \zh{马库萨} (Mǎkùsà).
\item \textbf{Formules masquées} : le professeur cache les formules clés (\zh{我觉得}, \zh{我不太同意}, \zh{你说得对}, \zh{不过}, \zh{你呢？}) : retrouvez-les de mémoire sans lire le texte.
\item \textbf{Improvisation guidée} : remplacez le thème « musique » par « cuisine », « films » ou « fêtes » (vocabulaire du tableau de substitution). Terminez par votre propre phrase de conclusion universelle, par exemple : \zh{美食没有国界。} \emph{Měishí méiyǒu guójiè.} « La gastronomie n'a pas de frontières. » ou \zh{友谊没有国界。} \emph{Yǒuyì méiyǒu guójiè.} « L'amitié n'a pas de frontières. »
\end{enumerate}
\end{experience}

\coursec{Leçon 21 — Expression orale : vie culturelle au Cameroun et en Chine}

Cette leçon constitue l'aboutissement du Module 3 « Citoyenneté et ouverture au monde ». Après l'acquisition du lexique culturel (Section 1), l'étude de deux dialogues modèles (Sections 2 et 5), la maîtrise des formules fonctionnelles pour présenter, opiner et relancer (Section 3) et l'entraînement phonétique sur les tons (Section 4), nous passons à la mise en situation authentique. L'objectif est de permettre à l'élève de produire un discours oral cohérent, prononcé correctement et interactif, dans le cadre de l'épreuve orale du Baccalauréat (LV2).

\courssub{Jeux de rôle et débat guidé}

\courssub{Organisation générale de la séance}

\begin{popfigure}
\begin{tikzpicture}[
    node distance=1.3cm and 1.8cm,
    every node/.style={font=\small, align=center},
    box/.style={rectangle, draw=popblue, thick, fill=popblueL, rounded corners=6pt, minimum width=3.5cm, minimum height=1.8cm},
    arrow/.style={-Stealth, thick, popdark}
]
\node[box, align=center] (prep){Préparer\\ (2 min)\\\textit{Notes, mots-clés,}\\\textit{structure}};
\node[box, right=of prep, align=center] (play){Jouer / Parler\\ (2--3 min)\\\textit{Regard vers l'auditoire,}\\\textit{tons, formules}};
\node[box, right=of play, align=center] (listen){Écouter les autres\\ \& Grille d'observation\\\textit{Tons ? Vocabulaire ?}\\\textit{Formules ?}};
\node[box, below=of play, align=center] (vote){Vote / Retour\\ (1 min)\\\textit{Critères de réussite}};
\node[box, left=of vote, align=center] (switch){Changer de rôle\\ \& Recommencer};

\draw[arrow] (prep) -- (play) node[midway, above] {Prêt ?};
\draw[arrow] (play) -- (listen) node[midway, above] {Suivant};
\draw[arrow] (listen) -- (vote) node[midway, right] {Évaluer};
\draw[arrow] (vote) -- (switch) node[midway, below] {Tour suivant};
\draw[arrow] (switch.west) |- (prep.west) node[midway, left] {Boucle itérative};
\end{tikzpicture}
\legende{Organigramme du déroulement d'un jeu de rôle / mini-exposé (boucle itérative)}
\end{popfigure}

\courssub{Jeu de rôle 1 : « Au stand culturel » (立文化展位)}

\begin{textebox}[Situation de départ]
\textbf{Contexte :} La Semaine culturelle sino-camerounaise à Yaoundé. Un stand présente les fêtes du Cameroun.
\textbf{Rôle A (Élève camerounais) :} Vous tenez le stand. Présentez une fête (Ngondo, Ngoun, Medumba, Fête du Maïs, etc.) à un visiteur chinois.
\textbf{Rôle B (Visiteur chinois) :} Vous êtes curieux. Posez des questions précises (date, lieu, sens, plats, vêtements).
\end{textebox}

\begin{methode}[Structure type pour le Rôle A (Présentateur)]
\begin{enumerate}
    \item \textbf{Accroche \& Salut :} \zh{你好！欢迎来到喀麦隆文化展位。} (Nǐ hǎo! Huānyíng lái dào Kāmàilóng wénhuà zhǎnwèi.) — Bonjour ! Bienvenue au stand culturel du Cameroun.
    \item \textbf{Annonce du sujet :} \zh{我来介绍一下……节。} (Wǒ lái jièshào yīxià … jié.) — Je vais vous présenter la fête …
    \item \textbf{Trois informations clés (Quand / Où / Quoi) :}
    \begin{itemize}
        \item \zh{时间：每年……月。} (Shíjiān: měinián … yuè.) — Date : chaque année … mois.
        \item \zh{地点：在……地区。} (Dìdiǎn: zài … dìqū.) — Lieu : dans la région de …
        \item \zh{活动：有赛船、跳舞、吃……} (Huódòng: yǒu sàichuán, tiàowǔ, chī …) — Activités : courses de pirogues, danses, manger …
    \end{itemize}
    \item \textbf{Avis personnel :} \zh{我觉得这个节日很有意义，因为……} (Wǒ juéde zhège jiérì hěn yǒu yìyì, yīnwèi …) — Je trouve cette fête très significative, car …
    \item \textbf{Conclusion \& Invitation :} \zh{欢迎你将来去体验一下！谢谢！} (Huānyíng nǐ jiānglái qù tǐyàn yīxià! Xièxie!) — Bienvenue pour aller l'expérimenter un jour ! Merci !
\end{enumerate}
\end{methode}

\begin{exemplebox}[Exemple joué — Fête du Ngondo (诺贡多节)]
\begin{tabular}{@{}p{0.15\linewidth}p{0.8\linewidth}@{}}
\textbf{A :} & \zh{你好！欢迎来到喀麦隆文化展位。我来介绍一下诺贡多节。} (Nǐ hǎo! Huānyíng lái dào Kāmàilóng wénhuà zhǎnwèi. Wǒ lái jièshào yīxià Nuògòngduō jié.) \\
\textbf{B :} & \zh{诺贡多节？听起来很有趣。请问是什么时候？} (Nuògòngduō jié? Tīng qǐlái hěn yǒuqù. Qǐngwèn shì shénme shíhou?) \\
\textbf{A :} & \zh{每年十一月或十二月，在杜阿拉，瓦里河边。} (Měinián shíyī yuè huò shíèr yuè, zài Dù'ālā, Wǎlǐ hé biān.) \\
\textbf{B :} & \zh{有什么特别的活动吗？} (Yǒu shénme tèbié de huódòng ma?) \\
\textbf{A :} & \zh{有传统的赛船，非常壮观。还有舞蹈和烤鱼。我觉得这个节日展示了萨瓦人的团结精神。} (Yǒu chuántǒng de sàichuán, fēicháng zhuāngguān. Hái yǒu wǔdǎo hé kǎoyú. Wǒ juéde zhège jiérì zhǎnshì le Sàwǎ rén de tuánjié jīngshén.) \\
\textbf{B :} & \zh{太棒了！赛船听起来很刺激。我也想去看看。谢谢介绍！} (Tài bàng le! Sàichuán tīng qǐlái hěn cìjī. Wǒ yě xiǎng qù kànkan. Xièxie jièshào!) \\
\end{tabular}
\end{exemplebox}

\begin{attention}[Piège à éviter : « Lire son exposé »]
Au Baccalauréat, l'épreuve orale (Exposé + Entretien) sanctionne sévèrement la \textbf{lecture à voix haute} d'un texte préparé par cœur. Le jury attend une \textbf{prise de parole vivante} : regard vers l'auditoire, hésitations naturelles, reformulations.
\par \textbf{Consigne stricte :} Vous avez le droit à une \emph{fiche bristol} (format carte de visite) avec \textbf{5 mots-clés maximum} (ex: \zh{诺贡多节、赛船、团结、烤鱼、欢迎}). Interdiction d'écrire des phrases complètes en caractères ou en pinyin.
\end{attention}

\courssub{Jeu de rôle 2 : « Interview » (采访艺术家)}

\begin{textebox}[Situation de départ]
\textbf{Contexte :} Émission radio « Voix de la Jeunesse » (青春之声). Un journaliste interviewe un artiste.
\textbf{Rôle A (Journaliste) :} Utilisez les formules de relance (\zh{请问……}, \zh{能不能具体说说？}, \zh{为什么？}). Notez les réponses pour un résumé final.
\textbf{Rôle B (Artiste — Musicien / Danseur / Cuisinier) :} Choisissez une identité (Camerounais ou Chinois). Présentez votre art, votre parcours, une anecdote.
\end{textebox}

\begin{propriete}[Formules types pour le Journaliste (Relances)]
\begin{center}
\begin{tabularx}{\linewidth}{|X|X|X|}
\hline
\rowcolor{popblueL} \textbf{Fonction} & \textbf{Chinois (Pinyin)} & \textbf{Français} \\ \hline
Question ouverte & \zh{请问您是怎么开始学……的？} (Qǐngwèn nín shì zěnme kāishǐ xué … de?) & Comment avez-vous commencé à apprendre … ? \\
Demander détail & \zh{能不能具体说说？} (Néng bùnéng jùtǐ shuōshuo?) & Pourriez-vous donner des détails ? \\
Demander raison & \zh{为什么选择这个领域？} (Wèishénme xuǎnzé zhège lǐngyù?) & Pourquoi choisir ce domaine ? \\
Réaction / Intérêt & \zh{真有意思！/ 太厉害了！} (Zhēn yǒu yìsi! / Tài lìhai le!) & Vraiment intéressant ! / Trop fort ! \\
Transition & \zh{最后一个问题：……} (Zuìhòu yīgè wèntí: …) & Dernière question : … \\ \hline
\end{tabularx}
\end{center}
\end{propriete}

\begin{exoresolu}[Simulation d'interview — Cuisinier chinois (饺子)]
\textbf{Énoncé :} Complétez les répliques du journaliste (A) pour relancer naturellement le cuisinier (B).
\begin{itemize}
    \item A: \zh{你好！请问您叫什么名字？}
    \item B: \zh{我叫王师傅，做饺子三十年了。}
    \item A: \zh{\textbf{三十年？真厉害！请问您是怎么开始学做饺子的？}} (Relance : durée + question origine)
    \item B: \zh{小时候跟奶奶学的，家里穷，饺子是过年才能吃的好东西。}
    \item A: \zh{\textbf{能不能具体说说饺子的馅儿怎么做？}} (Relance : détail technique)
    \item B: \zh{猪肉白菜最经典，关键是要把水挤干，放香油……}
    \item A: \zh{\textbf{听得我都饿了！最后一个问题：您觉得好吃的秘诀是什么？}} (Transition + question avis)
    \item B: \zh{用心，还有鲜肉。谢谢！}
    \item A: \zh{\textbf{谢谢王师傅的分享，祝您生意兴隆！}} (Conclusion polie)
\end{itemize}
\tcblower
\textbf{Solution détaillée :} Les relances en gras montrent l'écoute active (reprise de « 30 ans » → « 真厉害 »), l'usage de \zh{请问} pour la politesse, \zh{能不能具体说说} pour le détail, et une conclusion formelle \zh{祝您生意兴隆} (Souhaitez prospérité aux affaires).
\end{exoresolu}

\courssub{Débat guidé : « Faut-il préférer les fêtes traditionnelles ou modernes ? » (传统节日 vs 现代节日)}

\begin{methode}[Déroulement du débat (2 camps, 10 min total)]
\begin{enumerate}
    \item \textbf{Distribution des cartes de rôle (1 min) :} Chaque élève tire une carte imposant une opinion à défendre (voir ci-dessous). \emph{Important :} On défend la carte, pas son avis personnel.
    \item \textbf{Préparation individuelle (2 min) :} Rédigez 3 arguments en mots-clés + 3 formules d'opinion/relance obligatoires.
    \item \textbf{Débat en « Ping-Pong » (5 min) :} Camp A (Tradition) argue → Camp B (Moderne) répond/relance → Camp A relance… Le professeur note les formules utilisées au tableau.
    \item \textbf{Synthèse / Vote (2 min) :} Un élève par camp fait une phrase de conclusion. Vote à main levée : « Qui a été le plus convaincant linguistiquement ? » (tons, vocabulaire, structure).
\end{enumerate}
\end{methode}

\begin{definition}[Cartes de rôle — Opinions imposées]
Imprimez et découpez (format carte de visite). Chaque élève en pioche une.
\begin{center}
\begin{tabularx}{\linewidth}{|X|X|}
\hline
\rowcolor{poporangeL} \textbf{Camp A : Tradition (传统派)} & \textbf{Camp B : Modernité (现代派)} \\ \hline
\zh{你认为传统节日是文化的根，必须保护。} (Racines culturelles, protection nécessaire) & \zh{你认为现代节日(如双十一、情人节)更适应年轻人生活。} (Adaptés aux jeunes) \\ \hline
\zh{你觉得传统节日能增强家庭凝聚力。} (Renforcent cohésion familiale) & \zh{你觉得传统节日太繁琐，现代节日轻松快乐。} (Traditionnels complexes, modernes simples) \\ \hline
\zh{你主张学校应开设传统礼仪课。} (École devrait enseigner étiquette trad.) & \zh{你主张节日应创新，结合科技(如VR拜年)。} (Innover, technologie/VR) \\ \hline
\zh{你担心年轻人只会过「洋节」，忘本。} (Peur : fêtes étrangères, oubli racines) & \zh{你认为文化要发展，不能只守旧。} (Culture doit évoluer, pas stagner) \\ \hline
\zh{你提议用短视频传播传统节日习俗。} (Proposition : TikTok/Douyin pour coutumes) & \zh{你提议设立「文化融合节」，中喀混合。} (Proposition : Fête fusion Sino-Cam) \\ \hline
\end{tabularx}
\end{center}
\end{definition}

\begin{aretenir}[Obligations linguistiques du débat]
Chaque intervenant \textbf{DOIT} utiliser au minimum :
\begin{itemize}
    \item \textbf{3 formules d'opinion} différentes (ex: \zh{我认为……}, \zh{在我看来……}, \zh{我赞成……}, \zh{我反对……}, \zh{我不得不说……}).
    \item \textbf{2 relances} adressées à l'adversaire (ex: \zh{但是……}, \zh{不过你有没有想过……?}, \zh{请问你怎么看……?}).
    \item \textbf{4 mots du lexique} de la leçon (ex: \zh{传承}, \zh{创新}, \zh{仪式感}, \zh{商业化}, \zh{认同感}, \zh{代沟}).
\end{itemize}
\end{aretenir}

\begin{exemplebox}[Extrait de débat modèle]
\begin{tabular}{@{}p{0.15\linewidth}p{0.8\linewidth}@{}}
\textbf{A (Trad)} : & \zh{我认为传统节日是文化的根。比如春节的团圆饭，有很强的仪式感，能增强家庭认同感。} (Racines, repas réunion, sens rituel, identité familiale) \\
\textbf{B (Mod)} : & \zh{我赞成传统有价值，\textbf{但是}你有没有想过，现在很多年轻人加班太累，没时间准备繁琐的仪式？} (Accord partiel + Relance \textbf{但是} + question) \\
\textbf{A (Trad)} : & \zh{我不得不说商业化确实严重。\textbf{不过}我们可以创新形式，比如用短视频教包饺子，让传统活起来。} (Concession + Relance \textbf{不过} + proposition) \\
\textbf{B (Mod)} : & \zh{这个主意不错！我也赞成「文化融合节」，中喀青年一起过节，增进友谊。请问你怎么看？} (Appui + Proposition + Relance \textbf{请问}) \\
\end{tabular}
\end{exemplebox}

\courssub{Tournoi de mini-exposés : « Ma fête préférée » (我最喜欢的节日)}

\begin{experience}[Règles du tournoi]
\begin{enumerate}
    \item \textbf{Préparation (3 min) :} Chaque élève prépare un exposé de \textbf{5 phrases exactes} sur sa fête préférée (camerounaise ou chinoise). Utilisez la structure de la \textbf{Méthode} p. 2.
    \item \textbf{Passage (1 min/élève) :} Debout, sans notes (ou fiche 5 mots), regard vers la classe.
    \item \textbf{Évaluation par les pairs :} Chaque élève remplit la \textbf{Grille d'observation} (ci-dessous) pour 3 camarades tirés au sort.
    \item \textbf{Vote final :} La classe vote pour le « Prix de la Clarté » (prononciation + structure + regard).
\end{enumerate}
\end{experience}

\begin{textebox}[Grille d'auto-évaluation / Observation par les pairs]
\emph{À remplir pendant l'écoute (cocher OUI / NON). Nom de l'orateur : \underline{\hspace{4cm}}}
\begin{center}
\begin{tabularx}{\linewidth}{|X|c|c|}
\hline
\rowcolor{popgreenL} \textbf{Critère de réussite} & \textbf{OUI} & \textbf{NON} \\ \hline
1. A salué correctement (\zh{大家好 / 你好}) & \square & \square \\ \hline
2. A annoncé le sujet clairement (\zh{我来介绍一下……}) & \square & \square \\ \hline
3. A donné au moins 3 infos (date, lieu, activité, plat) & \square & \square \\ \hline
4. A exprimé son avis avec \zh{我觉得 / 我认为 / 在我看来} & \square & \square \\ \hline
5. A utilisé au moins \textbf{4 mots du lexique} de la leçon & \square & \square \\ \hline
6. A utilisé au moins \textbf{2 formules fonctionnelles} (relance, transition) & \square & \square \\ \hline
7. \textbf{Tons corrects} (pas de ton 3 qui traîne, ton 4 net) & \square & \square \\ \hline
8. \textbf{Regard vers l'auditoire} (ne lit pas ses notes) & \square & \square \\ \hline
9. Fluidité : phrases liées, pas de longs silences & \square & \square \\ \hline
10. Conclusion polie (\zh{谢谢大家 / 谢谢聆听}) & \square & \square \\ \hline
\end{tabularx}
\end{center}
\textbf{Note globale (sur 10) :} \underline{\hspace{2cm}} \quad \textbf{Commentaire constructif :} \underline{\hspace{6cm}}
\end{textebox}

\begin{savaistu}[Le risque payant à l'oral]
À l'examen du Baccalauréat (LV2), l'épreuve orale dure environ 10 minutes (Exposé 3-4 min + Entretien 6-7 min). Les correcteurs valorisent la \textbf{prise de risque maîtrisée} : 
\begin{itemize}
    \item Mieux vaut dire \zh{我喜欢春节，因为吃饺子，全家在一起，很开心。} (Wǒ xǐhuan Chūnjié, yīnwèi chī jiǎozi, quánjiā zài yīqǐ, hěn kāixīn.) — 100\% compréhensible, tons justes, structure SVO simple. 
    \item Que d'essayer \zh{由于春节具有深厚的文化内涵以及家庭团聚的社会功能，使得我对其产生浓厚的情感依恋……} (Yóuyú Chūnjié jùyǒu shēnhòu de wénhuà nèihán yǐjí jiātíng tuánjù de shèhuì gōngnéng, shǐde wǒ duì qí chǎnshēng nónghòu de qínggǎn yīliàn……) — Risque d'erreurs de tons, d'ordre des mots \zh{由于……使得}, de registre trop écrit, de bégaiement.
\end{itemize}
\cle{Stratégie gagnante :} Maîtrisez 10 « phrases-modèles » parfaites (tons, ordre, vocabulaire HSK 3) et combinez-les. La simplicité bien prononcée > La complexité mal prononcée.

\end{savaistu}

\courssub{Lexique récapitulatif de la leçon (Vie culturelle)}

\begin{center}
\begin{tabularx}{\linewidth}{|X|X|X|X|}
\hline
\rowcolor{popblueL} \textbf{Caractères} & \textbf{Pinyin} & \textbf{Nature} & \textbf{Traduction} \\ \hline
诺贡多节 & Nuògòngduō jié & n. & Fête du Ngondo (Cameroun) \\ \hline
中秋节 & Zhōngqiū jié & n. & Fête de la Mi-Automne (Lune) \\ \hline
春节 & Chūnjié & n. & Nouvel An chinois (Fête du Printemps) \\ \hline
国庆节 & Guóqìng jié & n. & Fête nationale \\ \hline
赛船 & sàichuán & v./n. & Course de pirogues / bateaux \\ \hline
团圆 & tuányuán & adj./v. & Réunion (familiale), se réunir \\ \hline
仪式感 & yíshǐ gǎn & n. & Sens du rituel, solennité \\ \hline
传承 & chuánchéng & v. & Transmettre (héritage culturel) \\ \hline
创新 & chuàngxīn & v./n. & Innover, innovation \\ \hline
商业化 & shāngyèhuà & n./v. & Commercialisation, commercialiser \\ \hline
认同感 & rèntóng gǎn & n. & Sentiment d'appartenance / d'identité \\ \hline
代沟 & dàigōu & n. & Fossé des générations \\ \hline
非遗 & fēiyí & n. & Patrimoine culturel immatériel (abbr.) \\ \hline
饺子 & jiǎozi & n. & Raviolis (jiaozi) \\ \hline
月饼 & yuèbǐng & n. & Gâteau de lune \\ \hline
红包 & hóngbāo & n. & Enveloppe rouge (argent) \\ \hline
舞狮 & wǔshī & v./n. & Danse du lion \\ \hline
贴春联 & tiē chūnlián & v.+o. & Afficher les couplets du Printemps \\ \hline
烤鱼 & kǎoyú & n. & Poisson grillé \\ \hline
瓦里河 & Wǎlǐ hé & n. & Fleuve Wouri (Douala) \\ \hline
萨瓦人 & Sàwǎ rén & n. & Peuple Sawa (côte camerounaise) \\ \hline
\end{tabularx}
\end{center}

\courssub{Exercices de consolidation (Entraînement autonome)}

\begin{exercice}[Préparation fiche bristol — 5 mots-clés]
Pour chacun des sujets ci-dessous, écrivez les \textbf{5 mots-clés maximum} (en caractères) que vous mettriez sur votre fiche bristol pour faire un exposé de 2 minutes.
\begin{enumerate}
    \item La fête du Ngondo (Cameroun).
    \item La Fête de la Mi-Automne / Lune (Chine) — \zh{中秋节}.
    \item Le Nouvel An chinois (Chine) — \zh{春节}.
    \item La Fête nationale du Cameroun (20 mai) — \zh{喀麦隆国庆节}.
\end{enumerate}
\end{exercice}

\begin{corrige}[Exercice : Préparation fiche bristol]
\begin{enumerate}
    \item \zh{诺贡多、赛船、瓦里河、团结、烤鱼} (Ngondo, course pirogues, fleuve Wouri, unité, poisson grillé)
    \item \zh{中秋节、月饼、赏月、团圆、嫦娥} (Mi-Automne, gâteaux de lune, admirer lune, réunion, Chang'e)
    \item \zh{春节、饺子、红包、舞狮、贴春联} (Nouvel An, raviolis, enveloppes rouges, danse lion, afficher couplets)
    \item \zh{五二零、游行、总统、雅温得、独立} (20 mai, défilé, président, Yaoundé, indépendance)
\end{enumerate}
\emph{Rappel :} Pas de particules (\zh{的}, \zh{了}), pas de verbes conjugués, juste des noms / verbes d'action / lieux / concepts forts.
\end{corrige}

\begin{exercice}[Transformation : Écrit $\to$ Oral]
Transformez ces phrases « style écrit / dissertation » en phrases « style oral / exposé » (plus courtes, sujet explicite, marqueurs de parole).
\begin{enumerate}
    \item \zh{由于诺贡多节具有极高的文化价值，所以它被列入非遗名录。}
    \item \zh{我认为应该加强对传统节日的保护和传承工作。}
    \item \zh{现代节日的商业化倾向日益严重，导致文化内涵淡化。}
\end{enumerate}
\end{exercice}

\begin{corrige}[Exercice : Transformation]
\begin{enumerate}
    \item \zh{诺贡多节文化价值很高，所以入选了非遗名录。} (Nuògòngduō jié wénhuà jiàzhí hěn gāo, suǒyǐ rùxuǎn le fēiyí mínglù.) — Suppression \zh{由于/所以} lourd, phrase courte.
    \item \zh{我觉得啊，咱们得好好保护传统节日，传下去。} (Wǒ juéde a, zánmen děi hǎohao bǎohù chuántǒng jiérì, chuán xiàqù.) — Ajout \zh{我觉得啊、咱们、好好}, registre oral.
    \item \zh{现在的现代节日太商业了，文化味儿越来越淡。} (Xiànzài de xiàndài jiérì tài shāngyè le, wénhuà wèir yuè lái yuè dàn.) — Vocabulaire oral \zh{太…了、味儿、淡}.
\end{enumerate}
\end{corrige}

\begin{attention}[Derniers rappels avant l'évaluation sommative]
\begin{itemize}
    \item \textbf{Tons :} Révisez les paires tonales difficiles : \zh{节日} (jié rì, 2-4) vs \zh{结实} (jiē shi, 1-2) ; \zh{传统} (chuán tǒng, 2-3) vs \zh{船通} (chuán tōng, 2-1).
    \item \textbf{Classificateurs :} \zh{一个节日} (yí gè jiérì), \zh{一场赛船} (yī chǎng sàichuán), \zh{一位艺术家} (yī wèi yìshùjiā).
    \item \textbf{Connecteurs logiques oraux :} Préférez \zh{然后} (ránhòu, ensuite), \zh{还有} (háiyǒu, aussi), \zh{不过} (búguò, mais) à \zh{此外} (cǐwài), \zh{因此} (yīncǐ), \zh{然而} (rán'ér) (registre écrit).
\end{itemize}
\end{attention}

\newpage

\coursec{Activité d'intégration}

\courssub{Objectifs de l'activité}

À la fin de cette activité d'intégration, tu seras capable de :

\begin{itemize}
\item présenter oralement, en chinois, un aspect de la vie culturelle du Cameroun ou de la Chine (fête, musique, cuisine, cinéma) en 6 à 8 phrases complètes ;
\item utiliser les formules fonctionnelles pour \cle{présenter} (\zh{我给大家介绍……}), \cle{donner son avis} (\zh{我觉得……}) et \cle{relancer} le dialogue (\zh{你觉得……？为什么？}) ;
\item poser des questions simples à un camarade et réagir à ses propos (\zh{我不太同意……}) ;
\item prononcer correctement les tons dans les mots-clés du lexique culturel ;
\item écouter attentivement un exposé et l'évaluer à l'aide d'une grille simple (tons, formules, fluidité).
\end{itemize}

\courssub{Situation}

Dans le cadre de la \cle{Semaine culturelle} du lycée de Nkol-Eton (Yaoundé), ton club de chinois organise un \textbf{forum culturel Cameroun--Chine}. Des élèves de Première A tiennent des stands : fêtes, musique, cuisine, cinéma. Des visiteurs (autres élèves, professeurs) circulent de stand en stand. Chaque exposant doit présenter son thème en chinois, répondre aux questions et convaincre les visiteurs. À la fin, la classe élit le \zh{最好的文化大使} (\emph{zuì hǎo de wénhuà dàshǐ}, « le meilleur ambassadeur culturel »).

Voici l'affiche officielle du forum, rédigée par le club :

\begin{textebox}[L'affiche du forum culturel]
喀麦隆—中国文化论坛 \\
\emph{Kāmàilóng — Zhōngguó wénhuà lùntán} \\
« Forum culturel Cameroun--Chine » \\[4pt]
时间：星期五下午两点到五点 \\
\emph{Shíjiān : xīngqīwǔ xiàwǔ liǎng diǎn dào wǔ diǎn} \\
« Horaire : vendredi, de 14 h à 17 h » \\[4pt]
地点：学校大厅 \\
\emph{Dìdiǎn : xuéxiào dàtīng} \\
« Lieu : le grand hall de l'école » \\[4pt]
内容：节日、音乐、美食、电影 \\
\emph{Nèiróng : jiérì, yīnyuè, měishí, diànyǐng} \\
« Contenu : fêtes, musique, gastronomie, cinéma » \\[4pt]
欢迎大家来听、来问、来说！ \\
\emph{Huānyíng dàjiā lái tīng, lái wèn, lái shuō !} \\
« Bienvenue à tous : venez écouter, questionner et parler ! »
\end{textebox}

\courssub{Consignes}

\begin{enumerate}
\item \textbf{Préparation du stand (15 min).} Par groupes de 3 ou 4, choisissez un thème : \zh{节日} (fêtes), \zh{音乐} (musique), \zh{美食} (cuisine) ou \zh{电影} (cinéma), côté Cameroun \emph{ou} côté Chine. Rédigez ensemble un mini-exposé de 6 à 8 phrases en respectant la structure vue en classe : salutation → présentation du sujet → trois informations → avis personnel → conclusion et remerciement.
\item \textbf{Entraînement des tons (5 min).} Relisez votre exposé à voix haute. Soulignez les mots difficiles pour les tons (\zh{音乐} yīnyuè, \zh{电影} diànyǐng, \zh{好吃} hǎochī) et répétez-les trois fois.
\item \textbf{Première tournée (15 min).} La moitié de la classe tient les stands, l'autre moitié circule. Chaque visiteur écoute l'exposé, pose \textbf{au moins deux questions} avec \zh{你觉得……？} ou \zh{为什么……？}, puis donne son avis avec \zh{我觉得……} ou \zh{我不太同意……}.
\item \textbf{Fiche d'écoute.} Pendant la visite, remplissez la fiche : l'exposant a-t-il bien prononcé les tons ? a-t-il utilisé au moins deux formules fonctionnelles ? son exposé était-il fluide ? Notez de 1 à 3 pour chaque critère.
\item \textbf{Inversion des rôles (15 min).} Les visiteurs deviennent exposants et inversement.
\item \textbf{Élection (5 min).} La classe vote pour le \zh{最好的文化大使} en justifiant son choix en une phrase : \zh{我觉得……最好，因为……}.
\item \textbf{Grand débat collectif (10 min).} L'enseignant lance la question finale : \zh{中国的文化有意思还是喀麦隆的文化有意思？} (\emph{Zhōngguó de wénhuà yǒu yìsi háishi Kāmàilóng de wénhuà yǒu yìsi ?} « La culture chinoise ou la culture camerounaise est-elle plus intéressante ? »). Chacun donne son avis en une ou deux phrases avec \zh{我觉得……，因为……}.
\end{enumerate}

\courssub{Ressources à mobiliser}

\begin{center}
\begin{tabularx}{\linewidth}{|X|X|X|X|}
\hline
\rowcolor{popblueL} Caractères & Pinyin & Nature & Traduction \\
\hline
文化 & wénhuà & nom & culture \\
\hline
节日 & jiérì & nom & fête, festival \\
\hline
音乐 & yīnyuè & nom & musique \\
\hline
美食 & měishí & nom & gastronomie, bons plats \\
\hline
电影 & diànyǐng & nom & film, cinéma \\
\hline
有意思 & yǒu yìsi & adj. & intéressant, amusant \\
\hline
介绍 & jièshào & verbe & présenter \\
\hline
同意 & tóngyì & verbe & être d'accord \\
\hline
\end{tabularx}
\end{center}

\begin{aretenir}[Boîte à outils du forum]
\begin{itemize}
\item \textbf{Présenter} : \zh{大家好！我给大家介绍……} (\emph{Dàjiā hǎo ! Wǒ gěi dàjiā jièshào……} « Bonjour à tous ! Je vous présente… »).
\item \textbf{Donner son avis} : \zh{我觉得……很有意思。} (\emph{Wǒ juéde…… hěn yǒu yìsi.} « Je trouve que… c'est très intéressant. »).
\item \textbf{Relancer} : \zh{你觉得……怎么样？} (\emph{Nǐ juéde…… zěnmeyàng ?} « Que penses-tu de… ? ») ; \zh{为什么？} (\emph{Wèishénme ?} « Pourquoi ? »).
\item \textbf{Nuancer} : \zh{我不太同意……} (\emph{Wǒ bú tài tóngyì……} « Je ne suis pas tout à fait d'accord… »).
\item \textbf{Conclure} : \zh{谢谢大家！} (\emph{Xièxie dàjiā !} « Merci à tous ! »).
\end{itemize}
\end{aretenir}

\begin{attention}[Pièges de tons à surveiller]
Ne confonds pas \zh{问} (\emph{wèn}, 4\textsuperscript{e} ton, « demander ») et \zh{文} (\emph{wén}, 2\textsuperscript{e} ton, dans \zh{文化}). Attention aussi à \zh{买} (\emph{mǎi}, 3\textsuperscript{e} ton, « acheter ») et \zh{卖} (\emph{mài}, 4\textsuperscript{e} ton, « vendre ») si tu parles du marché. Dans \zh{有意思}, le troisième caractère \zh{思} se prononce au ton neutre : \emph{yìsi}.
\end{attention}

\courssub{Proposition de production et corrigé commenté}

\begin{exoresolu}[Mini-exposé modèle : stand « cuisine camerounaise »]
Présente le stand « gastronomie du Cameroun » en 6 à 8 phrases, avec salutation, trois informations, un avis et une conclusion.
\tcblower
\textbf{Production modèle :} \\[4pt]
大家好！我给大家介绍喀麦隆的美食。 \\
\emph{Dàjiā hǎo ! Wǒ gěi dàjiā jièshào Kāmàilóng de měishí.} \\
« Bonjour à tous ! Je vous présente la gastronomie du Cameroun. » \\[4pt]
喀麦隆有很多好吃的菜，比如恩多莱。 \\
\emph{Kāmàilóng yǒu hěn duō hǎochī de cài, bǐrú ndolé.} \\
« Le Cameroun a beaucoup de bons plats, par exemple le ndolé. » \\[4pt]
我们也喜欢吃烤鱼和木薯。 \\
\emph{Wǒmen yě xǐhuan chī kǎo yú hé mùshǔ.} \\
« Nous aimons aussi manger le poisson braisé et le manioc. » \\[4pt]
吃饭以前，我们常常说“请吃”。 \\
\emph{Chī fàn yǐqián, wǒmen chángcháng shuō “qǐng chī”.} \\
« Avant de manger, nous disons souvent “bon appétit”. » \\[4pt]
我觉得喀麦隆的美食又好吃又有意思。 \\
\emph{Wǒ juéde Kāmàilóng de měishí yòu hǎochī yòu yǒu yìsi.} \\
« Je trouve la gastronomie camerounaise à la fois délicieuse et intéressante. » \\[4pt]
欢迎大家来尝一尝！谢谢大家！ \\
\emph{Huānyíng dàjiā lái cháng yi cháng ! Xièxie dàjiā !} \\
« Bienvenue à tous pour goûter ! Merci à tous ! » \\[6pt]
\textbf{Commentaire des choix de langue :} l'exposé s'ouvre sur la salutation rituelle \zh{大家好}, indispensable à l'oral. La formule de présentation \zh{我给大家介绍……} annonce clairement le sujet. Les trois informations utilisent des structures simples du niveau HSK 2 : \zh{有} pour l'existence, \zh{比如} (« par exemple ») pour illustrer, \zh{也} pour ajouter une idée, et la structure temporelle \zh{……以前} (« avant de… »). L'avis personnel mobilise la structure parallèle \zh{又……又……} (« à la fois… et… »), très appréciée à l'oral. La conclusion combine une invitation (\zh{尝一尝}, verbe redoublé qui adoucit le ton, « goûter un peu ») et le remerciement final. Les mots à tons délicats (\zh{美食} měishí, \zh{好吃} hǎochī) sont placés en position forte dans la phrase pour être bien articulés.
\end{exoresolu}

\courssub{Bilan de l'activité}

À l'issue du forum, vérifie que tu sais :

\begin{itemize}
\item construire un mini-exposé oral complet en cinq étapes (salutation, présentation, informations, avis, conclusion) ;
\item employer sans hésiter les formules \zh{我给大家介绍……}, \zh{我觉得……}, \zh{你觉得……怎么样？}, \zh{我不太同意……} ;
\item poser des questions ouvertes avec \zh{为什么} et réagir à la réponse d'un camarade ;
\item prononcer correctement les tons des mots-clés du lexique culturel (\zh{文化}, \zh{音乐}, \zh{电影}, \zh{美食}, \zh{有意思}) ;
\item écouter un exposé en repérant les informations principales et évaluer la production d'un pair selon des critères précis.
\end{itemize}

\begin{methode}[Pour aller plus loin]
Chez toi, enregistre ton mini-exposé avec ton téléphone. Réécoute-toi en te concentrant uniquement sur les tons : note les mots qui « sonnent faux », corrige-les avec le pinyin, puis réenregistre. Compare les deux versions : la seconde doit être plus fluide et plus juste. Tu peux aussi préparer une version « Chine » de ton exposé pour le prochain forum.
\end{methode}

\newpage

\coursec{Méthodes et exercices résolus}

\coursec{Leçon 21 — Expression orale : vie culturelle au Cameroun et en Chine}

\courssub{Lexique de la vie culturelle}

\begin{textebox}[Mise en situation : une conversation au foyer des élèves]
\zh{— 小李，这个周末有什么活动？} \emph{— Xiǎo Lǐ, zhège zhōumò yǒu shénme huódòng?} « Petit Li, y a-t-il des activités ce week-end ? » \\
\zh{— 我们学校有文化节，还有马科萨音乐会。你想去吗？} \emph{— Wǒmen xuéxiào yǒu wénhuà jié, hái yǒu Mǎkēsà yīnyuèhuì. Nǐ xiǎng qù ma?} « Notre école organise une fête culturelle, et il y a un concert de makossa. Tu veux y aller ? » \\
\zh{— 好啊！我最喜欢看非洲舞蹈了。} \emph{— Hǎo a! Wǒ zuì xǐhuan kàn Fēizhōu wǔdǎo le.} « Super ! J'adore regarder la danse africaine. »
\end{textebox}

\begin{center}
\begin{tabularx}{\linewidth}{|X|X|X|X|}
\hline
\rowcolor{popblueL} \textbf{Caractères} & \textbf{Pinyin} & \textbf{Nature} & \textbf{Traduction} \\
\hline
文化 & wénhuà & n. & culture \\
节日 & jiérì & n. & fête, festival \\
青年节 & Qīngnián Jié & n. & fête de la jeunesse (11 fév.) \\
春节 & Chūn Jié & n. & Fête du Printemps (Nouvel An chinois) \\
国庆节 & Guóqìng Jié & n. & Fête nationale (1er oct. Chine) \\
圣诞节 & Shèngdàn Jié & n. & Noël \\
音乐 & yīnyuè & n. & musique \\
马科萨 & Mǎkēsà & n. & makossa (genre musical camerounais) \\
舞蹈 & wǔdǎo & n. & danse \\
电影 & diànyǐng & n. & film, cinéma \\
足球 & zúqiú & n. & football \\
比赛 & bǐsài & n./v. & match, compétition / concourir \\
音乐会 & yīnyuèhuì & n. & concert \\
展览 & zhǎnlǎn & n. & exposition \\
博物馆 & bówùguǎn & n. & musée \\
\hline
\rowcolor{poporangeL} \textbf{Activités \& Verbes} & & & \\
\hline
唱歌 & chàng gē & v. & chanter \\
跳舞 & tiào wǔ & v. & danser \\
看 & kàn & v. & regarder, voir \\
参加 & cānjiā & v. & participer à \\
介绍 & jièshào & v. & présenter, introduire \\
邀请 & yāoqǐng & v. & inviter \\
\hline
\rowcolor{popgreenL} \textbf{Adjectifs \& Opinions} & & & \\
\hline
有意思 & yǒu yìsi & adj. & intéressant \\
精彩 & jīngcǎi & adj. & splendide, palpitant \\
热闹 & rènào & adj. & animé, bruyant (fête) \\
传统 & chuántǒng & adj. & traditionnel \\
现代 & xiàndài & adj. & moderne \\
重要 & zhòngyào & adj. & important \\
\hline
\rowcolor{poppurpleL} \textbf{Connecteurs \& Mots-outils} & & & \\
\hline
首先 & shǒuxiān & adv. & d'abord \\
然后 & ránhòu & adv. & ensuite \\
最后 & zuìhòu & adv. & enfin \\
因为 & yīnwèi & conj. & parce que \\
所以 & suǒyǐ & conj. & donc, par conséquent \\
可是 & kěshì & conj. & mais \\
不过 & bùguò & conj. & cependant \\
觉得 & juéde & v. & trouver, penser que \\
认为 & rènwéi & v. & considérer que \\
\hline
\end{tabularx}
\end{center}

\begin{aretenir}[Mots-clés à mémoriser prioritairement]
\zh{文化、节日、音乐、舞蹈、足球、比赛、音乐会、参加、介绍、有意思、精彩、首先、然后、因为、所以、可是、觉得、认为}.
\end{aretenir}

\courssub{Dialogue modèle 1 : présenter une fête}

\begin{textebox}[Dialogue : présenter la Fête de la Jeunesse au Cameroun]
\zh{A : 大家好！今天我要给大家介绍喀麦隆的青年节。} \emph{Dàjiā hǎo! Jīntiān wǒ yào gěi dàjiā jièshào Kāmàilóng de Qīngnián Jié.} \\
\zh{B : 青年节是什么时候？} \emph{Qīngnián Jié shì shénme shíhou?} \\
\zh{A : 青年节在二月十一日。首先，学生们不上课。然后，大家去体育场唱歌、跳舞、踢足球。最后，晚上有音乐会。} \emph{Qīngnián Jié zài èryuè shíyī rì. Shǒuxiān, xuéshengmen bù shàngkè. Ránhòu, dàjiā qù tǐyùchǎng chàng gē, tiào wǔ, tī zúqiú. Zuìhòu, wǎnshang yǒu yīnyuèhuì.} \\
\zh{B : 听起来很有意思！你觉得这个节日重要吗？} \emph{Tīng qǐlái hěn yǒu yìsi! Nǐ juéde zhège jiérì zhòngyào ma?} \\
\zh{A : 我觉得很重要，因为年轻人是国家的未来。谢谢大家！} \emph{Wǒ juéde hěn zhòngyào, yīnwèi niánqīngrén shì guójiā de wèilái. Xièxie dàjiā!}
\end{textebox}

\courssub{Formules fonctionnelles : présenter, opiner, relancer}

\begin{propriete}[Structure d'un exposé oral simple : 3 temps]
\textbf{1. Ouverture (开场白)} \\
\zh{大家好！今天我要给大家介绍……} \emph{Dàjiā hǎo! Jīntiān wǒ yào gěi dàjiā jièshào...} \\
\zh{各位同学好，我的题目是……} \emph{Gèwèi tóngxué hǎo, wǒ de tímù shì...} « Chers camarades, mon sujet est... »

\textbf{2. Développement (主体内容) — 3 à 5 phrases avec connecteurs} \\
\zh{首先，……} \emph{Shǒuxiān...} « D'abord... » \\
\zh{然后，……} \emph{Ránhòu...} « Ensuite... » \\
\zh{最后，……} \emph{Zuìhòu...} « Enfin... » \\
\emph{Ordre canonique :} Sujet + Temps + Lieu + Verbe + Complément. \\
\zh{青年节在二月十一日。} (Temps avant verbe) \checkmark

\textbf{3. Conclusion (结束语) — Avis + Remerciement} \\
\zh{我觉得……很有意思 / 很重要。} \emph{Wǒ juéde... hěn yǒu yìsi / hěn zhòngyào.} \\
\zh{我认为……很有意义。} \emph{Wǒ rènwéi... hěn yǒu yìyì.} \\
\zh{谢谢大家！ / 谢谢聆听！} \emph{Xièxie dàjiā! / Xièxie língtīng!}
\end{propriete}

\begin{propriete}[Formules pour le débat : Opiner, Réagir, Relancer]
\begin{center}
\begin{tabularx}{\linewidth}{|X|X|X|}
\hline
\rowcolor{popblueL} \textbf{Fonction} & \textbf{Formules chinoises (Pinyin)} & \textbf{Traduction} \\
\hline
\zh{给出观点} & \zh{我觉得……} \emph{Wǒ juéde...} & Je pense que... / Je trouve que... \\
& \zh{我认为……} \emph{Wǒ rènwéi...} & Je considère que... (plus formel) \\
& \zh{对我来说，……} \emph{Duì wǒ lái shuō...} & Pour moi... \\
\hline
\zh{买} \emph{mǎi} & \zh{我同意你的看法。} \emph{Wǒ tóngyì nǐ de kànfǎ.} & Je suis d'accord avec ton avis. \\
\zh{赞同/反对} & \zh{我不同意，因为……} \emph{Wǒ bù tóngyì, yīnwèi...} & Je ne suis pas d'accord, car... \\
& \zh{你说得对，可是……} \emph{Nǐ shuō de duì, kěshì...} & Tu as raison, mais... (désaccord poli) \\
& \zh{我不太赞成，因为……} \emph{Wǒ bù tài zànchéng, yīnwèi...} & Je ne suis pas très pour, car... \\
\hline
\zh{追问/邀请对方} & \zh{你呢？} \emph{Nǐ ne?} & Et toi ? \\
& \zh{你觉得怎么样？} \emph{Nǐ juéde zěnme yàng?} & Qu'en penses-tu ? \\
& \zh{你为什么这么想？} \emph{Nǐ wèishénme zhème xiǎng?} & Pourquoi penses-tu cela ? \\
& \zh{还有其他的看法吗？} \emph{Hái yǒu qítā de kànfǎ ma?} & D'autres avis ? \\
\hline
\zh{论证} & \zh{因为……所以……} \emph{Yīnwèi... suǒyǐ...} & Parce que... donc... (les deux mots s'utilisent ensemble) \\
\hline
\end{tabularx}
\end{center}
\end{propriete}

\begin{exoresolu}[Exposé guidé : La Fête de la Jeunesse]
\textbf{Énoncé.} Préparez un mini-exposé de 5 phrases pour présenter la fête de la jeunesse (\zh{青年节}) à un correspondant chinois. Contraintes : salutation + annonce, 2 connecteurs (\zh{首先}, \zh{然后}), activités (chanter, danser, football), avis personnel, remerciement.
\tcblower
\textbf{Solution détaillée.}
\begin{enumerate}
\item \textbf{Ouverture} : \zh{大家好！今天我要给大家介绍喀麦隆的青年节。} \emph{Dàjiā hǎo! Jīntiān wǒ yào gěi dàjiā jièshào Kāmàilóng de Qīngnián Jié.}
\item \textbf{Développement} : \zh{首先，青年节在二月十一日。} \emph{Shǒuxiān, Qīngnián Jié zài èryuè shíyī rì.} (Règle : Temps/Lieu avant verbe). \\
\zh{然后，很多年轻人在体育场唱歌、跳舞、踢足球。} \emph{Ránhòu, hěn duō niánqīngrén zài tǐyùchǎng chàng gē, tiào wǔ, tī zúqiú.} (Lieu \zh{在体育场} avant verbe ; énumération verbes avec 、).
\item \textbf{Conclusion} : \zh{我觉得这个节日很有意思，因为大家很开心。谢谢大家！} \emph{Wǒ juéde zhège jiérì hěn yǒu yìsi, yīnwèi dàjiā hěn kāixīn. Xièxie dàjiā!}
\end{enumerate}
\textbf{Analyse} : 6 phrases, structure respectée, connecteurs utilisés, ordre des mots correct (Temps/Lieu avant V), avis justifié par \zh{因为}.
\end{exoresolu}

\begin{exoresolu}[Débat guidé : Musique camerounaise vs Chinoise]
\textbf{Énoncé.} Complétez la réplique de B : exprimer un désaccord poli, justifier avec \zh{因为}, relancer par une question ouverte.
\zh{A : 我觉得中国音乐很好听。你呢？} \emph{Wǒ juéde Zhōngguó yīnyuè hěn hǎotīng. Nǐ ne?} \\
\zh{B : \_\_\_\_\_\_\_\_\_\_\_\_}
\tcblower
\textbf{Proposition de réponse.} \\
\zh{你说得对，可是我更喜欢喀麦隆的马科萨音乐，因为它很有节奏，让人想跳舞。你最喜欢哪种音乐？} \\
\emph{Nǐ shuō de duì, kěshì wǒ gèng xǐhuan Kāmàilóng de Mǎkēsà yīnyuè, yīnwèi tā hěn yǒu jiézòu, ràng rén xiǎng tiào wǔ. Nǐ zuì xǐhuan nǎ zhǒng yīnyuè?} \\
« Tu as raison, mais je préfère le makossa camerounais, car il a du rythme, ça donne envie de danser. Quel genre de musique préfères-tu ? »
\textbf{Points clés} : \zh{更} (encore plus) avant le verbe ; \zh{因为}... (cause) ; \zh{让人想跳舞} (complément de résultat \zh{让} + \zh{想}) ; question en \zh{哪种} (quel type) sans inversion.
\end{exoresolu}

\courssub{Les tons : entraînement ciblé}

\begin{methode}[Règles de prononciation : tons, sandhi et ton neutre]
\begin{enumerate}
\item \textbf{Les 4 tons + ton neutre} : 1\textsuperscript{er} haut plat (\zh{妈} \emph{mā}), 2\textsuperscript{e} montant (\zh{麻} \emph{má}), 3\textsuperscript{e} descendant-montant (\zh{马} \emph{mǎ}), 4\textsuperscript{e} descendant sec (\zh{骂} \emph{mà}), neutre léger/court (\zh{吗} \emph{ma}).
\item \textbf{Sandhi du 3\textsuperscript{e} ton (règle des deux 3\textsuperscript{es} tons)} : 3\textsuperscript{e} + 3\textsuperscript{e} $\rightarrow$ 2\textsuperscript{e} + 3\textsuperscript{e} à l'oral. Ex. : \zh{你好} \emph{ní hǎo}, \zh{很好} \emph{hén hǎo}, \zh{我很喜欢} \emph{wó hén xǐhuan} (trois 3\textsuperscript{es} tons $\rightarrow$ 2-2-3).
\item \textbf{Sandhi de \zh{不} \emph{bù}} : \emph{bù} (4\textsuperscript{e}) $\rightarrow$ \emph{bú} (2\textsuperscript{e}) devant un 4\textsuperscript{e} ton. Ex. : \zh{不是} \emph{bú shì}, \zh{不要} \emph{bú yào}, \zh{不去} \emph{bú qù}. Devant 1, 2, 3, neutre : reste \emph{bù} (\zh{不买} \emph{bù mǎi}).
\item \textbf{Sandhi de \zh{一} \emph{yī}} : seul/chiffre/ordinal = 1\textsuperscript{er} ton (\emph{yī}) ; devant 4\textsuperscript{e} ton = 2\textsuperscript{e} ton (\emph{yí} : \zh{一个} \emph{yí ge}) ; devant 1, 2, 3, neutre = 4\textsuperscript{e} ton (\emph{yì} : \zh{一样} \emph{yìyàng}).
\item \textbf{Ton neutre fréquent} : particules (\zh{了} \emph{le}, \zh{吗} \emph{ma}, \zh{吧} \emph{ba}), suffixes (\zh{子} \emph{zi} : \zh{桌子} \emph{zhuōzi}), 2\textsuperscript{e} syllabe de certains mots (\zh{欢} dans \zh{喜欢} \emph{xǐhuan}, \zh{国} dans \zh{中国} \emph{Zhōngguó} — attention, \zh{国} garde souvent son ton 2 en dictionnaire mais s'allège en flux).
\end{enumerate}
\end{methode}

\begin{popfigure}
\begin{tikzpicture}[scale=0.9, every node/.style={font=\small}]
% Tone contour diagram
\draw[popline, thick] (0,0) -- (6,0) node[right] {Temps};
\draw[popline, thick] (0,0) -- (0,3) node[above] {Hauteur};
% Tone 1
\draw[popblue, very thick] (0.5,2.2) -- (1.5,2.2) node[midway, above] {1\textsuperscript{er} \emph{mā}};
% Tone 2
\draw[poporange, very thick] (2,0.5) .. controls (2.5,1.5) and (3,2.5) .. (3.5,2.5) node[midway, above right] {2\textsuperscript{e} \emph{má}};
% Tone 3
\draw[popgreen, very thick] (4,1.5) .. controls (4.3,0.5) and (4.7,0.5) .. (5,1.5) node[midway, below] {3\textsuperscript{e} \emph{mǎ}};
% Tone 4
\draw[poppink, very thick] (5.5,2.5) -- (6.5,0.5) node[midway, right] {4\textsuperscript{e} \emph{mà}};
% Neutral
\node at (7, 1.2) {Neutre \emph{ma} (court, bas)};
\end{tikzpicture}
\legende{Contours des quatre tons et ton neutre (schématique).}
\end{popfigure}

\begin{exoresolu}[Identification et correction des tons]
\textbf{Énoncé.} (1) Donnez les tons (dictionnaire) de : \zh{文化}, \zh{音乐}, \zh{电影}, \zh{喀麦隆}. (2) Écrivez la prononciation réelle (sandhi appliqué) de : \zh{我很好}, \zh{不想去}, \zh{一个音乐会}.
\tcblower
\textbf{Solution.}
\begin{enumerate}
\item \zh{文化} \emph{wén huà} (2-4) ; \zh{音乐} \emph{yīn yuè} (1-4) ; \zh{电影} \emph{diàn yǐng} (4-3) ; \zh{喀麦隆} \emph{Kā mài lóng} (1-4-2).
\item \zh{我很好} : dictionnaire \emph{wǒ hěn hǎo} (3-3-3) $\rightarrow$ oral \emph{wó hén hǎo} (2-2-3). \\
\zh{不想去} : \zh{不} (4) + \zh{想} (3) + \zh{去} (4). \zh{不} devant \zh{想} (3\textsuperscript{e} ton) $\rightarrow$ reste \emph{bù}. Prononcé \emph{bù xiǎng qù}. (Pas de sandhi \emph{bú} car \zh{想} n'est pas 4\textsuperscript{e} ton). \\
\zh{一个音乐会} : \zh{一} (1) + \zh{个} (4) + \zh{音} (1) + \zh{乐} (4) + \zh{会} (4). \zh{一} devant \zh{个} (4\textsuperscript{e} ton) $\rightarrow$ \emph{yí}. Prononcé \emph{yí ge yīn yuè huì}.
\end{enumerate}
\textbf{Attention} : Le pinyin écrit garde les tons d'origine (citation form). Seule la prononciation change.
\end{exoresolu}

\courssub{Dialogue modèle 2 : donner son avis et débattre}

\begin{textebox}[Dialogue : Débat sur les loisirs du week-end]
\zh{A : 周末你想做什么？} \emph{Zhōumò nǐ xiǎng zuò shénme?} « Que veux-tu faire le week-end ? » \\
\zh{B : 我想去看足球比赛。喀麦隆对阵埃及，很精彩。你呢？} \emph{Wǒ xiǎng qù kàn zúqiú bǐsài. Kāmàilóng duìzhèn Aījí, hěn jīngcǎi. Nǐ ne?} « Je veux aller voir un match de foot. Cameroun contre Égypte, c'est palpitant. Et toi ? » \\
\zh{A : 我不同意。我觉得看电影更舒服，因为天气太热。} \emph{Wǒ bù tóngyì. Wǒ juéde kàn diànyǐng gèng shūfu, yīnwèi tiānqì tài rè.} « Je ne suis pas d'accord. Je trouve le ciné plus confortable, car il fait trop chaud. » \\
\zh{B : 你说得对，可是在体育场气氛很好，大家一起唱歌。} \emph{Nǐ shuō de duì, kěshì zài tǐyùchǎng qìfēn hěn hǎo, dàjiā yìqǐ chàng gē.} « Tu as raison, mais l'ambiance au stade est top, tout le monde chante ensemble. » \\
\zh{A : 好吧，那我们一起去吧！} \emph{Hǎo ba, nà wǒmen yìqǐ qù ba!} « D'accord, alors on y va ensemble ! »
\end{textebox}

\courssub{Jeux de rôle, débat guidé et comparaison culturelle}

\begin{methode}[Préparer et réussir un jeu de rôle (binôme)]
\begin{enumerate}
\item \textbf{Analyser la situation} : Qui ? Où ? Sujet ? Relation (amis, prof/élève, vendeur/client) $\rightarrow$ détermine le registre (familier \zh{你} / poli \zh{您}).
\item \textbf{Structurer en 6--8 répliques} : Salutation $\rightarrow$ Proposition/Question $\rightarrow$ Réaction (avis) $\rightarrow$ Négociation/Détails (heure, lieu) $\rightarrow$ Décision $\rightarrow$ Clôture.
\item \textbf{Mobiliser le lexique du thème} (voir tableau §1) et les formules §3.
\item \textbf{Répéter à voix haute} : soigner les tons (sandhi), la fluidité (pas de pause mot à mot), le contact visuel.
\item \textbf{Clôture naturelle} : \zh{下次再见！} \emph{Xià cì zài jiàn!} / \zh{我们一起去吧！} \emph{Wǒmen yìqǐ qù ba!}
\end{enumerate}
\end{methode}

\begin{exoresolu}[Jeu de rôle : Invitation au concert de Makossa]
\textbf{Énoncé.} A invite B à un concert de makossa à Douala samedi soir. B hésite (loin/cher), A insiste (ambiance), B accepte et propose resto après.
\tcblower
\textbf{Modèle de dialogue (8 répliques).} \\
\zh{A : 小王，星期六晚上杜阿拉有马科萨音乐会，你去吗？} \emph{Xiǎo Wáng, Xīngqīliù wǎnshang Dù'ālā yǒu Mǎkēsà yīnyuèhuì, nǐ qù ma?} \\
\zh{B : 杜阿拉太远了，而且票很贵吧？} \emph{Dù'ālā tài yuǎn le, érqiě piào hěn guì ba?} \\
\zh{A : 可是现场气氛特别好，大家一起跳舞，很有意思！} \emph{Kěshì xiànchǎng qìfēn tèbié hǎo, dàjiā yìqǐ tiào wǔ, hěn yǒu yìsi!} \\
\zh{B : 真的吗？那我考虑一下。几点开始？} \emph{Zhēn de ma? Nà wǒ kǎolǜ yí xià. Jǐ diǎn kāishǐ?} \\
\zh{A : 晚上八点。我可以送你一程。} \emph{Wǎnshang bā diǎn. Wǒ kěyǐ sòng nǐ yí chéng.} \\
\zh{B : 太好了！听完音乐会，我们去吃烤鱼吧，听说那里很好吃。} \emph{Tài hǎo le! Tīng wán yīnyuèhuì, wǒmen qù chī kǎo yú ba, tīngshuō nàli hěn hǎochī.} \\
\zh{A : 好主意！那就这么说定了。} \emph{Hǎo zhǔyì! Nà jiù zhème shuō dìng le.} \\
\zh{B : 好，周六见！} \emph{Hǎo, Xīngqīliù jiàn!}
\textbf{Vérif. grammaire} : \zh{太……了} (excès) ; \zh{而且} (de plus) ; \zh{特别} (surtout) ; \zh{考虑一下} (réfléchir un peu) ; \zh{送...一程} (raccompagner) ; \zh{听完} (après avoir fini d'écouter) ; \zh{好吃} (bon à manger) ; \zh{这么说定了} (c'est dit/entendu).
\end{exoresolu}

\begin{propriete}[Comparer avec \zh{比} et \zh{没有} : vie culturelle Chine / Cameroun]
\textbf{1. Supériorité (A > B) :} \zh{A + 比 + B + Adj.} \textbf{Jamais de \zh{很} avant l'adjectif.}
\zh{中国的春节比喀麦隆的圣诞节长。} \emph{Zhōngguó de Chūn Jié bǐ Kāmàilóng de Shèngdàn Jié cháng.} (Le Nouvel An chinois est plus long que Noël au Cameroun.)
\zh{足球比篮球有意思得多。} \emph{Zúqiú bǐ lánqiú yǒu yìsi de duō.} (Le foot est bien plus intéressant que le basket.)

\textbf{2. Infériorité / Égalité (A < B ou A = B) :} \zh{A + 没有 + B + Adj.} (A n'est pas aussi Adj que B).
\zh{喀麦隆的冬天没有中国的冬天冷。} \emph{Kāmàilóng de dōngtiān méiyǒu Zhōngguó de dōngtiān lěng.} (L'hiver camerounais n'est pas aussi froid que l'hiver chinois.)
\zh{这部电影没有那部有意思。} \emph{Zhè bù diànyǐng méiyǒu nà bù yǒu yìsi.} (Ce film n'est pas aussi intéressant que celui-là.)

\textbf{3. Nuances :} \zh{一点儿} (un peu), \zh{得多} (beaucoup), \zh{多了} (beaucoup, oral).
\zh{北京的鸭子比雅温得的鸡好吃一点儿。} (Le canard de Pékin est un peu meilleur que le poulet de Yaoundé.)
\end{propriete}

\begin{savaistu}[Focus culture : Fêtes nationales et jeunesse]
\textbf{Cameroun} : \zh{青年节} (11 février) — Défilés, sports, concerts, discours du Président. Journée chômée, forte mobilisation scolaire.
\textbf{Chine} : \zh{青年节} (4 mai) — Commémore le Mouvement du 4 mai 1919 (patriotisme, science, démocratie). Activités : bénévolat, conférences, spectacles. Pas férié pour tous (demi-journée pour les 14-28 ans).
\textbf{Point commun} : La jeunesse est célébrée comme avenir de la nation.
\textbf{Différence} : Au Cameroun, fête festive et sportive ; en Chine, mémoire historique et engagement civique.
\end{savaistu}

\begin{exoresolu}[Comparaison écrite et orale]
\textbf{Énoncé.} (1) Traduisez : « Au Cameroun, le football est plus populaire que le basket-ball. » (2) Corrigez : \zh{喀麦隆的音乐比中国的音乐很有节奏。}
\tcblower
\textbf{Solution.}
\begin{enumerate}
\item \zh{在喀麦隆，足球比篮球受欢迎。} \emph{Zài Kāmàilóng, zúqiú bǐ lánqiú shòu huānyíng.} (Lieu en tête ; \zh{受欢迎} « être populaire » ; pas de \zh{很}).
\item \textbf{Erreur} : \zh{很} interdit après \zh{比} + Adj. \textbf{Correct} : \zh{喀麦隆的音乐比中国的音乐有节奏。} \emph{Kāmàilóng de yīnyuè bǐ Zhōngguó de yīnyuè yǒu jiézòu.} Ou : \zh{喀麦隆的音乐比中国的音乐有节奏得多。} (Bien plus rythmé).
\end{enumerate}
\end{exoresolu}

\begin{attention}[Les 5 erreurs fatales à l'oral (Bac / Contrôle continu)]
\begin{enumerate}
\item \textbf{Tons fantômes} : Dire \emph{bù shì} au lieu de \emph{bú shì} (\zh{不是}) ; dire \emph{wǒ hěn hǎo} au lieu de \emph{wó hén hǎo} (\zh{我很好}). $\rightarrow$ \cle{Solution} : Répéter les paires minimales (\zh{买}/\zh{卖}, \zh{问}/\zh{闻}) et les phrases-types (\zh{我不去}, \zh{一个}) chaque jour.
\item \textbf{Calque \zh{比……很……}} : \zh{春节比圣诞节很长} $\times$. $\rightarrow$ \cle{Règle} : \zh{比} exprime déjà le comparatif ; l'adjectif est nu.
\item \textbf{Désordre Temps/Lieu} : \zh{我去音乐会星期六} $\times$. $\rightarrow$ \cle{Règle} : Sujet + \textbf{Temps} + \textbf{Lieu} + Verbe. \zh{我星期六去音乐会} \checkmark.
\item \textbf{Classificateur fantôme} : \zh{一音乐会}, \zh{这电影} $\times$. $\rightarrow$ \cle{Règle} : Nombre/Démonstratif + \textbf{Classif.} + Nom. \zh{一场音乐会} / \zh{一张票} / \zh{这部电影} \checkmark.
\item \textbf{Caractères « jumeaux » confondus} : \zh{乐} (yuè/lè) vs \zh{东} (dōng) ; \zh{节} (jié) vs \zh{书} (shū) ; \zh{文} (wén) vs \zh{又} (yòu) ; \zh{舞} (wǔ) vs \zh{无} (wú). $\rightarrow$ \cle{Entraînement} : Écrire 5 fois chaque caractère en respectant l'ordre des traits (haut$\rightarrow$bas, gauche$\rightarrow$droite, horizontal$\rightarrow$vertical).
\end{enumerate}
\end{attention}

\courssub{Exercices d'entraînement (Travail personnel / Devoirs)}

\begin{exercice}[Exercice 1 : Vocabulaire et Classificateurs]
Complétez avec le bon classificateur : \zh{个, 场, 张, 部, 首, 次}.
\begin{enumerate}
\item \zh{一\_\_\_\_音乐会} (un concert) \quad 2. \zh{三\_\_\_\_电影} (trois films) \quad 3. \zh{两\_\_\_\_票} (deux billets)
\item \zh{一\_\_\_\_比赛} (un match) \quad 5. \zh{几\_\_\_\_歌} (quelques chansons) \quad 6. \zh{下\_\_\_\_} (une fois / la prochaine fois)
\end{enumerate}
\end{exercice}

\begin{exercice}[Exercice 2 : Grammaire — Comparatif \zh{比} / \zh{没有}]
Traduisez en chinois (caractères + pinyin).
\begin{enumerate}
\item Le makossa est plus rythmé que la musique pop.
\item L'hiver à Yaoundé n'est pas aussi froid qu'à Harbin (哈尔滨 \emph{Hā'ěrbīn}).
\item Ce concert est un peu plus intéressant que celui d'hier.
\item Pour moi, danser est plus fatigant (\zh{累} \emph{lèi}) que chanter.
\end{enumerate}
\end{exercice}

\begin{exercice}[Exercice 3 : Expression orale — Préparation d'exposé]
Préparez un exposé de 1 minute (environ 8 phrases) sur l'un des sujets suivants. Écrivez le texte complet (caractères + pinyin) en respectant la structure : Ouverture $\rightarrow$ Développement (3 connecteurs min.) $\rightarrow$ Conclusion (avis + \zh{因为}).
\begin{enumerate}
\item \zh{介绍喀麦隆的国庆节 (5月20日)} (Présenter la Fête nationale du Cameroun, 20 mai).
\item \zh{比较中国春节和喀麦隆圣诞节} (Comparer Nouvel An chinois et Noël camerounais).
\item \zh{我的周末文化生活} (Ma vie culturelle du week-end).
\end{enumerate}
\end{exercice}

\begin{exercice}[Exercice 4 : Jeu de rôle — Débat]
En binôme. Sujet : \zh{“看现场音乐会” vs “在家看直播”} (Concert sur place vs Regarder le direct à la maison).
A : Préfère le direct (confort, gratuit, voir les détails).
B : Préfère le sur place (ambiance, chanter ensemble, rencontrer des gens).
\textbf{Contraintes} : 6 répliques chacun. Utiliser : \zh{我觉得}, \zh{可是/不过}, \zh{因为...所以...}, \zh{你觉得怎么样？}, \zh{比}, \zh{没有}. Enregistrez-vous ou jouez devant la classe.
\end{exercice}

\begin{corrige}[Exercice 1]
1. \zh{一场音乐会} \emph{yí chǎng yīnyuèhuì} \quad 2. \zh{三部电影} \emph{sān bù diànyǐng} \quad 3. \zh{两张票} \emph{liǎng zhāng piào} \\
4. \zh{一场比赛} \emph{yí chǎng bǐsài} \quad 5. \zh{几首歌} \emph{jǐ shǒu gē} \quad 6. \zh{下次} \emph{xìa cì}
\end{corrige}

\begin{corrige}[Exercice 2]
1. \zh{马科萨音乐比流行音乐有节奏。} \emph{Mǎkēsà yīnyuè bǐ liúxíng yīnyuè yǒu jiézòu.} \\
2. \zh{雅温得的冬天没有哈尔滨的冬天冷。} \emph{Yǎwēndé de dōngtiān méiyǒu Hā'ěrbīn de dōngtiān lěng.} \\
3. \zh{这场音乐会比昨天的有意思一点儿。} \emph{Zhè chǎng yīnyuèhuì bǐ zuótiān de yǒu yìsi yìdiǎnr.} \\
4. \zh{对我来说，跳舞比唱歌累。} \emph{Duì wǒ lái shuō, tiào wǔ bǐ chàng gē lèi.}
\end{corrige}

\begin{corrige}[Exercice 3 — Modèle pour Sujet 1 : Fête nationale 20 mai]
\zh{大家好！今天我要给大家介绍喀麦隆的国庆节。} \\
\zh{首先，国庆节在五月二十日。} \\
\zh{然后，早上有盛大的阅兵式，军人、警察、学生都参加。} \\
\zh{接着，下午有文艺表演，人们唱歌、跳舞，很热闹。} \\
\zh{最后，晚上有烟火，全家人一起看。} \\
\zh{我觉得这个节日很重要，因为它象征国家的团结和独立。} \\
\zh{谢谢大家！}
\end{corrige}

\begin{corrige}[Exercice 4 — Pistes pour le débat]
\zh{A : 我觉得在家看直播比去现场方便，因为不用买票，也不挤。而且电视画面清楚，细节都能看见。} \\
\zh{B : 你说得对，可是现场的气氛没有电视那么冷淡。大家一起唱歌、挥手，很有意思。你不觉得现场更精彩吗？} \\
\zh{A : 可是票很贵，而且回家很晚，很累。我觉得性价比不高。} \\
\zh{B : 没关系，体验比价格重要。下次有演唱会，我们一起去现场吧！} \\
\zh{A : 好，下次再说吧。}
\end{corrige}

\newpage

\coursec{Exercices}

\courssub{Je vérifie mes connaissances}

\begin{exercice}[QCM : la vie culturelle]
Pour chaque question, entoure la bonne réponse.

1. \zh{节日} (jiérì) signifie :
\begin{itemize}
\item A. la musique \quad B. la fête \quad C. le film \quad D. la danse
\end{itemize}

2. Pour dire « à mon avis », on utilise :
\begin{itemize}
\item A. \zh{你呢} \quad B. \zh{比如} \quad C. \zh{我觉得} \quad D. \zh{我同意}
\end{itemize}

3. \zh{喀麦隆的音乐比中国的音乐热闹。} signifie :
\begin{itemize}
\item A. La musique chinoise est plus animée que la musique camerounaise.
\item B. La musique camerounaise est plus animée que la musique chinoise.
\item C. Les deux musiques sont identiques.
\item D. La musique camerounaise est moins animée.
\end{itemize}

4. Quelle phrase permet de relancer la parole dans un débat ?
\begin{itemize}
\item A. \zh{我不同意。} \quad B. \zh{你呢？} \quad C. \zh{春节是中国的新年。} \quad D. \zh{我喜欢跳舞。}
\end{itemize}
\end{exercice}

\begin{exercice}[Vrai ou faux ? Justifie]
Dis si chaque affirmation est vraie (V) ou fausse (F), puis justifie ta réponse en français.

1. \zh{我觉得} (\emph{Wǒ juéde}) sert à donner son opinion personnelle.
2. \zh{比如} (\emph{bǐrú}) sert à comparer deux choses avec l'adjectif.
3. Dans un débat, \zh{我同意} (\emph{Wǒ tóngyì}) exprime le désaccord.
4. La structure de comparaison est : A + \zh{比} + B + adjectif.
5. \zh{虽然……但是……} (\emph{suīrán... dànshì...}) exprime une opposition entre deux idées.
\end{exercice}

\begin{exercice}[Vocabulaire : complète le tableau]
Complète le tableau suivant avec le caractère chinois, le pinyin ou la traduction manquante.

\begin{center}
\begin{tabularx}{\linewidth}{|X|X|X|X|}
\hline
\rowcolor{popblueL} Caractères & Pinyin & Nature & Traduction \\
\hline
文化 & \ldots\ldots\ldots & nom & la culture \\
\hline
\ldots\ldots\ldots & yīnyuè & nom & la musique \\
\hline
跳舞 & \ldots\ldots\ldots & verbe & danser \\
\hline
\ldots\ldots\ldots & biǎoyǎn & nom / verbe & le spectacle ; se produire \\
\hline
传统 & chuántǒng & nom & \ldots\ldots\ldots \\
\hline
\ldots\ldots\ldots & rènao & adjectif & animé, bruyant de joie \\
\hline
\end{tabularx}
\end{center}
\end{exercice}

\begin{exercice}[Pinyin : associe]
Relie chaque caractère ou mot à son pinyin. Attention aux tons !

\begin{center}
\begin{tabularx}{\linewidth}{|X|X|}
\hline
\rowcolor{popblueL} Caractères & Pinyin \\
\hline
1. 买 & a. mài \\
\hline
2. 卖 & b. cháng \\
\hline
3. 唱 & c. mǎi \\
\hline
4. 长 & d. chàng \\
\hline
5. 音乐 & e. wǔdǎo \\
\hline
6. 舞蹈 & f. yīnyuè \\
\hline
\end{tabularx}
\end{center}
\end{exercice}

\courssub{Je m'entraîne}

\begin{exercice}[Traduis en chinois]
Traduis les phrases suivantes en chinois (caractères + pinyin).

1. Je pense que la fête du Ngondo est très intéressante.
2. La musique camerounaise est plus animée que la musique chinoise.
3. Bien que j'aime les films chinois, je préfère la danse traditionnelle camerounaise.
4. Et toi, quelle fête est-ce que tu préfères ?
5. Par exemple, à Douala, on danse et on chante pendant les fêtes.
\end{exercice}

\begin{exercice}[Transforme les phrases]
Transforme chaque phrase selon la consigne entre parenthèses.

1. \zh{中国的春节很有意思。} (donne ton avis avec \zh{我觉得})
2. \zh{喀麦隆的音乐热闹。中国的音乐不热闹。} (fais une comparaison avec \zh{比})
3. \zh{我喜欢跳舞。我喜欢唱歌。} (relie les deux idées avec \zh{虽然……但是……} en exprimant une préférence)
4. \zh{你说得对。} (exprime l'accord avec \zh{我同意})
5. \zh{我喜欢中国的电影。} (pose la question à ton camarade avec \zh{你呢？})
\end{exercice}

\begin{exercice}[Dialogue à trous : les formules fonctionnelles]
Complète le dialogue avec les formules suivantes : \zh{我觉得} / \zh{我同意} / \zh{你呢} / \zh{比如} / \zh{虽然} / \zh{但是}. Puis joue le dialogue à deux.

\begin{textebox}[À la cantine du lycée]
A：你喜欢喀麦隆的文化吗？\\
\emph{Nǐ xǐhuan Kāmàilóng de wénhuà ma ?}\\
Aimes-tu la culture camerounaise ?\\[4pt]
B：当然喜欢！（1）\ldots\ldots\ldots 喀麦隆的音乐特别好听。\\
\emph{Dāngrán xǐhuan ! ... Kāmàilóng de yīnyuè tèbié hǎotīng.}\\
Bien sûr ! (1) \ldots la musique camerounaise est vraiment belle.\\[4pt]
A：（2）\ldots\ldots\ldots ！（3）\ldots\ldots\ldots ，mákòsà 和 bikutsi 都很有名。\\
\emph{... ! ..., makossa hé bikutsi dōu hěn yǒumíng.}\\
(2) \ldots ! (3) \ldots, le makossa et le bikutsi sont tous les deux célèbres.\\[4pt]
B：（4）\ldots\ldots\ldots 中国的音乐很优美，（5）\ldots\ldots\ldots 我更喜欢喀麦隆的音乐。（6）\ldots\ldots\ldots ？\\
\emph{... Zhōngguó de yīnyuè hěn yōuměi, ... wǒ gèng xǐhuan Kāmàilóng de yīnyuè. ... ?}\\
(4) \ldots la musique chinoise est très gracieuse, (5) \ldots je préfère la musique camerounaise. (6) \ldots ?
\end{textebox}
\end{exercice}

\begin{exercice}[Remets le débat dans l'ordre]
Remets dans l'ordre logique les répliques de ce débat sur la musique, puis joue-le avec ton propre avis.

1. \zh{我不同意！我觉得中国的京剧也很有趣。}\\
\emph{Wǒ bù tóngyì ! Wǒ juéde Zhōngguó de Jīngjù yě hěn yǒuqù.}\\
2. \zh{你呢？你更喜欢哪一种音乐？}\\
\emph{Nǐ ne ? Nǐ gèng xǐhuan nǎ yì zhǒng yīnyuè ?}\\
3. \zh{我觉得喀麦隆的音乐比中国的音乐热闹。}\\
\emph{Wǒ juéde Kāmàilóng de yīnyuè bǐ Zhōngguó de yīnyuè rènao.}\\
4. \zh{我更喜欢 bikutsi，因为它让我想跳舞！}\\
\emph{Wǒ gèng xǐhuan bikutsi, yīnwèi tā ràng wǒ xiǎng tiàowǔ !}\\
\end{exercice}

\courssub{Je me prépare au Bac}

\begin{exercice}[Compréhension de texte et questions de langue]
Lis le texte suivant, puis réponds aux questions.

\begin{textebox}[La fête de mon village]
在我的家乡，每年十二月有一个大节日。虽然这个节日没有春节那么有名，但是我们全村的人都非常喜欢。节日那天，大家穿上漂亮的传统衣服，又唱歌又跳舞。我觉得我们村的舞蹈比城里的舞蹈更有意思。比如，老人们讲以前的故事，年轻人表演传统的舞蹈。晚上，大家一起吃饭，非常热闹。中国的朋友，你们那里有什么节日？\\
\emph{Zài wǒ de jiāxiāng, měi nián shí'èr yuè yǒu yí ge dà jiérì. Suīrán zhège jiérì méiyǒu Chūnjié nàme yǒumíng, dànshì wǒmen quán cūn de rén dōu fēicháng xǐhuan. Jiérì nà tiān, dàjiā chuān shang piàoliang de chuántǒng yīfu, yòu chànggē yòu tiàowǔ. Wǒ juéde wǒmen cūn de wǔdǎo bǐ chénglǐ de wǔdǎo gèng yǒu yìsi. Bǐrú, lǎorénmen jiǎng yǐqián de gùshi, niánqīngrénmen biǎoyǎn chuántǒng de wǔdǎo. Wǎnshang, dàjiā yìqǐ chīfàn, fēicháng rènao. Zhōngguó de péngyou, nǐmen nàli yǒu shénme jiérì ?}\\
Dans mon village natal, il y a chaque année en décembre une grande fête. Bien que cette fête ne soit pas aussi célèbre que la Fête du Printemps, tous les gens de notre village l'aiment beaucoup. Le jour de la fête, tout le monde met de beaux vêtements traditionnels, chante et danse. Je trouve que les danses de notre village sont plus intéressantes que celles de la ville. Par exemple, les anciens racontent des histoires d'autrefois, les jeunes présentent des danses traditionnelles. Le soir, tout le monde mange ensemble, c'est très animé. Amis chinois, quelles fêtes y a-t-il chez vous ?
\end{textebox}

\textbf{Questions de compréhension} (réponds en français) :
\begin{enumerate}
\item Quand a lieu la fête décrite dans le texte ?
\item Que font les habitants le jour de la fête ? Cite deux activités.
\item Quelle comparaison l'auteur fait-il entre son village et la ville ?
\end{enumerate}

\textbf{Questions de langue} :
\begin{enumerate}
\item Relève dans le texte une phrase avec la structure \zh{虽然……但是……} et traduis-la.
\item Relève la comparaison avec \zh{比} et explique sa structure.
\item Quelle formule l'auteur utilise-t-il pour donner son avis ? Recopie-la avec son pinyin.
\end{enumerate}
\end{exercice}

\begin{exercice}[Thème et version]
\textbf{A. Thème.} Traduis en chinois (caractères + pinyin) :
\begin{enumerate}
\item Je pense que la culture camerounaise est très riche.
\item Bien que le Nouvel An chinois soit célèbre, je préfère les fêtes de mon village.
\item Et toi, est-ce que tu aimes la danse traditionnelle ?
\end{enumerate}

\textbf{B. Version.} Traduis en français :
\begin{enumerate}
\item \zh{我觉得中国的春节比喀麦隆的圣诞节更热闹。}\\
\emph{Wǒ juéde Zhōngguó de Chūnjié bǐ Kāmàilóng de Shèngdànjié gèng rènao.}
\item \zh{虽然京剧很难，但是很多年轻人喜欢。}\\
\emph{Suīrán Jīngjù hěn nán, dànshì hěn duō niánqīngrén xǐhuan.}
\item \zh{比如，在雅温得，学生们常常表演节目。}\\
\emph{Bǐrú, zài Yǎwēndé, xuéshengmen chángcháng biǎoyǎn jiémù.}
\end{enumerate}
\end{exercice}

\begin{exercice}[Expression écrite puis orale : exposé type Bac]
\textbf{Sujet :} Un correspondant chinois visite ton lycée. Présente-lui une fête traditionnelle camerounaise.

\textbf{Consignes :}
\begin{enumerate}
\item Rédige d'abord un mini-exposé écrit d'au moins 8 phrases (environ 60 à 80 caractères chinois) intitulé \zh{我最喜欢的节日} (\emph{Wǒ zuì xǐhuan de jiérì} — Ma fête préférée).
\item Utilise obligatoirement : au moins deux formules d'opinion (\zh{我觉得}, \zh{我想}...), une comparaison avec \zh{比}, une phrase avec \zh{虽然……但是……}, et termine par une relance avec \zh{你呢？}.
\item Prépare ensuite la présentation orale de ton exposé : travaille la prononciation des tons, parle lentement et regarde ton auditoire.
\item Ton camarade chinois (joué par un élève) doit te poser au moins une question à la fin ; réponds avec une formule fonctionnelle (\zh{我同意}, \zh{比如}...).
\end{enumerate}
\end{exercice}

\newpage

\coursec{Corrigés des exercices}

\begin{corrige}[Exercice 1 — QCM : la vie culturelle]
\textbf{1. B.} \zh{节日} (\emph{jiérì}) signifie « la fête ». Le mot est composé de \zh{节} (fête) et \zh{日} (jour) : littéralement « jour de fête ». Attention à ne pas confondre avec \zh{节目} (\emph{jiémù}) « le programme, le numéro de spectacle ».

\textbf{2. C.} \zh{我觉得} (\emph{Wǒ juéde}) = « je pense, je trouve » : c'est la formule d'opinion personnelle. \zh{你呢} sert à relancer, \zh{比如} à donner un exemple, \zh{我同意} à exprimer l'accord.

\textbf{3. B.} \zh{喀麦隆的音乐比中国的音乐热闹。} (\emph{Kāmàilóng de yīnyuè bǐ Zhōngguó de yīnyuè rènao.}) = « La musique camerounaise est plus animée que la musique chinoise. » Dans la structure A + \zh{比} + B + adjectif, c'est A qui possède davantage la qualité exprimée par l'adjectif.

\textbf{4. B.} \zh{你呢？} (\emph{Nǐ ne ?}) = « Et toi ? » : c'est la formule de relance qui rend la parole à l'interlocuteur. Les autres répliques ne relancent pas le débat.
\end{corrige}

\begin{corrige}[Exercice 2 — Vrai ou faux ? Justifie]
\textbf{1. Vrai.} \zh{我觉得} (\emph{Wǒ juéde}) introduit un jugement personnel : « je pense que, je trouve que ». Exemple : \zh{我觉得春节很有意思。} (\emph{Wǒ juéde Chūnjié hěn yǒu yìsi.}) « Je trouve la Fête du Printemps très intéressante. »

\textbf{2. Faux.} \zh{比如} (\emph{bǐrú}) signifie « par exemple » : il sert à illustrer une idée par un exemple, pas à comparer. La comparaison se fait avec \zh{比} (\emph{bǐ}) : A + \zh{比} + B + adjectif. Attention au faux ami visuel : \zh{比如} contient \zh{比} mais n'exprime pas la comparaison.

\textbf{3. Faux.} \zh{我同意} (\emph{Wǒ tóngyì}) exprime l'accord (« je suis d'accord »). Le désaccord s'exprime avec la négation : \zh{我不同意} (\emph{Wǒ bù tóngyì}) « je ne suis pas d'accord ».

\textbf{4. Vrai.} La structure de comparaison est bien : A + \zh{比} + B + adjectif. Exemple : \zh{中国的春节比喀麦隆的圣诞节热闹。} (\emph{Zhōngguó de Chūnjié bǐ Kāmàilóng de Shèngdànjié rènao.}) « La Fête du Printemps chinoise est plus animée que Noël au Cameroun. »

\textbf{5. Vrai.} \zh{虽然……但是……} (\emph{suīrán... dànshì...}) « bien que... mais... » exprime une opposition (concession) entre deux idées. En chinois, contrairement au français, on peut — et on préfère — utiliser les deux connecteurs dans la même phrase.
\end{corrige}

\begin{corrige}[Exercice 3 — Vocabulaire : complète le tableau]
\begin{center}
\begin{tabularx}{\linewidth}{|X|X|X|X|}
\hline
\rowcolor{popblueL} Caractères & Pinyin & Nature & Traduction \\
\hline
文化 & wénhuà & nom & la culture \\
\hline
音乐 & yīnyuè & nom & la musique \\
\hline
跳舞 & tiàowǔ & verbe & danser \\
\hline
表演 & biǎoyǎn & nom / verbe & le spectacle ; se produire \\
\hline
传统 & chuántǒng & nom / adjectif & la tradition ; traditionnel \\
\hline
热闹 & rènao & adjectif & animé, bruyant de joie \\
\hline
\end{tabularx}
\end{center}

\textbf{Points de vigilance :}
\begin{itemize}
\item \zh{音乐} : le premier caractère \zh{音} (son) porte le 1\ier{} ton ; attention à ne pas écrire \zh{因} (cause), homophone.
\item \zh{跳舞} : verbe séparable — on dit \zh{跳一个舞} (\emph{tiào yí ge wǔ}) « danser une danse ».
\item \zh{热闹} : \zh{闹} porte le ton léger dans ce mot (\emph{rènao}), pas le 4\ieme{} ton plein. À l'oral, un ton plein sur \zh{闹} est une erreur fréquente des francophones.
\end{itemize}
\end{corrige}

\begin{corrige}[Exercice 4 — Pinyin : associe]
\begin{center}
\begin{tabularx}{\linewidth}{|X|X|}
\hline
\rowcolor{popblueL} Caractères & Pinyin \\
\hline
1. 买 & c. mǎi (3\ieme{} ton — acheter) \\
\hline
2. 卖 & a. mài (4\ieme{} ton — vendre) \\
\hline
3. 唱 & d. chàng (4\ieme{} ton — chanter) \\
\hline
4. 长 & b. cháng (2\ieme{} ton — long) \\
\hline
5. 音乐 & f. yīnyuè \\
\hline
6. 舞蹈 & e. wǔdǎo \\
\hline
\end{tabularx}
\end{center}

\textbf{Points de vigilance :}
\begin{itemize}
\item \zh{买} \emph{mǎi} / \zh{卖} \emph{mài} : même syllabe, seul le ton change le sens (acheter / vendre). C'est un piège classique pour les francophones. Moyen mnémotechnique : \zh{卖} a un « chapeau » \zh{十} en plus — celui qui vend a quelque chose en plus à offrir.
\item \zh{长} se lit \emph{cháng} « long » mais aussi \emph{zhǎng} « grandir » : c'est un caractère à plusieurs lectures (polyphonique).
\item \zh{唱} \emph{chàng} (4\ieme{} ton descendant) ne doit pas être confondu avec \zh{长} \emph{cháng} (2\ieme{} ton montant).
\end{itemize}
\end{corrige}

\begin{corrige}[Exercice 5 — Traduis en chinois]
\textbf{1.} \zh{我觉得恩贡多节很有意思。}\\
\emph{Wǒ juéde Ēngòngduō jié hěn yǒu yìsi.}\\
Structure : \zh{我觉得} + proposition ; l'adjectif \zh{有意思} est précédé de \zh{很} (nécessaire pour lier un adjectif au sujet dans une phrase affirmative simple, sans valeur intensive réelle).

\textbf{2.} \zh{喀麦隆的音乐比中国的音乐热闹。}\\
\emph{Kāmàilóng de yīnyuè bǐ Zhōngguó de yīnyuè rènao.}\\
Structure : A + \zh{比} + B + adjectif. Attention : pas de \zh{很} devant l'adjectif dans une comparaison ; pour renforcer, on utilise \zh{更} ou \zh{还}.

\textbf{3.} \zh{虽然我喜欢中国的电影，但是我更喜欢喀麦隆的传统舞蹈。}\\
\emph{Suīrán wǒ xǐhuan Zhōngguó de diànyǐng, dànshì wǒ gèng xǐhuan Kāmàilóng de chuántǒng wǔdǎo.}\\
\zh{更} (\emph{gèng}) « davantage » renforce la préférence.

\textbf{4.} \zh{你呢？你最喜欢什么节日？}\\
\emph{Nǐ ne ? Nǐ zuì xǐhuan shénme jiérì ?}\\
La relance \zh{你呢？} se place en début de tour de parole ; \zh{最} (\emph{zuì}) marque le superlatif « le plus » ; \zh{什么} est le mot interrogatif « quel, quoi ».

\textbf{5.} \zh{比如，在杜阿拉，过节的时候大家又跳舞又唱歌。}\\
\emph{Bǐrú, zài Dù'ālā, guò jié de shíhou dàjiā yòu tiàowǔ yòu chànggē.}\\
La structure \zh{又……又……} (\emph{yòu... yòu...}) exprime deux actions simultanées : « à la fois... et... ». \zh{过节} (\emph{guò jié}) = « fêter, célébrer la fête » ; \zh{的时候} (\emph{de shíhou}) = « au moment où, quand ».
\end{corrige}

\begin{corrige}[Exercice 6 — Transforme les phrases]
\textbf{1.} \zh{我觉得中国的春节很有意思。}\\
\emph{Wǒ juéde Zhōngguó de Chūnjié hěn yǒu yìsi.}\\
On place simplement \zh{我觉得} en tête de phrase ; le reste ne change pas.

\textbf{2.} \zh{喀麦隆的音乐比中国的音乐热闹。}\\
\emph{Kāmàilóng de yīnyuè bǐ Zhōngguó de yīnyuè rènao.}\\
Les deux informations « animée / pas animée » se fondent en une seule comparaison : A + \zh{比} + B + adjectif. Ne pas ajouter \zh{很} devant \zh{热闹}.

\textbf{3.} \zh{虽然我喜欢跳舞，但是我更喜欢唱歌。}\\
\emph{Suīrán wǒ xǐhuan tiàowǔ, dànshì wǒ gèng xǐhuan chànggē.}\\
Pour exprimer une préférence, on ajoute \zh{更} dans la seconde proposition. (Solution inverse acceptée : \zh{虽然我喜欢唱歌，但是我更喜欢跳舞。})

\textbf{4.} \zh{我同意，你说得对。}\\
\emph{Wǒ tóngyì, nǐ shuō de duì.}\\
\zh{我同意} renforce la formule \zh{你说得对} « tu as raison ». Attention : dans \zh{说得对}, le mot-outil \zh{得} (\emph{de}, ton léger) introduit le complément de manière \zh{对} — ne pas écrire \zh{的}.

\textbf{5.} \zh{我喜欢中国的电影。你呢？}\\
\emph{Wǒ xǐhuan Zhōngguó de diànyǐng. Nǐ ne ?}\\
La relance \zh{你呢？} suit directement l'énoncé de son propre avis.
\end{corrige}

\begin{corrige}[Exercice 7 — Dialogue à trous : les formules fonctionnelles]
\textbf{Solutions :} (1) \zh{我觉得} \quad (2) \zh{我同意} \quad (3) \zh{比如} \quad (4) \zh{虽然} \quad (5) \zh{但是} \quad (6) \zh{你呢}

\textbf{Dialogue complet :}

\begin{textebox}[À la cantine du lycée — version corrigée]
A：你喜欢喀麦隆的文化吗？\\
\emph{Nǐ xǐhuan Kāmàilóng de wénhuà ma ?}\\
Aimes-tu la culture camerounaise ?\\[4pt]
B：当然喜欢！我觉得喀麦隆的音乐特别好听。\\
\emph{Dāngrán xǐhuan ! Wǒ juéde Kāmàilóng de yīnyuè tèbié hǎotīng.}\\
Bien sûr ! Je trouve la musique camerounaise vraiment belle.\\[4pt]
A：我同意！比如，马科萨和比库西都很有名。\\
\emph{Wǒ tóngyì ! Bǐrú, Mǎkēsà hé Bǐkùxī dōu hěn yǒumíng.}\\
Je suis d'accord ! Par exemple, le makossa et le bikutsi sont tous les deux célèbres.\\[4pt]
B：虽然中国的音乐很优美，但是我更喜欢喀麦隆的音乐。你呢？\\
\emph{Suīrán Zhōngguó de yīnyuè hěn yōuměi, dànshì wǒ gèng xǐhuan Kāmàilóng de yīnyuè. Nǐ ne ?}\\
Bien que la musique chinoise soit très gracieuse, je préfère la musique camerounaise. Et toi ?
\end{textebox}

\textbf{Justifications :} (1) opinion personnelle après une affirmation ; (2) accord avec le point de vue précédent ; (3) introduction d'exemples (makossa, bikutsi) ; (4)-(5) paire concessive inséparable \zh{虽然……但是……} ; (6) relance finale pour rendre la parole.

\textbf{Remarque :} les noms des danses camerounaises se transcrivent en caractères : \zh{马科萨} (\emph{Mǎkēsà}) « makossa » et \zh{比库西} (\emph{Bǐkùxī}) « bikutsi ». Dans un texte chinois, on évite de mélanger l'alphabet latin et les caractères.
\end{corrige}

\begin{corrige}[Exercice 8 — Remets le débat dans l'ordre]
\textbf{Ordre logique : 3 → 1 → 2 → 4.}

\begin{enumerate}
\item[3.] \zh{我觉得喀麦隆的音乐比中国的音乐热闹。} — Ouverture du débat : un avis avec comparaison.
\item[1.] \zh{我不同意！我觉得中国的京剧也很有趣。} — Réaction : désaccord + contre-argument.
\item[2.] \zh{你呢？你更喜欢哪一种音乐？} — Relance : on rend la parole.
\item[4.] \zh{我更喜欢比库西，因为它让我想跳舞！} — Conclusion : préférence justifiée avec \zh{因为} (\emph{yīnwèi}) « parce que ».
\end{enumerate}

\textbf{Logique du débat :} opinion → désaccord → relance → préférence argumentée. C'est exactement la progression attendue dans un échange oral en classe de Première : présenter, réagir, relancer, conclure.

\textbf{Point de langue :} dans la phrase 4, \zh{它让我想跳舞} (\emph{tā ràng wǒ xiǎng tiàowǔ}) « elle me donne envie de danser » illustre la structure pivot : \zh{让} + personne + verbe.
\end{corrige}

\begin{corrige}[Exercice 9 — Compréhension de texte et questions de langue]
\textbf{Questions de compréhension :}

\textbf{1.} La fête a lieu chaque année au mois de décembre (\zh{每年十二月}, \emph{měi nián shí'èr yuè}).

\textbf{2.} Le jour de la fête, les habitants mettent de beaux vêtements traditionnels, chantent et dansent (\zh{大家穿上漂亮的传统衣服，又唱歌又跳舞}). On peut aussi citer : les anciens racontent des histoires, les jeunes présentent des danses traditionnelles, et le soir tout le monde mange ensemble.

\textbf{3.} L'auteur compare les danses de son village à celles de la ville : il trouve que les danses de son village sont plus intéressantes que celles de la ville (\zh{我们村的舞蹈比城里的舞蹈更有意思}).

\textbf{Questions de langue :}

\textbf{1.} \zh{虽然这个节日没有春节那么有名，但是我们全村的人都非常喜欢。}\\
\emph{Suīrán zhège jiérì méiyǒu Chūnjié nàme yǒumíng, dànshì wǒmen quán cūn de rén dōu fēicháng xǐhuan.}\\
« Bien que cette fête ne soit pas aussi célèbre que la Fête du Printemps, tous les gens de notre village l'aiment beaucoup. » Structure : \zh{虽然} + concession, \zh{但是} + idée principale. On note aussi la comparaison d'égalité niée : \zh{没有} + B + \zh{那么} + adjectif « ne pas être aussi... que ».

\textbf{2.} \zh{我觉得我们村的舞蹈比城里的舞蹈更有意思。} Structure : A (\zh{我们村的舞蹈}) + \zh{比} + B (\zh{城里的舞蹈}) + adjectif (\zh{有意思}), renforcé par \zh{更} « encore plus ». C'est A qui possède davantage la qualité.

\textbf{3.} L'auteur utilise \zh{我觉得} (\emph{Wǒ juéde}) « je pense, je trouve » pour donner son avis.

\textbf{Barème indicatif (sur 20) :} compréhension 8 pts (2 + 4 + 2) ; questions de langue 12 pts (4 + 4 + 4). Exiger la phrase relevée en caractères exacts et une traduction fidèle.
\end{corrige}

\begin{corrige}[Exercice 10 — Thème et version]
\textbf{A. Thème.}

\textbf{1.} \zh{我觉得喀麦隆的文化很丰富。}\\
\emph{Wǒ juéde Kāmàilóng de wénhuà hěn fēngfù.}\\
\zh{丰富} (\emph{fēngfù}) « riche, abondant » qualifie la culture, la nourriture, l'expérience.

\textbf{2.} \zh{虽然中国的新年很有名，但是我更喜欢我们村的节日。}\\
\emph{Suīrán Zhōngguó de xīnnián hěn yǒumíng, dànshì wǒ gèng xǐhuan wǒmen cūn de jiérì.}\\
On peut aussi dire \zh{春节} pour « le Nouvel An chinois ».

\textbf{3.} \zh{你呢？你喜欢传统舞蹈吗？}\\
\emph{Nǐ ne ? Nǐ xǐhuan chuántǒng wǔdǎo ma ?}\\
Question fermée avec la particule finale \zh{吗}.

\textbf{B. Version.}

\textbf{1.} « Je trouve que la Fête du Printemps chinoise est plus animée que Noël au Cameroun. » (\zh{更} = « encore plus », à rendre par le comparatif.)

\textbf{2.} « Bien que l'opéra de Pékin soit difficile, beaucoup de jeunes l'aiment. » (\zh{京剧} = l'opéra de Pékin, théâtre traditionnel chinois.)

\textbf{3.} « Par exemple, à Yaoundé, les élèves présentent souvent des spectacles. » (\zh{雅温得} \emph{Yǎwēndé} = Yaoundé ; \zh{节目} = numéro, programme, spectacle.)

\textbf{Points de vigilance :} ne pas traduire \zh{但是} par un « mais » lourd en début de phrase française quand « bien que » suffit ; en thème, ne jamais mettre \zh{很} devant l'adjectif d'une comparaison avec \zh{比}.

\textbf{Barème indicatif (sur 20) :} thème 10 pts (3 + 4 + 3) ; version 10 pts (3 + 4 + 3). Une faute de ton dans le pinyin = 0.25pt ; une faute de structure (ordre des mots, \zh{比} mal placé) = 1 pt.
\end{corrige}

\begin{corrige}[Exercice 11 — Expression écrite puis orale : exposé type Bac]
\textbf{Proposition de rédaction modèle :}

\begin{textebox}[我最喜欢的节日]
我最喜欢的节日是杜阿拉的恩贡多节。我觉得这个节日非常有意思。每年十二月，大家都穿上漂亮的传统衣服，又唱歌又跳舞。虽然恩贡多节没有春节那么有名，但是我认为它更热闹。比如，水上的划船比赛特别精彩，比足球比赛还让人激动！我想，传统节日让我们记住自己的文化。你呢？你最喜欢什么节日？\\
\emph{Wǒ zuì xǐhuan de jiérì shì Dù'ālā de Ēngòngduō jié. Wǒ juéde zhège jiérì fēicháng yǒu yìsi. Měi nián shí'èr yuè, dàjiā dōu chuān shang piàoliang de chuántǒng yīfu, yòu chànggē yòu tiàowǔ. Suīrán Ēngòngduō jié méiyǒu Chūnjié nàme yǒumíng, dànshì wǒ rènwéi tā gèng rènao. Bǐrú, shuǐ shang de huáchuán bǐsài tèbié jīngcǎi, bǐ zúqiú bǐsài hái ràng rén jīdòng ! Wǒ xiǎng, chuántǒng jiérì ràng wǒmen jìzhù zìjǐ de wénhuà. Nǐ ne ? Nǐ zuì xǐhuan shénme jiérì ?}\\
Ma fête préférée est la fête du Ngondo à Douala. Je trouve cette fête très intéressante. Chaque année en décembre, tout le monde met de beaux vêtements traditionnels, chante et danse. Bien que le Ngondo ne soit pas aussi célèbre que la Fête du Printemps, je pense qu'il est plus animé. Par exemple, les courses de pirogues sur l'eau sont vraiment spectaculaires, encore plus passionnantes qu'un match de football ! Je pense que les fêtes traditionnelles nous font nous souvenir de notre culture. Et toi ? Quelle est ta fête préférée ?
\end{textebox}

\textbf{Vérification des contraintes :} deux formules d'opinion (\zh{我觉得}, \zh{我认为}, \zh{我想}) ; une comparaison avec \zh{比} (\zh{比足球比赛还让人激动}) ; une concession \zh{虽然……但是……} ; une relance finale \zh{你呢？}.

\textbf{Exemple de question du camarade et de réponse :}\\
— \zh{恩贡多节在哪个月？} (\emph{Ēngòngduō jié zài nǎ ge yuè ?}) « En quel mois a lieu le Ngondo ? »\\
— \zh{在十二月。我同意，这个时候天气很好，比如，大家可以看划船比赛。} (\emph{Zài shí'èr yuè. Wǒ tóngyì, zhège shíhou tiānqì hěn hǎo, bǐrú, dàjiā kěyǐ kàn huáchuán bǐsài.})

\textbf{Points de vigilance :}
\begin{itemize}
\item Pas de \zh{很} devant l'adjectif dans la comparaison avec \zh{比}.
\item \zh{虽然} et \zh{但是} s'emploient ensemble en chinois, contrairement au français.
\item À l'oral : soigner le 4\ieme{} ton de \zh{杜阿拉} (\emph{Dù'ālā}), le ton léger de \zh{热闹} (\emph{rènao}) et le 3\ieme{} ton de \zh{我} (\emph{wǒ}).
\item Parler lentement, regarder l'auditoire, annoncer le plan : présentation de la fête → activités → avis personnel → relance.
\end{itemize}

\textbf{Barème indicatif (sur 20) :} respect des contraintes et contenu 6 pts ; correction grammaticale (structures \zh{比}, \zh{虽然……但是……}, ordre des mots) 6 pts ; richesse du vocabulaire et caractères 4 pts ; prestation orale (tons, fluidité, regard, réponse à la question) 4 pts.
\end{corrige}

\newpage

\coursec{Fiche bilan}

\courssub{Carte mentale de la leçon}
\begin{popfigure}
\begin{tikzpicture}[
  mindmap,
  grow cyclic,
  every node/.style={concept, execute at begin node=\hskip0pt},
  concept color=popblue,
  level 1/.append style={level distance=4.5cm, sibling angle=60},
  level 2/.append style={level distance=3cm, sibling angle=45},
  branch/.style={concept color=#1, edge from parent/.style={draw=#1, thick}}
]
\node[concept, text width=3.2cm, align=center, font=\bfseries\large] {Vie culturelle\\Cameroun \& Chine (ZHA21)}
  child[branch=poporange] {node[concept, text width=2.8cm, align=center] {Lexique clé\\Fêtes, arts, coutumes}}
  child[branch=popgreen] {node[concept, text width=2.8cm, align=center] {Dialogue 1\\Présenter une fête}}
  child[branch=poppurple] {node[concept, text width=2.8cm, align=center] {Formules\\Présenter / Opiner / Relancer}}
  child[branch=poppink] {node[concept, text width=2.8cm, align=center] {Tons\\Entraînement ciblé}}
  child[branch=popgold] {node[concept, text width=2.8cm, align=center] {Dialogue 2\\Débat \& Avis}}
  child[branch=popdark] {node[concept, text width=2.8cm, align=center] {Jeux de rôle\\Débat guidé}};
\end{tikzpicture}
\legende{Vue d'ensemble de la leçon ZHA21 : Expression orale — Vie culturelle}
\end{popfigure}

\courssub{Bilan synthétique}
\begin{bilan}

\textbf{1. Lexique clé (Vie culturelle — Fêtes, Arts, Coutumes)}

\begin{center}
\begin{tabularx}{\linewidth}{|X|X|X|X|}
\hline
\rowcolor{popblueL} \textbf{Caractères} & \textbf{Pinyin} & \textbf{Nature} & \textbf{Traduction} \\ \hline
春节 & Chūnjié & n. & Fête du Printemps (Nouvel An chinois) \\ \hline
中秋节 & Zhōngqiūjié & n. & Fête de la Mi-Automne \\ \hline
端午节 & Duānwǔjié & n. & Fête des Bateaux-Dragons \\ \hline
国庆节 & Guóqìngjié & n. & Fête Nationale (1er oct.) \\ \hline
圣诞节 & Shèngdànjié & n. & Noël \\ \hline
新年 & Xīnnián & n. & Nouvel An \\ \hline
传统节日 & chuántǒng jiérì & n. & Fête traditionnelle \\ \hline
习俗 & xísú & n. & Coutume, tradition \\ \hline
饺子 & jiǎozi & n. & Raviolis (jiaozi) \\ \hline
月饼 & yuèbǐng & n. & Gâteau de lune \\ \hline
粽子 & zòngzi & n. & Zongzi (riz gluant feuille bambou) \\ \hline
舞狮 & wǔshī & v./n. & Danse du lion \\ \hline
舞龙 & wǔlóng & v./n. & Danse du dragon \\ \hline
放鞭炮 & fàng biānpào & v. & Lancer des pétards \\ \hline
赏月 & shǎngyuè & v. & Admirer la lune \\ \hline
团圆 & tuányuán & adj./v. & Réunion (familiale), se réunir \\ \hline
喀麦隆 & Kāmàilóng & n. & Cameroun \\ \hline
恩多莱 & ēnduólái & n. & Ndolé (plat national) \\ \hline
马卡西 & mǎkǎxī & n. & Makossa (danse/musique) \\ \hline
足球 & zúqiú & n. & Football \\ \hline
文化 & wénhuà & n. & Culture \\ \hline
交流 & jiāoliú & v./n. & Échanger, échange \\ \hline
介绍 & jièshào & v. & Présenter, introduire \\ \hline
认为 & rènwéi & v. & Penser que, estimer que \\ \hline
赞成 & zànchéng & v. & Être d'accord, approuver \\ \hline
反对 & fǎnduì & v. & Être contre, s'opposer \\ \hline
建议 & jiànyì & v./n. & Suggérer, suggestion \\ \hline
\end{tabularx}
\end{center}

\textbf{2. Règles de grammaire / Structures fonctionnelles à connaître par cœur}

\begin{center}
\begin{tabularx}{\linewidth}{|X|X|X|}
\hline
\rowcolor{popblueL} \textbf{Structure / Fonction} & \textbf{Exemple (Caractères + Pinyin)} & \textbf{Traduction} \\ \hline
\zh{先……再……} (Ordre chronologique) & \zh{我们先吃饺子，再看春晚。} (Wǒmen xiān chī jiǎozi, zài kàn Chūnwǎn.) & Nous mangeons d'abord des raviolis, puis regardons le Gala. \\ \hline
\zh{不仅……而且……} (Progression) & \zh{春节不仅吃饺子，而且放鞭炮。} (Chūnjié bùjǐn chī jiǎozi, érqiě fàng biānpào.) & Au Nouvel An, on ne mange pas seulement des raviolis, mais on lance aussi des pétards. \\ \hline
\zh{我觉得……} / \zh{我认为……} (Exprimer opinion) & \zh{我觉得中秋节最有意义。} (Wǒ juéde Zhōngqiūjié zuì yǒu yìyì.) & Je trouve que la Fête de la Mi-Automne est la plus significative. \\ \hline
\zh{你怎么看？} / \zh{你觉得呢？} (Relancer / Demander avis) & \zh{关于端午节，你怎么看？} (Guānyú Duānwǔjié, nǐ zěnme kàn?) & Concernant la Fête des Bateaux-Dragons, qu'en penses-tu ? \\ \hline
\zh{一方面……另一方面……} (Nuancer) & \zh{一方面热闹，另一方面太累。} (Yīfāngmiàn rènào, lìngyī fāngmiàn tài lèi.) & D'un côté c'est animé, de l'autre c'est trop fatigant. \\ \hline
\zh{建议 + V / 建议 + 对象 + V} (Suggérer) & \zh{我建议我们一起包粽子。} (Wǒ jiànyì wǒmen yìqǐ bāo zòngzi.) & Je suggère qu'on fasse des zongzi ensemble. \\ \hline
\zh{同意 / 不同意} (Accord/Désaccord) & \zh{我同意你的看法。} (Wǒ tóngyì nǐ de kànfǎ.) & Je suis d'accord avec ton point de vue. \\ \hline
\end{tabularx}
\end{center}

\textbf{3. Méthodes express}

\end{bilan}

\begin{methode}[Réussir un exposé oral court (1-2 min)]
\begin{enumerate}
\item \textbf{Préparer} : 3 mots-clés (Sujet, Détails culturels, Sentiment personnel).
\item \textbf{Structurer} : \zh{大家好，今天我介绍……} (Introduction) → \zh{首先……其次……最后……} (Corps) → \zh{总之，我觉得……谢谢。} (Conclusion).
\item \textbf{Prononcer} : Repérer les tons 3/3 (→ 2/3) et le \cle{neutre} sur les particules (\zh{了, 的, 呢}). S'enregistrer.
\item \textbf{Regarder} : Contact visuel, pas de lecture mot à mot.
\end{enumerate}
\end{methode}

\begin{methode}[Participer à un débat guidé]
\begin{enumerate}
\item \textbf{Écouter} : Noter le mot-clé de l'interlocuteur (\zh{文化, 传统, 现代}).
\item \textbf{Réagir} : \zh{我赞同……因为……} / \zh{我不太同意，因为……} / \zh{我想补充一点：……}.
\item \textbf{Relancer} : \zh{那……呢？} / \zh{你有什么建议？}.
\end{enumerate}
\end{methode}

\textbf{4. Points de vigilance (Francophones)}

\begin{attention}[Pièges fréquents]
\begin{itemize}
\item \textbf{Tons :} \zh{节} (jié, 2e ton) vs \zh{结} (jié/jiē). \zh{舞} (wǔ, 3e ton) → \zh{舞狮} (wǔshī, 3-1). Ne pas dire *wúshī.
\item \textbf{Ordre des mots :} Le complément de temps/lieu \textbf{avant} le verbe (\zh{在喀麦隆过圣诞节}, pas *\zh{过圣诞节在喀麦隆}).
\item \textbf{Classificateurs :} \zh{一个节日} (gè), \zh{一场舞狮} (chǎng), \zh{一家人} (jiā/kuò).
\item \zh{觉得} (sentiment subjectif) vs \zh{认为} (jugement réfléchi) vs \zh{想} (vouloir/planifier).
\item \zh{和} (nom + nom) ne relie pas des phrases ; utiliser \zh{而且} ou \zh{并且}.
\end{itemize}
\end{attention}



\courssub{Checklist avant le Bac}
\begin{aretenir}[Checklist ZHA21 — Expression orale : Vie culturelle]
$\square$ Je cite 5 fêtes chinoises + 2 fêtes camerounaises en chinois (caractères + pinyin + tons justes).\\
$\square$ Je décris une coutume (repas, danse, activité) avec \zh{先……再……} et \zh{不仅……而且……}.\\
$\square$ Je présente un exposé de 1 min : salutation, sujet, 2-3 détails, avis personnel, conclusion, au revoir.\\
$\square$ J'utilise \zh{我觉得 / 我认为} pour donner mon avis et \zh{你怎么看？} pour relancer.\\
$\square$ J'exprime accord (\zh{赞成/同意}) et désaccord (\zh{反对/不同意}) avec \zh{因为……}.\\
$\square$ Je prononce correctement les paires tonales : \zh{节日 (jiérì)}, \zh{传统 (chuántǒng)}, \zh{喀麦隆 (Kāmàilóng)}, \zh{团圆 (tuányuán)}.\\
$\square$ Je distingue \zh{在} (lieu/action en cours) et \zh{过} (expérience passée) dans mes phrases.\\
$\square$ J'écris correctement les caractères : \zh{节, 饺, 饼, 粽, 狮, 龙, 恩, 多, 莱}.\\
$\square$ Je compare une similitude et une différence culturelle Cameroun/Chine (ex: repas familial, danse).\\
$\square$ Je participe à un jeu de rôle : inviter un ami chinois à une fête camerounaise (ou inversement).\\
$\square$ Je gère mon trac : respiration, articulation, regard vers le public, débit modéré.\\
$\square$ Je connais les formules de politesse orale : \zh{请问}, \zh{不好意思}, \zh{没关系}, \zh{谢谢您的分享}.
\end{aretenir}

\findecours

\end{document}
